% Options for packages loaded elsewhere
\PassOptionsToPackage{unicode}{hyperref}
\PassOptionsToPackage{hyphens}{url}
\documentclass[
  english,
  11pt,
]{article}
\usepackage{xcolor}
\usepackage[margin=1in]{geometry}
\usepackage{amsmath,amssymb}
\setcounter{secnumdepth}{5}
\usepackage{iftex}
\ifPDFTeX
  \usepackage[T1]{fontenc}
  \usepackage[utf8]{inputenc}
  \usepackage{textcomp} % provide euro and other symbols
\else % if luatex or xetex
  \usepackage{unicode-math} % this also loads fontspec
  \defaultfontfeatures{Scale=MatchLowercase}
  \defaultfontfeatures[\rmfamily]{Ligatures=TeX,Scale=1}
\fi
\usepackage{lmodern}
\ifPDFTeX\else
  % xetex/luatex font selection
\fi
% Use upquote if available, for straight quotes in verbatim environments
\IfFileExists{upquote.sty}{\usepackage{upquote}}{}
\IfFileExists{microtype.sty}{% use microtype if available
  \usepackage[]{microtype}
  \UseMicrotypeSet[protrusion]{basicmath} % disable protrusion for tt fonts
}{}
\makeatletter
\@ifundefined{KOMAClassName}{% if non-KOMA class
  \IfFileExists{parskip.sty}{%
    \usepackage{parskip}
  }{% else
    \setlength{\parindent}{0pt}
    \setlength{\parskip}{6pt plus 2pt minus 1pt}}
}{% if KOMA class
  \KOMAoptions{parskip=half}}
\makeatother
\usepackage{color}
\usepackage{fancyvrb}
\newcommand{\VerbBar}{|}
\newcommand{\VERB}{\Verb[commandchars=\\\{\}]}
\DefineVerbatimEnvironment{Highlighting}{Verbatim}{commandchars=\\\{\}}
% Add ',fontsize=\small' for more characters per line
\usepackage{framed}
\definecolor{shadecolor}{RGB}{248,248,248}
\newenvironment{Shaded}{\begin{snugshade}}{\end{snugshade}}
\newcommand{\AlertTok}[1]{\textcolor[rgb]{0.94,0.16,0.16}{#1}}
\newcommand{\AnnotationTok}[1]{\textcolor[rgb]{0.56,0.35,0.01}{\textbf{\textit{#1}}}}
\newcommand{\AttributeTok}[1]{\textcolor[rgb]{0.13,0.29,0.53}{#1}}
\newcommand{\BaseNTok}[1]{\textcolor[rgb]{0.00,0.00,0.81}{#1}}
\newcommand{\BuiltInTok}[1]{#1}
\newcommand{\CharTok}[1]{\textcolor[rgb]{0.31,0.60,0.02}{#1}}
\newcommand{\CommentTok}[1]{\textcolor[rgb]{0.56,0.35,0.01}{\textit{#1}}}
\newcommand{\CommentVarTok}[1]{\textcolor[rgb]{0.56,0.35,0.01}{\textbf{\textit{#1}}}}
\newcommand{\ConstantTok}[1]{\textcolor[rgb]{0.56,0.35,0.01}{#1}}
\newcommand{\ControlFlowTok}[1]{\textcolor[rgb]{0.13,0.29,0.53}{\textbf{#1}}}
\newcommand{\DataTypeTok}[1]{\textcolor[rgb]{0.13,0.29,0.53}{#1}}
\newcommand{\DecValTok}[1]{\textcolor[rgb]{0.00,0.00,0.81}{#1}}
\newcommand{\DocumentationTok}[1]{\textcolor[rgb]{0.56,0.35,0.01}{\textbf{\textit{#1}}}}
\newcommand{\ErrorTok}[1]{\textcolor[rgb]{0.64,0.00,0.00}{\textbf{#1}}}
\newcommand{\ExtensionTok}[1]{#1}
\newcommand{\FloatTok}[1]{\textcolor[rgb]{0.00,0.00,0.81}{#1}}
\newcommand{\FunctionTok}[1]{\textcolor[rgb]{0.13,0.29,0.53}{\textbf{#1}}}
\newcommand{\ImportTok}[1]{#1}
\newcommand{\InformationTok}[1]{\textcolor[rgb]{0.56,0.35,0.01}{\textbf{\textit{#1}}}}
\newcommand{\KeywordTok}[1]{\textcolor[rgb]{0.13,0.29,0.53}{\textbf{#1}}}
\newcommand{\NormalTok}[1]{#1}
\newcommand{\OperatorTok}[1]{\textcolor[rgb]{0.81,0.36,0.00}{\textbf{#1}}}
\newcommand{\OtherTok}[1]{\textcolor[rgb]{0.56,0.35,0.01}{#1}}
\newcommand{\PreprocessorTok}[1]{\textcolor[rgb]{0.56,0.35,0.01}{\textit{#1}}}
\newcommand{\RegionMarkerTok}[1]{#1}
\newcommand{\SpecialCharTok}[1]{\textcolor[rgb]{0.81,0.36,0.00}{\textbf{#1}}}
\newcommand{\SpecialStringTok}[1]{\textcolor[rgb]{0.31,0.60,0.02}{#1}}
\newcommand{\StringTok}[1]{\textcolor[rgb]{0.31,0.60,0.02}{#1}}
\newcommand{\VariableTok}[1]{\textcolor[rgb]{0.00,0.00,0.00}{#1}}
\newcommand{\VerbatimStringTok}[1]{\textcolor[rgb]{0.31,0.60,0.02}{#1}}
\newcommand{\WarningTok}[1]{\textcolor[rgb]{0.56,0.35,0.01}{\textbf{\textit{#1}}}}
\usepackage{longtable,booktabs,array}
\usepackage{caption}
\captionsetup[table]{skip=6pt}
\newcounter{none} % for unnumbered tables
\usepackage{calc} % for calculating minipage widths
% Correct order of tables after \paragraph or \subparagraph
\usepackage{etoolbox}
\makeatletter
\patchcmd\longtable{\par}{\if@noskipsec\mbox{}\fi\par}{}{}
\makeatother
\usepackage{graphicx}
\makeatletter
\newsavebox\pandoc@box
\newcommand*\pandocbounded[1]{% scales image to fit in text height/width
  \sbox\pandoc@box{#1}%
  \Gscale@div\@tempa{\textheight}{\dimexpr\ht\pandoc@box+\dp\pandoc@box\relax}%
  \Gscale@div\@tempb{\linewidth}{\wd\pandoc@box}%
  \ifdim\@tempb\p@<\@tempa\p@\let\@tempa\@tempb\fi% select the smaller of both
  \ifdim\@tempa\p@<\p@\scalebox{\@tempa}{\usebox\pandoc@box}%
  \else\usebox{\pandoc@box}%
  \fi%
}
% Set default figure placement to htbp
\def\fps@figure{htbp}
\makeatother
\ifLuaTeX
\usepackage[bidi=basic,shorthands=off]{babel}
\else
\usepackage[bidi=default,shorthands=off]{babel}
\fi
\ifLuaTeX
  \usepackage{selnolig} % disable illegal ligatures
\fi
\setlength{\emergencystretch}{3em} % prevent overfull lines
\providecommand{\tightlist}{%
  \setlength{\itemsep}{0pt}\setlength{\parskip}{0pt}}
\usepackage{amsmath}
\usepackage{amssymb}
\usepackage{booktabs}
\usepackage{booktabs}
\usepackage{longtable}
\usepackage{array}
\usepackage{multirow}
\usepackage{wrapfig}
\usepackage{float}
\usepackage{colortbl}
\usepackage{pdflscape}
\usepackage{tabularx}
\usepackage{xltabular}
\usepackage{threeparttable}
\usepackage{threeparttablex}
\usepackage[normalem]{ulem}
\usepackage{makecell}
\usepackage{xcolor}
\usepackage{bookmark}
\IfFileExists{xurl.sty}{\usepackage{xurl}}{} % add URL line breaks if available
\urlstyle{same}
% fallback for those not using the hyperref driver hyperxmp:
\makeatletter
\@ifundefined{xmpquote}{\newcommand{\xmpquote}[1]{#1}}{}
\makeatother
\hypersetup{
  pdftitle={Modeling Insurance Claims Frequency and Severity},
  pdfauthor={LACTU2150 - Statistical Analysis of Insurance Data},
  pdflang={en},
  hidelinks,
  pdfcreator={LaTeX via pandoc}}

\title{Modeling Insurance Claims Frequency and Severity}
\author{LACTU2150 - Statistical Analysis of Insurance Data}
\date{}

\begin{document}
\maketitle

{
\setcounter{tocdepth}{2}
\tableofcontents
}
\section{Executive Summary}\label{executive-summary}

\section{Introduction}\label{introduction}

\subsection{Background}\label{background}

In Belgium, drivers are required to have Motor Third Party Liability
(MTPL) insurance to ensure that damages caused to third parties will be
covered financially. MTPL insurances are part of non-life insurance
which covers risks to goods, property, and individuals that are not
related to life insurance. These insurances are typically offered on a
yearly basis. The policyholder pays a premium that covers the expected
cost of claims during the policy period and the insurer assumes the risk
of those claims. The advantage that insurances have over individuals is
the concept of \texttt{mutualization}. When risks are pooled together:

\begin{itemize}
\tightlist
\item
  The claims cost is better distributed across years and policyholders,
  turning a punctual high accident cost into a predictable premium
  \footnote{This relies on the assumption that claims are independent
    across policyholders and time periods. This means that we do not
    consider catastrophic events or manufacturer defects that could lead
    to correlated claims.}.
\item
  Average losses are mathematically more predictable due to the law of
  large numbers, reducing the impact of individual high-cost claims.
\end{itemize}

The goal of mutualization is thus to support \texttt{chance\ solidarity}
among policyholders, ensuring them against bad luck. However, not all
policyholders have the same risk profile, and insurers will need to
account for this \texttt{heterogeneity} when pricing premiums. If they
don't, the policyholders with lower risk (technical gains) will end up
subsidizing those with higher risk (technical losses), leading to an
unfair distribution of premiums. They will tend to seek alternative
insurance options better aligned with their risk profile. The
higher-risk policyholders will on the contrary remain as a
\texttt{community\ rating}(uniform premium) is more interesting for them
compared to a risk-adjusted premium. The portfolio will thus tend to
become increasingly composed of higher-risk individuals over time. This
is known as \texttt{adverse\ selection} and is the main risk for
insurers if they fail to properly discriminate between different risk
profiles when pricing premiums.

The ideal for a commercial insurer \footnote{Community rating can be the
  best choice for a state insurer where the re-distribution mechanism is
  desirable and no alternative insurance options are available. Mutual
  insurances are a typical example of such arrangements.} is to
perfectly discriminate between different risk profiles to avoid
\texttt{transfers} between risk classes. They can then provide
\texttt{random\ solidarity} within each risk class without forcing
\texttt{income\ or\ risk\ solidarity} across different risk classes. The
commercial insurer will therefore start with a complex and detailed risk
classification system to get as close as possible to the real risk of
each policyholder. This is what we call \texttt{technical\ premium}
modelling. It uncovers risk transfer between policyholders. It is
however not possible nor desirable to use this complex model as a
commercial model because the premium setting is not transparent enough
and may not be acceptable to policyholders. Therefore, insurers often
simplify the model to balance risk discrimination with transparency and
marketability. We then talk about \texttt{commercial\ premium}
modelling.

The goal of this report is to analyze the Belgian Motor Third Party
Liability (MTPL) insurance portfolio for 1997 to explore the risk
factors and propose a suitable technical and commercial premium model.
As we will explain in the ``Approach'' section of this Introduction
chapter, we will study the frequency and severity of claims separately
to better understand the impact of different risk factors on each
component of the claims process.

\subsection{Portfolio Overview}\label{portfolio-overview}

\subsubsection{Records (Rows)}\label{records-rows}

\begin{Shaded}
\begin{Highlighting}[]
\FunctionTok{setwd}\NormalTok{(}\FunctionTok{paste0}\NormalTok{(}
    \StringTok{"C:/Users/MathieuDemarets/Documents/UCL Certificat/"}\NormalTok{,}
    \StringTok{"CodeRepositories/LACTU2150/project"}
\NormalTok{))}
\FunctionTok{load}\NormalTok{(}\StringTok{"../data/mtpl\_be.rda"}\NormalTok{)}
\NormalTok{data }\OtherTok{\textless{}{-}}\NormalTok{ mtpl\_be}

\NormalTok{n\_unique\_ids }\OtherTok{\textless{}{-}} \FunctionTok{length}\NormalTok{(}\FunctionTok{unique}\NormalTok{(mtpl\_be}\SpecialCharTok{$}\NormalTok{ID))}
\NormalTok{n\_rows }\OtherTok{\textless{}{-}} \FunctionTok{nrow}\NormalTok{(mtpl\_be)}

\NormalTok{index\_training }\OtherTok{\textless{}{-}} \FunctionTok{createDataPartition}\NormalTok{(mtpl\_be}\SpecialCharTok{$}\NormalTok{NCLAIMS, }\AttributeTok{times =} \DecValTok{1}\NormalTok{, }\AttributeTok{p =} \FloatTok{0.8}\NormalTok{, }\AttributeTok{list =} \ConstantTok{FALSE}\NormalTok{) }\CommentTok{\# Note that we make a single partition, cv will be done with the training set during modelling}
\NormalTok{data }\OtherTok{\textless{}{-}}\NormalTok{ mtpl\_be[index\_training, ]}
\NormalTok{holdout }\OtherTok{\textless{}{-}}\NormalTok{ mtpl\_be[}\SpecialCharTok{{-}}\NormalTok{index\_training, ]}
\end{Highlighting}
\end{Shaded}

The first essential step is to get familiar with the MTPL insurance
portfolio for 1997. The full portfolio consists of 163231 individual
policies. We ensured that each record represents a unique policy by
verifying that the number of unique \texttt{ID}s matches the total
number of rows (163231 unique IDs = 163231 rows). Before going any
further, we will split the data into a training set (80\% of the
policies = 130585) and a holdout set (20\% of the policies = 32646) for
model evaluation. This split is made with regards to the number of
claims (\texttt{NCLAIMS}) to ensure that both the training and holdout
sets have a similar distribution of claims. Splitting the data is
necessary for a proper independent evaluation of the models we will
develop. The portfolio overview and data exploration is already part of
the modelling process as it helps us understand the characteristics of
the data and will guide us in selecting appropriate features and
modeling techniques.

\subsubsection{Observations (Columns)}\label{observations-columns}

For each unique policies, we have a number of attributes that can be
classified as:

\begin{itemize}
\tightlist
\item
  \textbf{Policyholder A-Priori Attributes}: age (\texttt{AGEPH}),
  gender (\texttt{SEX}), localisation (\texttt{LAT}, \texttt{LONG},
  \texttt{PC}, and \texttt{TOWN})
\item
  \textbf{Policyholder A-Posteriori Attributes}: bonus-malus score
  (\texttt{BM}) \footnote{A-priori attributes are what you know about
    the policyholder before engaging in the insurance contract. The
    bonus-malus score (\texttt{BM}) is an a-posteriori attribute as it
    reflects the policyholder's claim history and behavior. It is most
    often observed after the policy has been in force for some time and
    is therefore used for experience rating.}
\item
  \textbf{Vehicle Attributes}: fuel type (\texttt{FUEL}), age
  (\texttt{AGEC}), power (\texttt{POWER}), usage (\texttt{USE}), fleet
  status (\texttt{FLEET})
\item
  \textbf{Policy Attributes}: exposure (\texttt{EXP}) and MTPL coverage
  (\texttt{COVERAGE})
\item
  \textbf{Target Variables}: frequency (\texttt{NCLAIMS}), severity
  (\texttt{AVG}), and total claim cost (\texttt{AMOUNT})
\end{itemize}

This classification helps us understand the data from a business
perspective. From a modeling perspective, we will rather differentiate:

\begin{itemize}
\tightlist
\item
  \textbf{Categorical Explanatory Variables}: \texttt{SEX}, \texttt{PC},
  \texttt{TOWN}, \texttt{FUEL}, \texttt{USE}, \texttt{FLEET},
  \texttt{COVERAGE} will be used as factors with eventual merged levels.
\item
  \textbf{Numerical Explanatory Variables}: \texttt{AGEPH},
  \texttt{LAT}, \texttt{LONG}, \texttt{AGEC}, \texttt{POWER},
  \texttt{BM} will be binned and treated as factors or will be used as
  continuous variables depending on the chosen model.
\item
  \textbf{Exposures}: \texttt{EXP} will be used as an offset to account
  for the varying lengths of policy coverage in the modeling process.
\item
  \textbf{Target Variables}: \texttt{NCLAIMS}, \texttt{AVG} will be used
  as the response variables in the modeling process, with
  \texttt{AMOUNT} being derived from them.
\end{itemize}

\paragraph{All Variables}\label{all-variables}

\begin{Shaded}
\begin{Highlighting}[]
\NormalTok{col\_types }\OtherTok{\textless{}{-}} \FunctionTok{sapply}\NormalTok{(data, class)}
\NormalTok{categories }\OtherTok{\textless{}{-}} \FunctionTok{list}\NormalTok{(}
    \AttributeTok{categorical =} \FunctionTok{c}\NormalTok{(}\StringTok{"SEX"}\NormalTok{, }\StringTok{"PC"}\NormalTok{, }\StringTok{"TOWN"}\NormalTok{, }\StringTok{"FUEL"}\NormalTok{, }\StringTok{"USE"}\NormalTok{, }\StringTok{"FLEET"}\NormalTok{, }\StringTok{"COVERAGE"}\NormalTok{),}
    \AttributeTok{numerical =} \FunctionTok{c}\NormalTok{(}\StringTok{"AGEPH"}\NormalTok{, }\StringTok{"LAT"}\NormalTok{, }\StringTok{"LONG"}\NormalTok{, }\StringTok{"AGEC"}\NormalTok{, }\StringTok{"POWER"}\NormalTok{, }\StringTok{"BM"}\NormalTok{),}
    \AttributeTok{exposure =} \FunctionTok{c}\NormalTok{(}\StringTok{"EXP"}\NormalTok{),}
    \AttributeTok{target =} \FunctionTok{c}\NormalTok{(}\StringTok{"NCLAIMS"}\NormalTok{, }\StringTok{"AVG"}\NormalTok{, }\StringTok{"AMOUNT"}\NormalTok{)}
\NormalTok{)}

\NormalTok{all\_vars }\OtherTok{\textless{}{-}} \FunctionTok{data.frame}\NormalTok{(}
    \AttributeTok{Type =}\NormalTok{ col\_types,}
    \AttributeTok{Category =} \FunctionTok{sapply}\NormalTok{(}\FunctionTok{names}\NormalTok{(data), }\ControlFlowTok{function}\NormalTok{(x) \{}
        \ControlFlowTok{if}\NormalTok{ (x }\SpecialCharTok{\%in\%}\NormalTok{ categories}\SpecialCharTok{$}\NormalTok{categorical) \{}
            \StringTok{"Categorical"}
\NormalTok{        \} }\ControlFlowTok{else} \ControlFlowTok{if}\NormalTok{ (x }\SpecialCharTok{\%in\%}\NormalTok{ categories}\SpecialCharTok{$}\NormalTok{numerical) \{}
            \StringTok{"Numerical"}
\NormalTok{        \} }\ControlFlowTok{else} \ControlFlowTok{if}\NormalTok{ (x }\SpecialCharTok{\%in\%}\NormalTok{ categories}\SpecialCharTok{$}\NormalTok{exposure) \{}
            \StringTok{"Exposure"}
\NormalTok{        \} }\ControlFlowTok{else} \ControlFlowTok{if}\NormalTok{ (x }\SpecialCharTok{\%in\%}\NormalTok{ categories}\SpecialCharTok{$}\NormalTok{target) \{}
            \StringTok{"Target"}
\NormalTok{        \} }\ControlFlowTok{else}\NormalTok{ \{}
            \StringTok{"Other"}
\NormalTok{        \}}
\NormalTok{    \}),}
    \AttributeTok{Center =} \FunctionTok{sapply}\NormalTok{(data, }\ControlFlowTok{function}\NormalTok{(x) \{}
        \ControlFlowTok{if}\NormalTok{ (}\FunctionTok{is.numeric}\NormalTok{(x)) \{}
            \FunctionTok{mean}\NormalTok{(x, }\AttributeTok{na.rm =} \ConstantTok{TRUE}\NormalTok{) }\SpecialCharTok{\%\textgreater{}\%} \FunctionTok{round}\NormalTok{(}\DecValTok{4}\NormalTok{)}
\NormalTok{        \} }\ControlFlowTok{else}\NormalTok{ \{}
\NormalTok{            counts }\OtherTok{\textless{}{-}} \FunctionTok{sort}\NormalTok{(}\FunctionTok{table}\NormalTok{(x), }\AttributeTok{decreasing =} \ConstantTok{TRUE}\NormalTok{)}
            \FunctionTok{names}\NormalTok{(counts)[}\DecValTok{1}\NormalTok{]}
\NormalTok{        \}}
\NormalTok{    \}),}
    \AttributeTok{Unique =} \FunctionTok{sapply}\NormalTok{(data, }\ControlFlowTok{function}\NormalTok{(x) }\FunctionTok{length}\NormalTok{(}\FunctionTok{unique}\NormalTok{(x))),}
    \AttributeTok{Missing =} \FunctionTok{sapply}\NormalTok{(data, }\ControlFlowTok{function}\NormalTok{(x) }\FunctionTok{sum}\NormalTok{(}\FunctionTok{is.na}\NormalTok{(x)))}
\NormalTok{) }\SpecialCharTok{\%\textgreater{}\%} \FunctionTok{arrange}\NormalTok{(Category)}

\NormalTok{all\_vars }\SpecialCharTok{\%\textgreater{}\%} \FunctionTok{kable}\NormalTok{(}
    \AttributeTok{caption =} \FunctionTok{paste0}\NormalTok{(}\StringTok{"Overview of all variables in the MTPL dataset ("}\NormalTok{, }\FunctionTok{nrow}\NormalTok{(data), }\StringTok{" rows, "}\NormalTok{, }\FunctionTok{ncol}\NormalTok{(data), }\StringTok{" columns)"}\NormalTok{),}
    \AttributeTok{booktabs =} \ConstantTok{TRUE}
\NormalTok{)}
\end{Highlighting}
\end{Shaded}

\begin{longtable}[]{@{}llllrr@{}}
\caption{Overview of all variables in the MTPL dataset (130585 rows, 18
columns)}\tabularnewline
\toprule\noalign{}
& Type & Category & Center & Unique & Missing \\
\midrule\noalign{}
\endfirsthead
\toprule\noalign{}
& Type & Category & Center & Unique & Missing \\
\midrule\noalign{}
\endhead
\bottomrule\noalign{}
\endlastfoot
COVERAGE & character & Categorical & TPL & 3 & 0 \\
FUEL & character & Categorical & gasoline & 2 & 0 \\
USE & character & Categorical & private & 2 & 0 \\
FLEET & character & Categorical & N & 2 & 0 \\
SEX & character & Categorical & male & 2 & 0 \\
PC & integer & Categorical & 5275.3966 & 583 & 0 \\
TOWN & character & Categorical & CHARLERO & 578 & 0 \\
EXP & numeric & Exposure & 0.8899 & 365 & 0 \\
AGEPH & integer & Numerical & 47.01 & 78 & 0 \\
BM & integer & Numerical & 3.2671 & 23 & 0 \\
AGEC & integer & Numerical & 7.3796 & 49 & 0 \\
POWER & integer & Numerical & 56.0367 & 174 & 0 \\
LONG & numeric & Numerical & 4.4065 & 583 & 0 \\
LAT & numeric & Numerical & 50.7588 & 583 & 0 \\
ID & integer & Other & 81699.2272 & 130585 & 0 \\
NCLAIMS & integer & Target & 0.1234 & 6 & 0 \\
AMOUNT & numeric & Target & 208.8667 & 6761 & 0 \\
AVG & numeric & Target & 1698.6839 & 6783 & 116015 \\
\end{longtable}

From this basic overview, we can note that:

\begin{itemize}
\tightlist
\item
  There is no missing data except for \texttt{AVG}. As we will see
  later, this is expected as an average on zero-claim policies would be
  undefined.
\item
  \texttt{ID} is indeed a unique identifier as there are as many unique
  values (130585) as rows in the dataset.
\item
  Categorical variables are currently represented as character and
  integers, they will be later converted to factor variables.
\item
  \texttt{PC} and \texttt{TOWN} have too many levels (583 and 578
  respectively) and will have to be transformed to be usable. They also
  seem to indicate the same geographical information.
\item
  \texttt{EXP} has 365 unique values which indicates that claims
  happened over all possible days in a year.
\item
  The average claim \texttt{AMOUNT} (208.8667) is not equal to the
  average \texttt{AVG} (1698.6839). We will explore this further next
  because the Center values should be adjusted for the exposure
  \texttt{EXP}.
\end{itemize}

\paragraph{\texorpdfstring{Deep Dive into Exposure
(\texttt{EXP})}{Deep Dive into Exposure (EXP)}}\label{deep-dive-into-exposure-exp}

The \texttt{exposure} is an important concept in frequency modelling. In
latter sections, we will try to model the number of claims over a given
period, this is the \texttt{claim\ frequency}. The given period, which
we call the exposure, is by default one year as non-life insurance
policies are typically renewed annually. However, a policy can be
stopped or canceled before the end of the year, resulting in an exposure
less than one. It is important to keep in mind that exposure is the unit
of risk and not the number of policies. Indeed, if we take the number of
policies instead of the exposure, we will always underestimate the true
claim frequency (as the exposure is between 0 and 1).

If we denote the number of claims by \texttt{NCLAIMS}, the exposure by
\texttt{EXP}, then the claim frequency
\texttt{\$\textbackslash{}hat\{\textbackslash{}lambda\}\$} (the expected
number of claims per unit of exposure) is given by:

\[\hat{\lambda} = \frac{\sum_{i=1}^{n} NCLAIMS_i}{\sum_{i=1}^{n} EXP_i} \leq \frac{\sum_{i=1}^{n} NCLAIMS_i}{n}\]

In the dataset, about 0.23 of the policies have an exposure less than 1,
it is therefore important to account for the actual exposure when
estimating the claim frequency. We also saw in the description of all
variables that the average exposure is 0.8899, which means that the
actual claim frequency will be underestimated by this factor if we were
to ignore the exposure and take the naive mean of \texttt{NCLAIMS}.

\begin{Shaded}
\begin{Highlighting}[]
\NormalTok{correct\_lambda }\OtherTok{\textless{}{-}} \FunctionTok{sum}\NormalTok{(data}\SpecialCharTok{$}\NormalTok{NCLAIMS) }\SpecialCharTok{/} \FunctionTok{sum}\NormalTok{(data}\SpecialCharTok{$}\NormalTok{EXP)}
\NormalTok{naive\_lambda }\OtherTok{\textless{}{-}} \FunctionTok{mean}\NormalTok{(data}\SpecialCharTok{$}\NormalTok{NCLAIMS)}
\NormalTok{se\_mean }\OtherTok{\textless{}{-}} \FunctionTok{sqrt}\NormalTok{(correct\_lambda }\SpecialCharTok{/} \FunctionTok{sum}\NormalTok{(data}\SpecialCharTok{$}\NormalTok{EXP))}
\NormalTok{ic\_99 }\OtherTok{\textless{}{-}}\NormalTok{ correct\_lambda }\SpecialCharTok{+} \FunctionTok{c}\NormalTok{(}\SpecialCharTok{{-}}\DecValTok{1}\NormalTok{, }\DecValTok{1}\NormalTok{) }\SpecialCharTok{*} \FunctionTok{qnorm}\NormalTok{(}\FloatTok{0.995}\NormalTok{) }\SpecialCharTok{*}\NormalTok{ se\_mean}

\FunctionTok{cat}\NormalTok{(}\FunctionTok{sprintf}\NormalTok{(}
    \StringTok{"Correct mean claim frequency is \%.3f\%\% (99\%\% CI: [\%.3f\%\%, \%.3f\%\%]) and not the naive mean \%.3f\%\%"}\NormalTok{,}
\NormalTok{    correct\_lambda }\SpecialCharTok{*} \DecValTok{100}\NormalTok{,}
\NormalTok{    ic\_99[}\DecValTok{1}\NormalTok{] }\SpecialCharTok{*} \DecValTok{100}\NormalTok{,}
\NormalTok{    ic\_99[}\DecValTok{2}\NormalTok{] }\SpecialCharTok{*} \DecValTok{100}\NormalTok{,}
\NormalTok{    naive\_lambda }\SpecialCharTok{*} \DecValTok{100}
\NormalTok{))}
\end{Highlighting}
\end{Shaded}

\begin{verbatim}
Correct mean claim frequency is 13.862% (99% CI: [13.581%, 14.144%]) and not the naive mean 12.336%
\end{verbatim}

We confirm numerically that the naive mean of \texttt{NCLAIMS}
underestimates the true claim frequency as it falls short of the 99\%
confidence interval\footnote{If we assume a poisson distribution for the
  number of claims, we'll show in the data exploration section that the
  standard error of the claim frequency can be estimated as \$
  \sqrt{\frac{\hat{\lambda}}{\sum_{i=1}^{n} EXP_i}} \$.} for the correct
mean claim frequency. It is therefore important to account for the
exposure through the offset\footnote{In frequency modelling, the
  exposure is typically included as an offset in the model to adjust for
  the varying lengths of policy coverage. This ensures that the
  estimated claim frequency accurately reflects the risk per unit of
  exposure rather than per policy.} in the frequency modelling.

\paragraph{\texorpdfstring{Deep Dive into Target Variables
(\texttt{NCLAIMS}, \texttt{AVG}, \texttt{AMOUNT}
)}{Deep Dive into Target Variables (NCLAIMS, AVG, AMOUNT )}}\label{deep-dive-into-target-variables-nclaims-avg-amount}

The expected cost of claims can be decomposed into the expected
frequency of claims and the expected severity of claims. The expected
total claim cost \texttt{\$\textbackslash{}pi\$} can be expressed in
terms of the expected frequency \texttt{\$\textbackslash{}lambda\$} and
the expected severity \texttt{\$\textbackslash{}mu\$} as:

\[\pi(x) = \lambda(x) \cdot \mu(x)\]

This split is useful to isolate the effects of explanatory variables on
the likehood of claims occurring (\texttt{\textbackslash{}lambda(x)})
from the impact of that claim occurring
(\texttt{\textbackslash{}mu(x)})\footnote{Similar decomposition are done
  in business with \texttt{Impact-Likelihood} matrices or in credit risk
  modelling with \texttt{Probability\ of\ Default} and
  \texttt{Loss\ Given\ Default}.}. The mitigation strategies are also
different depending on whether we aim to reduce the frequency of claims
(prevention and underwriting) or the severity of claims (reinsurance and
deductibles).

For our training dataset\footnote{Remember that the global portfolio was
  split to keep a holdout dataset of 20\% for validation.}, the
empirical expected total claim cost can be decomposed into the product
of the empirical expected claim frequency and the empirical expected
claim severity:

\[\frac{\sum_{i=1}^{n} AMOUNT_i}{\sum_{i=1}^{n} EXP_i} = \frac{\sum_{i=1}^{n} NCLAIMS_i}{\sum_{i=1}^{n} EXP_i} \cdot \frac{\sum_{i=1}^{n} AMOUNT_i}{\sum_{i=1}^{n} NCLAIMS_i}
\text{where } \frac{\sum_{i=1}^{n} AMOUNT_i}{\sum_{i=1}^{n} NCLAIMS_i} = AVG_i\]

\begin{Shaded}
\begin{Highlighting}[]
\NormalTok{sum\_exp }\OtherTok{\textless{}{-}} \FunctionTok{sum}\NormalTok{(data}\SpecialCharTok{$}\NormalTok{EXP)}
\NormalTok{sum\_nclaims }\OtherTok{\textless{}{-}} \FunctionTok{sum}\NormalTok{(data}\SpecialCharTok{$}\NormalTok{NCLAIMS)}
\NormalTok{sum\_amount }\OtherTok{\textless{}{-}} \FunctionTok{sum}\NormalTok{(data}\SpecialCharTok{$}\NormalTok{AMOUNT)}

\FunctionTok{cat}\NormalTok{(}\FunctionTok{sprintf}\NormalTok{(}\StringTok{"Expected Total Claim Cost: \%.2f = \%.2f * \%.2f"}\NormalTok{, sum\_amount }\SpecialCharTok{/}\NormalTok{ sum\_exp, sum\_nclaims }\SpecialCharTok{/}\NormalTok{ sum\_exp, sum\_amount }\SpecialCharTok{/}\NormalTok{ sum\_nclaims))}
\end{Highlighting}
\end{Shaded}

\begin{verbatim}
Expected Total Claim Cost: 234.71 = 0.14 * 1693.14
\end{verbatim}

\begin{Shaded}
\begin{Highlighting}[]
\FunctionTok{cat}\NormalTok{(}\FunctionTok{sprintf}\NormalTok{(}\StringTok{"With sum\_amount = \%.2f, sum\_nclaims = \%.2f, sum\_exp = \%.2f"}\NormalTok{, sum\_amount, sum\_nclaims, sum\_exp))}
\end{Highlighting}
\end{Shaded}

\begin{verbatim}
With sum_amount = 27274852.52, sum_nclaims = 16109.00, sum_exp = 116206.43
\end{verbatim}

This is not the full picture as an aggregated view hides the
heterogeneity within the portfolio. Looking at the distribution of
claims for frequency and the distribution of average claim cost for
severity provides a view of this heterogeneity that we will try to
capture in the rest of the report.

\begin{Shaded}
\begin{Highlighting}[]
\NormalTok{p\_n\_claims }\OtherTok{\textless{}{-}} \FunctionTok{ggplot}\NormalTok{(data, }\FunctionTok{aes}\NormalTok{(}\AttributeTok{x =}\NormalTok{ NCLAIMS, }\AttributeTok{fill =}\NormalTok{ NCLAIMS }\SpecialCharTok{==} \DecValTok{0}\NormalTok{)) }\SpecialCharTok{+}
    \FunctionTok{geom\_bar}\NormalTok{() }\SpecialCharTok{+}
    \FunctionTok{scale\_fill\_manual}\NormalTok{(}\AttributeTok{values =} \FunctionTok{c}\NormalTok{(}\StringTok{"TRUE"} \OtherTok{=} \StringTok{"lightgrey"}\NormalTok{, }\StringTok{"FALSE"} \OtherTok{=} \StringTok{"steelblue"}\NormalTok{)) }\SpecialCharTok{+}
    \FunctionTok{labs}\NormalTok{(}
        \AttributeTok{title =} \StringTok{"Frequency {-} Number of Claims"}\NormalTok{,}
        \AttributeTok{y =} \FunctionTok{sprintf}\NormalTok{(}\StringTok{"Count (All: \%d)"}\NormalTok{, }\FunctionTok{nrow}\NormalTok{(data)),}
        \AttributeTok{x =} \StringTok{"Number of Claims"}
\NormalTok{    ) }\SpecialCharTok{+}
    \FunctionTok{theme\_minimal}\NormalTok{()}

\NormalTok{avg\_plot }\OtherTok{\textless{}{-}}\NormalTok{ data }\SpecialCharTok{\%\textgreater{}\%} \FunctionTok{filter}\NormalTok{(}\SpecialCharTok{!}\FunctionTok{is.na}\NormalTok{(AVG))}

\NormalTok{p\_avg }\OtherTok{\textless{}{-}} \FunctionTok{ggplot}\NormalTok{(avg\_plot, }\FunctionTok{aes}\NormalTok{(}\AttributeTok{x =} \FunctionTok{log}\NormalTok{(AVG), }\AttributeTok{fill =}\NormalTok{ NCLAIMS }\SpecialCharTok{==} \DecValTok{0}\NormalTok{)) }\SpecialCharTok{+}
    \FunctionTok{geom\_histogram}\NormalTok{(}\AttributeTok{bins =} \DecValTok{30}\NormalTok{) }\SpecialCharTok{+}
    \FunctionTok{scale\_fill\_manual}\NormalTok{(}\AttributeTok{values =} \FunctionTok{c}\NormalTok{(}\StringTok{"TRUE"} \OtherTok{=} \StringTok{"lightgrey"}\NormalTok{, }\StringTok{"FALSE"} \OtherTok{=} \StringTok{"steelblue"}\NormalTok{)) }\SpecialCharTok{+}
    \FunctionTok{labs}\NormalTok{(}
        \AttributeTok{title =} \StringTok{"Severity {-} Average Claim Cost"}\NormalTok{,}
        \AttributeTok{x =} \StringTok{"Log of Average Claim Cost"}\NormalTok{,}
        \AttributeTok{y =} \FunctionTok{sprintf}\NormalTok{(}\StringTok{"Count (All: \%d)"}\NormalTok{, }\FunctionTok{nrow}\NormalTok{(avg\_plot))}
\NormalTok{    ) }\SpecialCharTok{+}
    \FunctionTok{geom\_vline}\NormalTok{(}\FunctionTok{aes}\NormalTok{(}\AttributeTok{xintercept =} \DecValTok{5}\NormalTok{), }\AttributeTok{linetype =} \StringTok{"dashed"}\NormalTok{, }\AttributeTok{color =} \StringTok{"orange"}\NormalTok{) }\SpecialCharTok{+}
    \FunctionTok{geom\_vline}\NormalTok{(}\FunctionTok{aes}\NormalTok{(}\AttributeTok{xintercept =} \FloatTok{7.5}\NormalTok{), }\AttributeTok{linetype =} \StringTok{"dashed"}\NormalTok{, }\AttributeTok{color =} \StringTok{"red"}\NormalTok{) }\SpecialCharTok{+}
    \FunctionTok{geom\_vline}\NormalTok{(}\FunctionTok{aes}\NormalTok{(}\AttributeTok{xintercept =} \DecValTok{10}\NormalTok{), }\AttributeTok{linetype =} \StringTok{"dashed"}\NormalTok{, }\AttributeTok{color =} \StringTok{"purple"}\NormalTok{) }\SpecialCharTok{+}
    \FunctionTok{theme\_minimal}\NormalTok{()}

\NormalTok{p\_n\_claims }\SpecialCharTok{|}\NormalTok{ p\_avg}
\end{Highlighting}
\end{Shaded}

\begin{center}\includegraphics{project_MD_files/figure-latex/heterogeneity-1} \end{center}

For the number of claims, we observe a clear imbalance in the data with
88.84\% of the policies having no claims. This is quite typical in
insurance data, where many policies have no claims or very low average
claim costs. We could tackle this data imbalance with dedicated
techniques such as:

\begin{itemize}
\tightlist
\item
  Synthetic oversampling techniques (e.g., SMOTE) create synthetic
  samples for the minority class to balance the dataset.
\item
  Hurdle models can be employed to separately model the zero and
  non-zero counts, effectively handling the excess zeros in the data and
  focusing on the non-zero records in the second part of the model.
\end{itemize}

These techniques will however break the adequation between the observed
frequency and the modeled frequency which is an important benefit of the
modeling approach that we will present in the next section. We will
therefore not apply these techniques in this analysis.

For policies with claims, we look at the log of the average claim cost
to better understand the distribution of claim severity in terms of
magnitude. We observe two peaks and a long right tail. This shows that
the average claim cost we computed above actually hides very different
types of claims:

\begin{itemize}
\tightlist
\item
  First peak corresponds to small dents and repairs: the log around 5
  represents claims of approximately 148 euros
\item
  Second peak corresponds to more significant repairs: the log around
  7.5 represents claims of approximately 1808 euros
\item
  Right tail corresponds to way bigger damages: the log around 10
  represents claims of approximately 22026 euros
\end{itemize}

Even if they are not so frequent, the larger claims in the right tail
can have a significant impact on the company profitability. A future
report could use \texttt{Extreme\ Value\ Theory} to better model and
understand these rare but impactful claims. We will not apply this
technique in the current analysis.

\subsection{Approach}\label{approach}

The goal of this report is to estimate the expected claim cost for
different policies, taking into account both the frequency and severity
of claims. The

\section{Data Preparation}\label{data-preparation}

\subsection{Data Transformation}\label{data-transformation}

\subsubsection{Categorical Variables}\label{categorical-variables}

As concluded at the end of the last section, we need to transform some
character variables to factor variables. We will also make sure that the
most frequent level (Center in the previous overview) is set as the
reference level for unordered factor variables.

\begin{Shaded}
\begin{Highlighting}[]
\ControlFlowTok{for}\NormalTok{ (col }\ControlFlowTok{in} \FunctionTok{c}\NormalTok{(}\StringTok{"COVERAGE"}\NormalTok{, }\StringTok{"FUEL"}\NormalTok{, }\StringTok{"USE"}\NormalTok{, }\StringTok{"FLEET"}\NormalTok{, }\StringTok{"SEX"}\NormalTok{, }\StringTok{"TOWN"}\NormalTok{, }\StringTok{"PC"}\NormalTok{)) \{}
\NormalTok{    data[[col]] }\OtherTok{\textless{}{-}} \FunctionTok{as.factor}\NormalTok{(data[[col]]) }\SpecialCharTok{\%\textgreater{}\%} \FunctionTok{fct\_infreq}\NormalTok{()}
\NormalTok{\}}
\end{Highlighting}
\end{Shaded}

We can then assert that the categorical variables (ordered or unordered)
have the correct reference levels. Note that our function also checks
for ordered factor variables, ensuring that the reference level is the
lowest level in this case. This does not happen in our dataset. The
variable \texttt{BM} could be considered as unordered as we know that a
higher \texttt{BM} score indicates a more risky driver but we don't know
if the jumps between the levels are constant. We decided to treat it as
numerical because the numeric variable will be binned with a weight of
evidence (WoE) approach and it will yield more meaningful bins than
simple singleton bins.

\begin{Shaded}
\begin{Highlighting}[]
\NormalTok{assert\_ref\_level }\OtherTok{\textless{}{-}} \ControlFlowTok{function}\NormalTok{(data, cols) \{}
    \ControlFlowTok{for}\NormalTok{ (col }\ControlFlowTok{in}\NormalTok{ cols) \{}
        \ControlFlowTok{if}\NormalTok{ (}\SpecialCharTok{!}\FunctionTok{is.factor}\NormalTok{(data[[col]])) }\FunctionTok{stop}\NormalTok{(}\FunctionTok{paste}\NormalTok{(col, }\StringTok{"is not a factor variable"}\NormalTok{))}
        \ControlFlowTok{if}\NormalTok{ (}\FunctionTok{is.ordered}\NormalTok{(data[[col]])) \{}
            \ControlFlowTok{if}\NormalTok{ (}\FunctionTok{as.numeric}\NormalTok{(}\FunctionTok{levels}\NormalTok{(data[[col]]))[}\DecValTok{1}\NormalTok{] }\SpecialCharTok{!=} \FunctionTok{min}\NormalTok{(}\FunctionTok{as.numeric}\NormalTok{(}\FunctionTok{as.character}\NormalTok{(data[[col]])), }\AttributeTok{na.rm =} \ConstantTok{TRUE}\NormalTok{)) \{}
            \CommentTok{\# Pay attention to turn labels as numeric and not the factor internal codes which are just integers starting from 1}
                \FunctionTok{stop}\NormalTok{(}\FunctionTok{paste}\NormalTok{(}\StringTok{"Reference level for"}\NormalTok{, col, }\StringTok{"is not the lowest level"}\NormalTok{))}
\NormalTok{            \}}
\NormalTok{        \} }\ControlFlowTok{else}\NormalTok{ \{}
            \ControlFlowTok{if}\NormalTok{ (}
                \FunctionTok{levels}\NormalTok{(data[[col]])[}\DecValTok{1}\NormalTok{] }\SpecialCharTok{!=}
                    \FunctionTok{fct\_count}\NormalTok{(data[[col]], }\AttributeTok{sort =} \ConstantTok{TRUE}\NormalTok{)}\SpecialCharTok{$}\NormalTok{f[}\DecValTok{1}\NormalTok{]}
\NormalTok{            ) \{}
                \FunctionTok{stop}\NormalTok{(}
                    \FunctionTok{paste}\NormalTok{(}
                        \StringTok{"Reference level for"}\NormalTok{, col,}
                        \StringTok{"is not the most frequent level"}
\NormalTok{                    )}
\NormalTok{                )}
\NormalTok{            \}}
\NormalTok{    \}}
\NormalTok{    \}}
    \FunctionTok{cat}\NormalTok{(}\StringTok{"Reference levels for"}\NormalTok{, }\FunctionTok{paste}\NormalTok{(cols, }\AttributeTok{collapse =} \StringTok{", "}\NormalTok{), }\StringTok{"are correctly set}\SpecialCharTok{\textbackslash{}n}\StringTok{"}\NormalTok{)}
\NormalTok{\}       }
\FunctionTok{assert\_ref\_level}\NormalTok{(data,  }\FunctionTok{c}\NormalTok{(}\StringTok{"COVERAGE"}\NormalTok{, }\StringTok{"FUEL"}\NormalTok{, }\StringTok{"USE"}\NormalTok{, }\StringTok{"FLEET"}\NormalTok{, }\StringTok{"SEX"}\NormalTok{, }\StringTok{"TOWN"}\NormalTok{, }\StringTok{"PC"}\NormalTok{))}
\end{Highlighting}
\end{Shaded}

\begin{verbatim}
Reference levels for COVERAGE, FUEL, USE, FLEET, SEX, TOWN, PC are correctly set
\end{verbatim}

\subsubsection{Numerical Variables
Binning}\label{numerical-variables-binning}

For the initial \texttt{GLM} modeling, the numerical variables can be
used but this assumes that:

\begin{itemize}
\tightlist
\item
  The effect of one increment of the numerical variable on the target is
  assumed to be linear (age 18 to 19 has the same effect as age 50 to 51
  which may not be realistic as experience marginal returns are likely
  diminishing).
\item
  The linearity assumption may not hold for all numerical variables, and
  discretizing them into bins can help capture non-linear relationships
  with the target.
\end{itemize}

For these reasons, we will discretize the continuous numerical variables
into bins. Defining optimal binning for a numerical variable is a
complex and important step. There are a set of naive approaches that can
be used:

\begin{itemize}
\tightlist
\item
  Singleton bins: Each unique value of the numerical variable is treated
  as a separate bin.
\item
  Equal-width bins: The range of the variable is divided into intervals
  of equal width.
\item
  Equal-frequency bins: Quantiles of the variable are used to create
  bins with approximately the same number of observations.
\end{itemize}

For this project, we will try to implement a more complex approach based
on weight of evidence (WoE) to automatically determine a suitable
binning for numerical variables. The idea of WoE binning is to group
values based on how strongly they are associated with the target
variable (binary outcome encoded as good if claim = 0 and bad if claim
\textgreater= 1):

\[\text{WoE} = \ln \frac{\text{proportion of Good Outcomes in the bin}}{\text{proportion of Bad Outcomes in the bin}}\]

Based on the WoE values, we can also compute the information value (IV)
for each variable, which measures the predictive power of the variable
with respect to the target. The IV is calculated as:

\[\text{IV} = \sum (\text{proportion of Good Outcomes in the bin} - \text{proportion of Bad Outcomes in the bin}) \times \text{WoE}\]

The IV value helps to assess the predictive power of a variable for the
binning we have chosen. The idea of the algorithm will be to:

\begin{itemize}
\tightlist
\item
  Start with an adapted quantile-based binning as a baseline with a
  maximum number of bins specified.
\item
  Iteratively merge bins based on bins with minimal information value
  (IV) contribution and assess total IV after each merge.
\item
  Find the elbow point where the marginal contribution of one additional
  bin to the total IV decreases significantly.
\end{itemize}

We will take that elbow point and its associated bins boundaries for the
final binning of the variable. We will detail each step of the algorithm
with an example using the \texttt{AGEPH} variable.

\paragraph{Baseline Quantile Binning}\label{baseline-quantile-binning}

Our function \texttt{qbinning} is a simple algorithm that will start
from the maximum number of bins specified by \texttt{max\_bins} and
iteratively reduce the number of bins until each bin contains at least
one claim. This will ensure that the weight of evidence (WoE) can be
calculated for each bin, as each bin will have at least one claim. We'll
then avoid having an undefined WoE.

\begin{Shaded}
\begin{Highlighting}[]
\NormalTok{qbinning }\OtherTok{\textless{}{-}} \ControlFlowTok{function}\NormalTok{(data, x, y, }\AttributeTok{max\_bins =} \DecValTok{10}\NormalTok{) \{}
    \StringTok{"Quantile{-}based binning for numerical variable x with respect to binary target y"}
    \ControlFlowTok{if}\NormalTok{ (}\FunctionTok{sum}\NormalTok{(data[[y]] }\SpecialCharTok{\textgreater{}=} \DecValTok{1}\NormalTok{) }\SpecialCharTok{==} \DecValTok{0}\NormalTok{) }\FunctionTok{stop}\NormalTok{(}\StringTok{"No claims in the dataset"}\NormalTok{)}
    \ControlFlowTok{if}\NormalTok{ (}\FunctionTok{sum}\NormalTok{(data[[y]] }\SpecialCharTok{==} \DecValTok{0}\NormalTok{) }\SpecialCharTok{==} \DecValTok{0}\NormalTok{) }\FunctionTok{stop}\NormalTok{(}\StringTok{"No non{-}claims in the dataset"}\NormalTok{)}
    \ControlFlowTok{if}\NormalTok{ (}\SpecialCharTok{!}\FunctionTok{is.numeric}\NormalTok{(data[[x]])) }\FunctionTok{stop}\NormalTok{(}\StringTok{"Variable x must be numeric"}\NormalTok{)}
\NormalTok{    raw\_data }\OtherTok{\textless{}{-}}\NormalTok{ data}
\NormalTok{    temp\_n }\OtherTok{\textless{}{-}}\NormalTok{ max\_bins}
\NormalTok{    temp\_binning }\OtherTok{\textless{}{-}} \ConstantTok{NULL}
\NormalTok{    baseline\_bins }\OtherTok{\textless{}{-}} \ConstantTok{NULL}
    \ControlFlowTok{while}\NormalTok{ (temp\_n }\SpecialCharTok{\textgreater{}} \DecValTok{1}\NormalTok{) \{}
\NormalTok{        temp\_binning }\OtherTok{\textless{}{-}} \FunctionTok{unique}\NormalTok{(}
            \FunctionTok{quantile}\NormalTok{(raw\_data[[x]], }\AttributeTok{probs =} \FunctionTok{seq}\NormalTok{(}\DecValTok{0}\NormalTok{, }\DecValTok{1}\NormalTok{, }\AttributeTok{length.out =}\NormalTok{ temp\_n }\SpecialCharTok{+} \DecValTok{1}\NormalTok{), }\AttributeTok{na.rm =} \ConstantTok{TRUE}\NormalTok{)}
\NormalTok{        )}
        \CommentTok{\# check if there is at least one claim in each bin}
\NormalTok{        bin\_claims }\OtherTok{\textless{}{-}} \FunctionTok{tapply}\NormalTok{(}
\NormalTok{            raw\_data[[y]],}
            \FunctionTok{cut}\NormalTok{(raw\_data[[x]], temp\_binning, }\AttributeTok{include.lowest =} \ConstantTok{TRUE}\NormalTok{),}
            \ControlFlowTok{function}\NormalTok{(z) }\FunctionTok{sum}\NormalTok{(z }\SpecialCharTok{\textgreater{}=} \DecValTok{1}\NormalTok{))}
        \ControlFlowTok{if}\NormalTok{ (}\FunctionTok{all}\NormalTok{(bin\_claims }\SpecialCharTok{\textgreater{}} \DecValTok{0}\NormalTok{)) \{}
\NormalTok{            baseline\_bins }\OtherTok{\textless{}{-}}\NormalTok{ temp\_binning}
            \ControlFlowTok{break}
\NormalTok{        \}}
\NormalTok{        temp\_n }\OtherTok{\textless{}{-}}\NormalTok{ temp\_n }\SpecialCharTok{{-}} \DecValTok{1}
\NormalTok{    \}}
\NormalTok{    baseline\_bins}
\NormalTok{\}}
\end{Highlighting}
\end{Shaded}

\begin{Shaded}
\begin{Highlighting}[]
\NormalTok{bins\_ageph }\OtherTok{\textless{}{-}} \FunctionTok{qbinning}\NormalTok{(data, }\AttributeTok{x =} \StringTok{"AGEPH"}\NormalTok{, }\AttributeTok{y =} \StringTok{"NCLAIMS"}\NormalTok{, }\AttributeTok{max\_bins =} \DecValTok{30}\NormalTok{)}
\NormalTok{counts }\OtherTok{\textless{}{-}} \FunctionTok{tapply}\NormalTok{(}
\NormalTok{    data[[}\StringTok{"NCLAIMS"}\NormalTok{]],}
    \FunctionTok{cut}\NormalTok{(data[[}\StringTok{"AGEPH"}\NormalTok{]], bins\_ageph, }\AttributeTok{include.lowest =} \ConstantTok{TRUE}\NormalTok{),}
    \ControlFlowTok{function}\NormalTok{(z) }\FunctionTok{sum}\NormalTok{(z }\SpecialCharTok{\textgreater{}=} \DecValTok{1}\NormalTok{)}
\NormalTok{)}

\NormalTok{df\_bins }\OtherTok{\textless{}{-}} \FunctionTok{data.frame}\NormalTok{(}
    \AttributeTok{bin =} \FunctionTok{factor}\NormalTok{(}\FunctionTok{names}\NormalTok{(counts), }\AttributeTok{levels =} \FunctionTok{names}\NormalTok{(counts)),}
    \AttributeTok{count =} \FunctionTok{as.numeric}\NormalTok{(counts)}
\NormalTok{)}

\NormalTok{p }\OtherTok{\textless{}{-}} \FunctionTok{ggplot}\NormalTok{(df\_bins, }\FunctionTok{aes}\NormalTok{(}\AttributeTok{x =}\NormalTok{ bin, }\AttributeTok{y =}\NormalTok{ count)) }\SpecialCharTok{+}
\FunctionTok{ggtitle}\NormalTok{(}\StringTok{"Baseline Quantile Binning (AGEPH)"}\NormalTok{) }\SpecialCharTok{+}
\FunctionTok{geom\_col}\NormalTok{(}\AttributeTok{fill =} \StringTok{"steelblue"}\NormalTok{) }\SpecialCharTok{+}
\FunctionTok{labs}\NormalTok{(}\AttributeTok{x =} \StringTok{"AGEPH baseline quantile bin"}\NormalTok{, }\AttributeTok{y =} \StringTok{"Policies with at least one claim"}\NormalTok{) }\SpecialCharTok{+}
\FunctionTok{theme\_minimal}\NormalTok{() }\SpecialCharTok{+}
\FunctionTok{theme}\NormalTok{(}\AttributeTok{axis.text.x =} \FunctionTok{element\_text}\NormalTok{(}\AttributeTok{angle =} \DecValTok{60}\NormalTok{, }\AttributeTok{hjust =} \DecValTok{1}\NormalTok{))}
\FunctionTok{ggplotly}\NormalTok{(p)}
\end{Highlighting}
\end{Shaded}

\begin{center}\includegraphics{project_MD_files/figure-latex/example_quantile_binning-1} \end{center}

Starting with 30 quantile bins, we apply our tweaked function
\texttt{qbinning} to the \texttt{AGEPH} variable and find these 30 bins
with some policies having at least one claim.

\paragraph{Weight of Evidence (WoE) and Information Value (IV)
Calculation}\label{weight-of-evidence-woe-and-information-value-iv-calculation}

Based on the formula presented above, the weight of evidence (WoE) for
each bin can be calculated as the natural logarithm of:

\begin{itemize}
\tightlist
\item
  \(\text{proportion of Good Outcomes in the bin} = \frac{\text{Number of Good Outcomes in the bin}}{\text{Total Number of Good Outcomes}}\)
\item
  \(\text{proportion of Bad Outcomes in the bin} = \frac{\text{Number of Bad Outcomes in the bin}}{\text{Total Number of Bad Outcomes}}\)
\end{itemize}

If the bin has as much good (no claim) as bad (at least one claim)
outcomes, the WoE will be zero as there is very little predictive power
in that bin. If the bin has more good outcomes than bad outcomes, the
WoE will be positive, indicating that the bin is associated with a lower
risk. Finally, if the bin has more bad outcomes than good outcomes, the
WoE will be negative, indicating that the bin is associated with a
higher risk. The IV quantifies the overall predictive power based on the
differences in the distribution of good and bad outcomes and the
corresponding WoE values.

\begin{Shaded}
\begin{Highlighting}[]
\NormalTok{woe\_calc }\OtherTok{\textless{}{-}} \ControlFlowTok{function}\NormalTok{(data, x, y, bins) \{}
    \StringTok{"Calculate Weight of Evidence (WoE) for each bin"}
\NormalTok{    woe\_df }\OtherTok{\textless{}{-}}\NormalTok{ data}
\NormalTok{    binned\_woe }\OtherTok{\textless{}{-}} \FunctionTok{tapply}\NormalTok{(}
\NormalTok{        woe\_df[[y]],}
        \FunctionTok{cut}\NormalTok{(woe\_df[[x]], bins, }\AttributeTok{include.lowest =} \ConstantTok{TRUE}\NormalTok{),}
        \ControlFlowTok{function}\NormalTok{(z) }\FunctionTok{c}\NormalTok{(}\AttributeTok{good =} \FunctionTok{sum}\NormalTok{(z }\SpecialCharTok{==} \DecValTok{0}\NormalTok{), }\AttributeTok{bad =} \FunctionTok{sum}\NormalTok{(z }\SpecialCharTok{\textgreater{}=} \DecValTok{1}\NormalTok{))}
\NormalTok{    )}
\NormalTok{    binned\_woe\_df }\OtherTok{\textless{}{-}} \FunctionTok{data.frame}\NormalTok{(}
        \AttributeTok{bin =} \FunctionTok{factor}\NormalTok{(}\FunctionTok{names}\NormalTok{(binned\_woe), }\AttributeTok{levels =} \FunctionTok{names}\NormalTok{(binned\_woe)),}
        \AttributeTok{good =} \FunctionTok{sapply}\NormalTok{(binned\_woe, }\ControlFlowTok{function}\NormalTok{(x) x[}\StringTok{"good"}\NormalTok{]),}
        \AttributeTok{bad =} \FunctionTok{sapply}\NormalTok{(binned\_woe, }\ControlFlowTok{function}\NormalTok{(x) x[}\StringTok{"bad"}\NormalTok{])}
\NormalTok{    ) }
    \FunctionTok{rownames}\NormalTok{(binned\_woe\_df) }\OtherTok{\textless{}{-}} \ConstantTok{NULL}
\NormalTok{    binned\_woe\_df}\SpecialCharTok{$}\NormalTok{prop\_good }\OtherTok{\textless{}{-}}\NormalTok{ binned\_woe\_df}\SpecialCharTok{$}\NormalTok{good }\SpecialCharTok{/} \FunctionTok{sum}\NormalTok{(binned\_woe\_df}\SpecialCharTok{$}\NormalTok{good)}
\NormalTok{    binned\_woe\_df}\SpecialCharTok{$}\NormalTok{prop\_bad }\OtherTok{\textless{}{-}}\NormalTok{ binned\_woe\_df}\SpecialCharTok{$}\NormalTok{bad }\SpecialCharTok{/} \FunctionTok{sum}\NormalTok{(binned\_woe\_df}\SpecialCharTok{$}\NormalTok{bad)}
\NormalTok{    binned\_woe\_df}\SpecialCharTok{$}\NormalTok{woe }\OtherTok{\textless{}{-}} \FunctionTok{log}\NormalTok{(binned\_woe\_df}\SpecialCharTok{$}\NormalTok{prop\_good }\SpecialCharTok{/}\NormalTok{ binned\_woe\_df}\SpecialCharTok{$}\NormalTok{prop\_bad)}
\NormalTok{    binned\_woe\_df}\SpecialCharTok{$}\NormalTok{iv }\OtherTok{\textless{}{-}}\NormalTok{ iv }\OtherTok{\textless{}{-}}\NormalTok{ (binned\_woe\_df}\SpecialCharTok{$}\NormalTok{prop\_good }\SpecialCharTok{{-}}\NormalTok{ binned\_woe\_df}\SpecialCharTok{$}\NormalTok{prop\_bad) }\SpecialCharTok{*}\NormalTok{ binned\_woe\_df}\SpecialCharTok{$}\NormalTok{woe}
    \FunctionTok{return}\NormalTok{(binned\_woe\_df)}
\NormalTok{\}}

\NormalTok{plot\_woe\_iv }\OtherTok{\textless{}{-}} \ControlFlowTok{function}\NormalTok{(woe\_df, }\AttributeTok{interactive =} \ConstantTok{TRUE}\NormalTok{) \{}
\NormalTok{    max\_woe }\OtherTok{\textless{}{-}} \FunctionTok{max}\NormalTok{(}\FunctionTok{abs}\NormalTok{(woe\_df}\SpecialCharTok{$}\NormalTok{woe))}
\NormalTok{    max\_iv }\OtherTok{\textless{}{-}} \FunctionTok{max}\NormalTok{(}\FunctionTok{abs}\NormalTok{(woe\_df}\SpecialCharTok{$}\NormalTok{iv))}
    \CommentTok{\# Avoid dividing by zero when all values are zero}
\NormalTok{    scale\_factor }\OtherTok{\textless{}{-}} \ControlFlowTok{if}\NormalTok{ (max\_woe }\SpecialCharTok{\textgreater{}} \DecValTok{0} \SpecialCharTok{\&\&}\NormalTok{ max\_iv }\SpecialCharTok{\textgreater{}} \DecValTok{0}\NormalTok{) \{}
\NormalTok{        max\_woe }\SpecialCharTok{/}\NormalTok{ max\_iv}
\NormalTok{    \} }\ControlFlowTok{else}\NormalTok{ \{}
        \DecValTok{1}
\NormalTok{    \}}
\NormalTok{    p }\OtherTok{\textless{}{-}} \FunctionTok{ggplot}\NormalTok{(woe\_df, }\FunctionTok{aes}\NormalTok{(}\AttributeTok{x =}\NormalTok{ bin)) }\SpecialCharTok{+}
        \FunctionTok{geom\_col}\NormalTok{(}
            \FunctionTok{aes}\NormalTok{(}\AttributeTok{y =}\NormalTok{ woe, }\AttributeTok{fill =}\NormalTok{ woe }\SpecialCharTok{\textgreater{}} \DecValTok{0}\NormalTok{),}
            \AttributeTok{width =} \FloatTok{0.8}\NormalTok{,}
            \AttributeTok{show.legend =} \ConstantTok{FALSE}
\NormalTok{        ) }\SpecialCharTok{+}
        \FunctionTok{scale\_fill\_manual}\NormalTok{(}
            \AttributeTok{values =} \FunctionTok{c}\NormalTok{(}\StringTok{"FALSE"} \OtherTok{=} \StringTok{"firebrick"}\NormalTok{, }\StringTok{"TRUE"} \OtherTok{=} \StringTok{"forestgreen"}\NormalTok{)}
\NormalTok{        ) }\SpecialCharTok{+}
        \FunctionTok{geom\_line}\NormalTok{(}
            \FunctionTok{aes}\NormalTok{(}\AttributeTok{y =}\NormalTok{ iv }\SpecialCharTok{*}\NormalTok{ scale\_factor, }\AttributeTok{group =} \DecValTok{1}\NormalTok{),}
            \AttributeTok{color =} \StringTok{"blue"}
\NormalTok{        ) }\SpecialCharTok{+}
        \FunctionTok{geom\_point}\NormalTok{(}
            \FunctionTok{aes}\NormalTok{(}\AttributeTok{y =}\NormalTok{ iv }\SpecialCharTok{*}\NormalTok{ scale\_factor),}
            \AttributeTok{color =} \StringTok{"blue"}\NormalTok{,}
\NormalTok{        ) }\SpecialCharTok{+}
        \FunctionTok{scale\_y\_continuous}\NormalTok{(}
            \AttributeTok{name =} \StringTok{"WoE and Scale{-}adjusted IV"}
\NormalTok{        ) }\SpecialCharTok{+}
        \FunctionTok{labs}\NormalTok{(}\AttributeTok{x =} \StringTok{"Bins"}\NormalTok{) }\SpecialCharTok{+}
        \FunctionTok{theme\_minimal}\NormalTok{() }\SpecialCharTok{+}
        \FunctionTok{theme}\NormalTok{(}
            \AttributeTok{axis.text.x =} \FunctionTok{element\_text}\NormalTok{(}\AttributeTok{angle =} \DecValTok{60}\NormalTok{, }\AttributeTok{hjust =} \DecValTok{1}\NormalTok{),}
\NormalTok{        ) }\SpecialCharTok{+}
        \FunctionTok{ggtitle}\NormalTok{(}
            \FunctionTok{paste0}\NormalTok{(}\StringTok{"WoE and IV (bins="}\NormalTok{, }\FunctionTok{length}\NormalTok{(woe\_df}\SpecialCharTok{$}\NormalTok{bin), }\StringTok{")"}\NormalTok{)}
\NormalTok{        )}
    \ControlFlowTok{if}\NormalTok{ (interactive) \{}
        \FunctionTok{ggplotly}\NormalTok{(p)}
\NormalTok{    \} }\ControlFlowTok{else}\NormalTok{ \{}
        \FunctionTok{return}\NormalTok{(p)}
\NormalTok{    \}}
\NormalTok{\}}
\end{Highlighting}
\end{Shaded}

\begin{Shaded}
\begin{Highlighting}[]
\NormalTok{woe\_ex }\OtherTok{\textless{}{-}} \FunctionTok{woe\_calc}\NormalTok{(data, }\AttributeTok{x =} \StringTok{"AGEPH"}\NormalTok{, }\AttributeTok{y =} \StringTok{"NCLAIMS"}\NormalTok{, }\AttributeTok{bins =}\NormalTok{ bins\_ageph)}
\FunctionTok{plot\_woe\_iv}\NormalTok{(woe\_ex)}
\end{Highlighting}
\end{Shaded}

\begin{center}\includegraphics{project_MD_files/figure-latex/example_woe_calculation-1} \end{center}

The example for \texttt{AGEPH} already shows very interesting patterns
for our baseline binning:

\begin{itemize}
\tightlist
\item
  The WoE values go from negative to positive, indicating that the
  younger age groups have a higher likelihood of bad outcomes (at least
  one claim) compared to the older age groups.
\item
  The IV line shows a very high contribution from the young and older
  age groups, indicating that these bins are the most predictive for the
  likelihood of claims.
\end{itemize}

Let's confirm visuallythe relationship between age and number of claims
by examining the boxplots and claim frequency by age bin.

\begin{Shaded}
\begin{Highlighting}[]
\CommentTok{\# boxplot of age for each number of claims}
\FunctionTok{boxplot}\NormalTok{(AGEPH }\SpecialCharTok{\textasciitilde{}}\NormalTok{ NCLAIMS, }\AttributeTok{data =}\NormalTok{ data, }\AttributeTok{xlab =} \StringTok{"NCLAIMS"}\NormalTok{, }\AttributeTok{ylab =} \StringTok{"AGEPH"}\NormalTok{, }\AttributeTok{main =} \StringTok{"Boxplot of AGEPH by NCLAIMS"}\NormalTok{)}
\end{Highlighting}
\end{Shaded}

\begin{center}\includegraphics{project_MD_files/figure-latex/age_claims_correlation-1} \end{center}

\begin{Shaded}
\begin{Highlighting}[]
\CommentTok{\# claim frequency in our example bins, group by bin}
\NormalTok{age\_quantiles }\OtherTok{\textless{}{-}} \FunctionTok{cut}\NormalTok{(data}\SpecialCharTok{$}\NormalTok{AGEPH, }\AttributeTok{breaks =}\NormalTok{ bins\_ageph, }\AttributeTok{include.lowest =} \ConstantTok{TRUE}\NormalTok{)}
\NormalTok{gpd\_data }\OtherTok{\textless{}{-}}\NormalTok{ data }\SpecialCharTok{\%\textgreater{}\%} 
    \FunctionTok{mutate}\NormalTok{(}\AttributeTok{age\_bin =}\NormalTok{ age\_quantiles) }\SpecialCharTok{\%\textgreater{}\%} 
    \FunctionTok{group\_by}\NormalTok{(age\_bin) }\SpecialCharTok{\%\textgreater{}\%} 
    \FunctionTok{summarise}\NormalTok{(}
        \AttributeTok{sum\_exp =} \FunctionTok{sum}\NormalTok{(EXP),}
        \AttributeTok{sum\_claims =} \FunctionTok{sum}\NormalTok{(NCLAIMS),}
        \AttributeTok{lambda =}\NormalTok{ sum\_claims }\SpecialCharTok{/}\NormalTok{ sum\_exp,}
        \AttributeTok{sd\_lambda =} \FunctionTok{sqrt}\NormalTok{(lambda }\SpecialCharTok{/}\NormalTok{ sum\_exp),}
        \AttributeTok{lambda\_95\_lower =}\NormalTok{ lambda }\SpecialCharTok{{-}} \FloatTok{1.96} \SpecialCharTok{*}\NormalTok{ sd\_lambda,}
        \AttributeTok{lambda\_95\_upper =}\NormalTok{ lambda }\SpecialCharTok{+} \FloatTok{1.96} \SpecialCharTok{*}\NormalTok{ sd\_lambda,}
        \AttributeTok{lambda\_99\_lower =}\NormalTok{ lambda }\SpecialCharTok{{-}} \FloatTok{2.58} \SpecialCharTok{*}\NormalTok{ sd\_lambda,}
        \AttributeTok{lambda\_99\_upper =}\NormalTok{ lambda }\SpecialCharTok{+} \FloatTok{2.58} \SpecialCharTok{*}\NormalTok{ sd\_lambda}
\NormalTok{    )}
\CommentTok{\# Plot the claim frequency (lambda) with filled error bars 95 and 99 percent confidence intervals}
\NormalTok{p2 }\OtherTok{\textless{}{-}} \FunctionTok{ggplot}\NormalTok{(gpd\_data, }\FunctionTok{aes}\NormalTok{(}\AttributeTok{x =}\NormalTok{ age\_bin, }\AttributeTok{group =} \DecValTok{1}\NormalTok{)) }\SpecialCharTok{+}
    \FunctionTok{geom\_ribbon}\NormalTok{(}
        \FunctionTok{aes}\NormalTok{(}\AttributeTok{ymin =}\NormalTok{ lambda\_99\_lower, }\AttributeTok{ymax =}\NormalTok{ lambda\_99\_upper),}
        \AttributeTok{fill =} \StringTok{"gray"}\NormalTok{, }\AttributeTok{alpha =} \FloatTok{0.3}\NormalTok{, }\AttributeTok{label =} \StringTok{"99\% CI"}
\NormalTok{    ) }\SpecialCharTok{+}
    \FunctionTok{geom\_ribbon}\NormalTok{(}
        \FunctionTok{aes}\NormalTok{(}\AttributeTok{ymin =}\NormalTok{ lambda\_95\_lower, }\AttributeTok{ymax =}\NormalTok{ lambda\_95\_upper),}
        \AttributeTok{fill =} \StringTok{"lightblue"}\NormalTok{, }\AttributeTok{alpha =} \FloatTok{0.5}\NormalTok{, }\AttributeTok{label =} \StringTok{"95\% CI"}
\NormalTok{    ) }\SpecialCharTok{+}
    \FunctionTok{geom\_line}\NormalTok{(}\FunctionTok{aes}\NormalTok{(}\AttributeTok{y =}\NormalTok{ lambda), }\AttributeTok{color =} \StringTok{"black"}\NormalTok{) }\SpecialCharTok{+}
    \FunctionTok{geom\_point}\NormalTok{(}\FunctionTok{aes}\NormalTok{(}\AttributeTok{y =}\NormalTok{ lambda), }\AttributeTok{color =} \StringTok{"black"}\NormalTok{) }\SpecialCharTok{+}
    \FunctionTok{theme\_minimal}\NormalTok{() }\SpecialCharTok{+}
    \FunctionTok{scale\_y\_continuous}\NormalTok{(}\AttributeTok{labels =}\NormalTok{ scales}\SpecialCharTok{::}\NormalTok{comma) }\SpecialCharTok{+}
    \FunctionTok{labs}\NormalTok{(}
        \AttributeTok{x =} \StringTok{"Age Bin"}\NormalTok{,}
        \AttributeTok{y =} \StringTok{"Claim Frequency (Lambda)"}\NormalTok{,}
        \AttributeTok{title =} \StringTok{"Claim Frequency by Age Bin"}
\NormalTok{    ) }\SpecialCharTok{+}
    \FunctionTok{theme}\NormalTok{(}\AttributeTok{axis.text.x =} \FunctionTok{element\_text}\NormalTok{(}\AttributeTok{angle =} \DecValTok{45}\NormalTok{, }\AttributeTok{hjust =} \DecValTok{1}\NormalTok{))}
\FunctionTok{ggplotly}\NormalTok{(p2)}
\end{Highlighting}
\end{Shaded}

\begin{center}\includegraphics{project_MD_files/figure-latex/age_claims_correlation-2} \end{center}

The boxplots show that the group of policyholders with lower amount of
claims do tend to be older. It is even more striking when we look at the
claim frequency by age bin with the confidence intervals\footnote{The
  confidence intervals are based on the MLE estimator for the claim
  frequency under the assumption of a Poisson distribution for the
  number of claims.}.

\paragraph{Iterative Merging Based on Information Value
(IV)}\label{iterative-merging-based-on-information-value-iv}

Remember that the Information Value (IV) is a measure based on the
Weight of Evidence (WoE) that quantifies the predictive power of a
feature in distinguishing between good and bad outcomes. It is
calculated as:

\[\text{IV} = \sum_{\text{bins}} (\text{proportion of Good Outcomes in the bin} - \text{proportion of Bad Outcomes in the bin}) \times \text{WoE of the bin}\]

We have an information value for each bin, which allows us to identify
the adjacent bins that contribute the least to the overall predictive
power and merge them iteratively. The goal is to recursively merge the
bin with lowest IV contribution with its neighboring bin, and evaluate
the resulting IV. Then we can find the elbow point where further merging
would lead to a significant loss of predictive power.

\begin{Shaded}
\begin{Highlighting}[]
\NormalTok{iv\_elbow }\OtherTok{\textless{}{-}} \ControlFlowTok{function}\NormalTok{(data, x, y, bins) \{}
    \StringTok{"Recursively merge the bin with the lowest IV contribution with its neighboring bin with minimum IV"}

\NormalTok{    temp\_bins }\OtherTok{\textless{}{-}}\NormalTok{ bins}
\NormalTok{    temp\_n\_bins }\OtherTok{\textless{}{-}} \FunctionTok{length}\NormalTok{(bins) }\SpecialCharTok{{-}} \DecValTok{1}
\NormalTok{    temp\_woe\_df }\OtherTok{\textless{}{-}} \FunctionTok{woe\_calc}\NormalTok{(data, x, y, temp\_bins)}
\NormalTok{    temp\_iv }\OtherTok{\textless{}{-}}\NormalTok{ temp\_woe\_df}\SpecialCharTok{$}\NormalTok{iv}
\NormalTok{    temp\_total\_iv }\OtherTok{\textless{}{-}} \FunctionTok{sum}\NormalTok{(temp\_iv)}

\NormalTok{    iter\_log }\OtherTok{\textless{}{-}} \FunctionTok{data.frame}\NormalTok{(}
        \AttributeTok{n\_bins =}\NormalTok{ temp\_n\_bins,}
        \AttributeTok{iv =}\NormalTok{ temp\_total\_iv,}
        \AttributeTok{bins =} \FunctionTok{paste}\NormalTok{(temp\_bins, }\AttributeTok{collapse =} \StringTok{","}\NormalTok{),}
        \AttributeTok{diff =} \ConstantTok{NA}\NormalTok{,}
        \AttributeTok{diff\_order =} \ConstantTok{NA}
\NormalTok{    )}

    \ControlFlowTok{while}\NormalTok{(temp\_n\_bins }\SpecialCharTok{\textgreater{}} \DecValTok{2}\NormalTok{) \{}
\NormalTok{        bin\_to\_merge }\OtherTok{\textless{}{-}} \FunctionTok{which.min}\NormalTok{(temp\_iv)}
        \ControlFlowTok{if}\NormalTok{ (bin\_to\_merge }\SpecialCharTok{==} \DecValTok{1}\NormalTok{) \{}
\NormalTok{            neighbor }\OtherTok{\textless{}{-}} \DecValTok{2}
\NormalTok{        \} }\ControlFlowTok{else} \ControlFlowTok{if}\NormalTok{ (bin\_to\_merge }\SpecialCharTok{==}\NormalTok{ temp\_n\_bins) \{}
\NormalTok{            neighbor }\OtherTok{\textless{}{-}}\NormalTok{ temp\_n\_bins }\SpecialCharTok{{-}} \DecValTok{1}
\NormalTok{        \} }\ControlFlowTok{else}\NormalTok{ \{}
\NormalTok{            neighbor }\OtherTok{\textless{}{-}} \FunctionTok{ifelse}\NormalTok{(temp\_iv[bin\_to\_merge }\SpecialCharTok{{-}} \DecValTok{1}\NormalTok{] }\SpecialCharTok{\textless{}}\NormalTok{ temp\_iv[bin\_to\_merge }\SpecialCharTok{+} \DecValTok{1}\NormalTok{], bin\_to\_merge }\SpecialCharTok{{-}} \DecValTok{1}\NormalTok{, bin\_to\_merge }\SpecialCharTok{+} \DecValTok{1}\NormalTok{)}
\NormalTok{        \}}
\NormalTok{        to\_merge }\OtherTok{\textless{}{-}} \FunctionTok{sort}\NormalTok{(}\FunctionTok{c}\NormalTok{(bin\_to\_merge, neighbor))}
\NormalTok{        temp\_bins }\OtherTok{\textless{}{-}}\NormalTok{ temp\_bins[}\SpecialCharTok{{-}}\NormalTok{to\_merge[}\DecValTok{2}\NormalTok{]]}
\NormalTok{        temp\_n\_bins }\OtherTok{\textless{}{-}} \FunctionTok{length}\NormalTok{(temp\_bins) }\SpecialCharTok{{-}} \DecValTok{1}
\NormalTok{        temp\_woe\_df }\OtherTok{\textless{}{-}} \FunctionTok{woe\_calc}\NormalTok{(data, x, y, temp\_bins)}
\NormalTok{        temp\_iv }\OtherTok{\textless{}{-}}\NormalTok{ temp\_woe\_df}\SpecialCharTok{$}\NormalTok{iv}
\NormalTok{        temp\_total\_iv }\OtherTok{\textless{}{-}} \FunctionTok{sum}\NormalTok{(temp\_iv)}

\NormalTok{        iter\_log }\OtherTok{\textless{}{-}} \FunctionTok{rbind}\NormalTok{(iter\_log, }\FunctionTok{data.frame}\NormalTok{(}
            \AttributeTok{n\_bins =}\NormalTok{ temp\_n\_bins,}
            \AttributeTok{iv =}\NormalTok{ temp\_total\_iv,}
            \AttributeTok{bins =} \FunctionTok{paste}\NormalTok{(temp\_bins, }\AttributeTok{collapse =} \StringTok{","}\NormalTok{),}
            \AttributeTok{diff =} \ConstantTok{NA}\NormalTok{,}
            \AttributeTok{diff\_order =} \ConstantTok{NA}
\NormalTok{        ))}
\NormalTok{    \}}
\NormalTok{    iter\_log }\OtherTok{\textless{}{-}}\NormalTok{ iter\_log }\SpecialCharTok{\%\textgreater{}\%} \FunctionTok{arrange}\NormalTok{(n\_bins)}
\NormalTok{    iter\_log}\SpecialCharTok{$}\NormalTok{diff }\OtherTok{\textless{}{-}} \FunctionTok{c}\NormalTok{(}\ConstantTok{NA}\NormalTok{, }\FunctionTok{abs}\NormalTok{(}\FunctionTok{diff}\NormalTok{(iter\_log}\SpecialCharTok{$}\NormalTok{iv)))}
\NormalTok{    iter\_log}\SpecialCharTok{$}\NormalTok{diff[}\FunctionTok{is.na}\NormalTok{(iter\_log}\SpecialCharTok{$}\NormalTok{diff)] }\OtherTok{\textless{}{-}} \DecValTok{0}
\NormalTok{    iter\_log}\SpecialCharTok{$}\NormalTok{diff\_order }\OtherTok{\textless{}{-}} \FunctionTok{rank}\NormalTok{(}\SpecialCharTok{{-}}\NormalTok{iter\_log}\SpecialCharTok{$}\NormalTok{diff)}
    \FunctionTok{return}\NormalTok{(iter\_log)}
\NormalTok{\}}
\NormalTok{plot\_iv\_elbow }\OtherTok{\textless{}{-}} \ControlFlowTok{function}\NormalTok{(iv\_elbow\_result, }\AttributeTok{interactive =} \ConstantTok{TRUE}\NormalTok{) \{}
\NormalTok{    first\_max\_diff }\OtherTok{\textless{}{-}} \FunctionTok{filter}\NormalTok{(iv\_elbow\_result, diff\_order }\SpecialCharTok{==} \DecValTok{1}\NormalTok{)}
\NormalTok{    p }\OtherTok{\textless{}{-}} \FunctionTok{ggplot}\NormalTok{(iv\_elbow\_result, }\FunctionTok{aes}\NormalTok{(}\AttributeTok{x =}\NormalTok{ n\_bins, }\AttributeTok{y =}\NormalTok{ iv)) }\SpecialCharTok{+}
        \FunctionTok{geom\_line}\NormalTok{() }\SpecialCharTok{+}
        \FunctionTok{geom\_point}\NormalTok{() }\SpecialCharTok{+}
        \FunctionTok{labs}\NormalTok{(}\AttributeTok{title =} \StringTok{"Information Value Elbow Plot"}\NormalTok{, }\AttributeTok{x =} \StringTok{"Number of Bins"}\NormalTok{, }\AttributeTok{y =} \StringTok{"Information Value"}\NormalTok{) }\SpecialCharTok{+}
        \FunctionTok{theme\_minimal}\NormalTok{() }\SpecialCharTok{+}
        \FunctionTok{geom\_vline}\NormalTok{(}\AttributeTok{xintercept =}\NormalTok{ first\_max\_diff}\SpecialCharTok{$}\NormalTok{n\_bins, }\AttributeTok{linetype =} \StringTok{"dashed"}\NormalTok{, }\AttributeTok{color =} \StringTok{"red"}\NormalTok{) }\SpecialCharTok{+}
        \FunctionTok{geom\_hline}\NormalTok{(}\AttributeTok{yintercept =}\NormalTok{ first\_max\_diff}\SpecialCharTok{$}\NormalTok{iv, }\AttributeTok{linetype =} \StringTok{"dashed"}\NormalTok{, }\AttributeTok{color =} \StringTok{"red"}\NormalTok{)}
    \ControlFlowTok{if}\NormalTok{ (interactive) \{}
        \FunctionTok{ggplotly}\NormalTok{(p)}
\NormalTok{    \} }\ControlFlowTok{else}\NormalTok{ \{}
        \FunctionTok{return}\NormalTok{(p)}
\NormalTok{    \}}
\NormalTok{\}}
\end{Highlighting}
\end{Shaded}

\begin{Shaded}
\begin{Highlighting}[]
\NormalTok{result\_named\_list }\OtherTok{\textless{}{-}} \FunctionTok{list}\NormalTok{()}
\NormalTok{n\_bins\_named\_list }\OtherTok{\textless{}{-}} \FunctionTok{list}\NormalTok{()}
\NormalTok{result\_ageph }\OtherTok{\textless{}{-}} \FunctionTok{iv\_elbow}\NormalTok{(data, }\StringTok{"AGEPH"}\NormalTok{, }\StringTok{"NCLAIMS"}\NormalTok{, bins\_ageph)}
\NormalTok{result\_named\_list[[}\StringTok{"AGEPH"}\NormalTok{]] }\OtherTok{\textless{}{-}}\NormalTok{ result\_ageph}
\FunctionTok{plot\_iv\_elbow}\NormalTok{(result\_ageph, }\AttributeTok{interactive =} \ConstantTok{TRUE}\NormalTok{)}
\end{Highlighting}
\end{Shaded}

\begin{center}\includegraphics{project_MD_files/figure-latex/iv_elbow_result-1} \end{center}

Looking at the IV elbow plot and taking the number of bins corresponding
to the first maximum difference in IV as the optimal choice, we can
rebuild the bins with only three bins. Adding a single bin seems
sensible as we capture more information with a single additional split.
We could also consider 7 or 12 bins that mark other levels in the IV -
Number of Bins relationship.

\begin{Shaded}
\begin{Highlighting}[]
\NormalTok{rebuild\_bins }\OtherTok{\textless{}{-}} \ControlFlowTok{function}\NormalTok{(iv\_elbow\_result, n) \{}
    \StringTok{"Rebuild bins based on the IV elbow result for the specified number of bins"}
\NormalTok{    target\_row }\OtherTok{\textless{}{-}}\NormalTok{ iv\_elbow\_result }\SpecialCharTok{\%\textgreater{}\%} \FunctionTok{filter}\NormalTok{(n\_bins }\SpecialCharTok{==}\NormalTok{ n)}
\NormalTok{    bins }\OtherTok{\textless{}{-}} \FunctionTok{as.numeric}\NormalTok{(}\FunctionTok{unlist}\NormalTok{(}\FunctionTok{strsplit}\NormalTok{(target\_row}\SpecialCharTok{$}\NormalTok{bins, }\StringTok{","}\NormalTok{)))}
    \FunctionTok{return}\NormalTok{(bins)}
\NormalTok{\}}
\end{Highlighting}
\end{Shaded}

\begin{Shaded}
\begin{Highlighting}[]
\NormalTok{plots }\OtherTok{\textless{}{-}} \FunctionTok{list}\NormalTok{()}
\ControlFlowTok{for}\NormalTok{ (i }\ControlFlowTok{in} \FunctionTok{c}\NormalTok{(}\DecValTok{3}\NormalTok{, }\DecValTok{4}\NormalTok{, }\DecValTok{7}\NormalTok{, }\DecValTok{12}\NormalTok{)) \{}
\NormalTok{    temp\_bins }\OtherTok{\textless{}{-}} \FunctionTok{rebuild\_bins}\NormalTok{(result\_ageph, i)}
\NormalTok{    woe\_temp }\OtherTok{\textless{}{-}} \FunctionTok{woe\_calc}\NormalTok{(data, }\AttributeTok{x =} \StringTok{"AGEPH"}\NormalTok{, }\AttributeTok{y =} \StringTok{"NCLAIMS"}\NormalTok{, }\AttributeTok{bins =}\NormalTok{ temp\_bins)}
\NormalTok{    plots[[i]] }\OtherTok{\textless{}{-}} \FunctionTok{plot\_woe\_iv}\NormalTok{(woe\_temp, }\AttributeTok{interactive =} \ConstantTok{FALSE}\NormalTok{)}
\NormalTok{\}}
\NormalTok{(plots[[}\DecValTok{3}\NormalTok{]] }\SpecialCharTok{|}\NormalTok{ plots[[}\DecValTok{4}\NormalTok{]] ) }\SpecialCharTok{/}\NormalTok{ (plots[[}\DecValTok{7}\NormalTok{]] }\SpecialCharTok{|}\NormalTok{ plots[[}\DecValTok{12}\NormalTok{]])}
\end{Highlighting}
\end{Shaded}

\begin{center}\includegraphics{project_MD_files/figure-latex/rebuild_bins_result-1} \end{center}

Based on the comparison above, the optimal amount of bins seems to be 7
because:

\begin{itemize}
\tightlist
\item
  We capture in fine granularity the effet of young ages with higher
  risk.
\item
  We skip the middle of the age distribution which brings no additional
  information (28 to 66).
\item
  We capture the effect of older ages with lower risk profiles with good
  granularity.
\end{itemize}

The complex pattern we noticed in our initial 30 bins seems to be well
captured with 7 bins only.

\begin{Shaded}
\begin{Highlighting}[]
\CommentTok{\# We add the ordered factor variable for the binned AGEPH}
\NormalTok{data}\SpecialCharTok{$}\NormalTok{AGEPH\_binned }\OtherTok{\textless{}{-}} \FunctionTok{cut}\NormalTok{(}
\NormalTok{    data}\SpecialCharTok{$}\NormalTok{AGEPH, }
    \AttributeTok{breaks =} \FunctionTok{rebuild\_bins}\NormalTok{(result\_ageph, }\DecValTok{7}\NormalTok{), }\AttributeTok{include.lowest =} \ConstantTok{TRUE}\NormalTok{, }\AttributeTok{ordered\_result =} \ConstantTok{TRUE}\NormalTok{)}
\CommentTok{\# Check the distribution of the binned variable}
\NormalTok{n\_bins\_named\_list[[}\StringTok{"AGEPH"}\NormalTok{]] }\OtherTok{\textless{}{-}} \DecValTok{7}
\FunctionTok{summary}\NormalTok{(data}\SpecialCharTok{$}\NormalTok{AGEPH\_binned) }\SpecialCharTok{\%\textgreater{}\%}
    \FunctionTok{kable}\NormalTok{(}
        \AttributeTok{caption =} \StringTok{"Distribution of binned AGEPH"}\NormalTok{,}
        \AttributeTok{booktabs =} \ConstantTok{TRUE}
\NormalTok{    )}
\end{Highlighting}
\end{Shaded}

\begin{longtable}[]{@{}lr@{}}
\caption{Distribution of binned AGEPH}\tabularnewline
\toprule\noalign{}
& x \\
\midrule\noalign{}
\endfirsthead
\toprule\noalign{}
& x \\
\midrule\noalign{}
\endhead
\bottomrule\noalign{}
\endlastfoot
{[}18,24{]} & 4466 \\
(24,26{]} & 4453 \\
(26,28{]} & 5052 \\
(28,66{]} & 100180 \\
(66,68{]} & 3478 \\
(68,71{]} & 4362 \\
(71,95{]} & 8594 \\
\end{longtable}

\paragraph{Bonus-Malus Binning (BM)}\label{bonus-malus-binning-bm}

\begin{Shaded}
\begin{Highlighting}[]
\NormalTok{baseline\_bm }\OtherTok{\textless{}{-}} \FunctionTok{qbinning}\NormalTok{(data, }\AttributeTok{x =} \StringTok{"BM"}\NormalTok{, }\AttributeTok{y =} \StringTok{"NCLAIMS"}\NormalTok{, }\AttributeTok{max\_bins =} \DecValTok{30}\NormalTok{)}
\NormalTok{result\_bm }\OtherTok{\textless{}{-}} \FunctionTok{iv\_elbow}\NormalTok{(data, }\StringTok{"BM"}\NormalTok{, }\StringTok{"NCLAIMS"}\NormalTok{, baseline\_bm)}
\NormalTok{result\_named\_list[[}\StringTok{"BM"}\NormalTok{]] }\OtherTok{\textless{}{-}}\NormalTok{ result\_bm}
\FunctionTok{plot\_iv\_elbow}\NormalTok{(result\_bm, }\AttributeTok{interactive =} \ConstantTok{TRUE}\NormalTok{)}
\end{Highlighting}
\end{Shaded}

\begin{center}\includegraphics{project_MD_files/figure-latex/bm_binning-1} \end{center}

Based on the Elbow plot for the Bonus-Malus (\texttt{BM}), let's try 3,
4, 5, and 11 bins.

\begin{Shaded}
\begin{Highlighting}[]
\NormalTok{plots\_bm }\OtherTok{\textless{}{-}} \FunctionTok{list}\NormalTok{()}
\ControlFlowTok{for}\NormalTok{ (i }\ControlFlowTok{in} \FunctionTok{c}\NormalTok{(}\DecValTok{3}\NormalTok{, }\DecValTok{4}\NormalTok{, }\DecValTok{5}\NormalTok{, }\DecValTok{11}\NormalTok{)) \{}
\NormalTok{    temp\_bins\_bm }\OtherTok{\textless{}{-}} \FunctionTok{rebuild\_bins}\NormalTok{(result\_bm, i)}
\NormalTok{    woe\_temp\_bm }\OtherTok{\textless{}{-}} \FunctionTok{woe\_calc}\NormalTok{(data, }\AttributeTok{x =} \StringTok{"BM"}\NormalTok{, }\AttributeTok{y =} \StringTok{"NCLAIMS"}\NormalTok{, }\AttributeTok{bins =}\NormalTok{ temp\_bins\_bm)}
\NormalTok{    plots\_bm[[i]] }\OtherTok{\textless{}{-}} \FunctionTok{plot\_woe\_iv}\NormalTok{(woe\_temp\_bm, }\AttributeTok{interactive =} \ConstantTok{FALSE}\NormalTok{)}
\NormalTok{\}}
\NormalTok{(plots\_bm[[}\DecValTok{3}\NormalTok{]] }\SpecialCharTok{|}\NormalTok{ plots\_bm[[}\DecValTok{4}\NormalTok{]] ) }\SpecialCharTok{/}\NormalTok{ (plots\_bm[[}\DecValTok{5}\NormalTok{]] }\SpecialCharTok{|}\NormalTok{ plots\_bm[[}\DecValTok{11}\NormalTok{]])}
\end{Highlighting}
\end{Shaded}

\begin{center}\includegraphics{project_MD_files/figure-latex/rebuild_bins_bm-1} \end{center}

We see that most information is captured by the first bin indicating
that those with a BM of 0 are significantly safer than the others. The
three bins are sufficient to capture the first group at bonus-malus
zero, then the first and second risk groups that are the first and
second half of the BM score. Adding bins just creates bins in between
with low informational value.

\begin{Shaded}
\begin{Highlighting}[]
\NormalTok{data}\SpecialCharTok{$}\NormalTok{BM\_binned }\OtherTok{\textless{}{-}} \FunctionTok{cut}\NormalTok{(}
\NormalTok{    data}\SpecialCharTok{$}\NormalTok{BM,}
    \AttributeTok{breaks =} \FunctionTok{rebuild\_bins}\NormalTok{(result\_bm, }\DecValTok{3}\NormalTok{), }\AttributeTok{include.lowest =} \ConstantTok{TRUE}\NormalTok{, }\AttributeTok{ordered\_result =} \ConstantTok{TRUE}\NormalTok{)}
\NormalTok{n\_bins\_named\_list[[}\StringTok{"BM"}\NormalTok{]] }\OtherTok{\textless{}{-}} \DecValTok{3}
\end{Highlighting}
\end{Shaded}

\paragraph{Age of the Car Binning
(AGEC)}\label{age-of-the-car-binning-agec}

\begin{Shaded}
\begin{Highlighting}[]
\NormalTok{baseline\_agec }\OtherTok{\textless{}{-}} \FunctionTok{qbinning}\NormalTok{(data, }\AttributeTok{x =} \StringTok{"AGEC"}\NormalTok{, }\AttributeTok{y =} \StringTok{"NCLAIMS"}\NormalTok{, }\AttributeTok{max\_bins =} \DecValTok{30}\NormalTok{)}
\NormalTok{result\_agec }\OtherTok{\textless{}{-}} \FunctionTok{iv\_elbow}\NormalTok{(data, }\StringTok{"AGEC"}\NormalTok{, }\StringTok{"NCLAIMS"}\NormalTok{, baseline\_agec)}
\NormalTok{result\_named\_list[[}\StringTok{"AGEC"}\NormalTok{]] }\OtherTok{\textless{}{-}}\NormalTok{ result\_agec}
\FunctionTok{plot\_iv\_elbow}\NormalTok{(result\_agec, }\AttributeTok{interactive =} \ConstantTok{TRUE}\NormalTok{)}
\end{Highlighting}
\end{Shaded}

\begin{center}\includegraphics{project_MD_files/figure-latex/agec_binning-1} \end{center}

For the age of the car (\texttt{AGEC}), we also see a plateau after
around 5 and 9 bins, let's compare 3, 5, 9 and 16 (max) bins.

\begin{Shaded}
\begin{Highlighting}[]
\NormalTok{plots\_agec }\OtherTok{\textless{}{-}} \FunctionTok{list}\NormalTok{()}
\ControlFlowTok{for}\NormalTok{ (i }\ControlFlowTok{in} \FunctionTok{c}\NormalTok{(}\DecValTok{3}\NormalTok{, }\DecValTok{5}\NormalTok{, }\DecValTok{9}\NormalTok{, }\DecValTok{16}\NormalTok{)) \{}
\NormalTok{    temp\_bins\_agec }\OtherTok{\textless{}{-}} \FunctionTok{rebuild\_bins}\NormalTok{(result\_agec, i)}
\NormalTok{    woe\_temp\_agec }\OtherTok{\textless{}{-}} \FunctionTok{woe\_calc}\NormalTok{(data, }\AttributeTok{x =} \StringTok{"AGEC"}\NormalTok{, }\AttributeTok{y =} \StringTok{"NCLAIMS"}\NormalTok{, }\AttributeTok{bins =}\NormalTok{ temp\_bins\_agec)}
\NormalTok{    plots\_agec[[i]] }\OtherTok{\textless{}{-}} \FunctionTok{plot\_woe\_iv}\NormalTok{(woe\_temp\_agec, }\AttributeTok{interactive =} \ConstantTok{FALSE}\NormalTok{)}
\NormalTok{\}}
\NormalTok{(plots\_agec[[}\DecValTok{3}\NormalTok{]] }\SpecialCharTok{|}\NormalTok{ plots\_agec[[}\DecValTok{5}\NormalTok{]]) }\SpecialCharTok{/}\NormalTok{ (plots\_agec[[}\DecValTok{9}\NormalTok{]] }\SpecialCharTok{|}\NormalTok{ plots\_agec[[}\DecValTok{16}\NormalTok{]])}
\end{Highlighting}
\end{Shaded}

\begin{center}\includegraphics{project_MD_files/figure-latex/rebuild_bins_agec-1} \end{center}

The maximum binning shows us a clear separation of risk between:

\begin{itemize}
\tightlist
\item
  Brand new cars (AGEC close to 0) have the highest risk.
\item
  Older cars (AGEC over 15 years) have a lower risk.
\end{itemize}

This is captured from 3 bins onward. Using 5 bins would allow us to
capture the low-risk group for cars of 1 to 3 years old separately from
the older low-risk cars. These cars are typically good new cars with
conductors that have had the time to get used to them. We will
thereforeuse 5 bins for the age of the car.

\begin{Shaded}
\begin{Highlighting}[]
\NormalTok{data}\SpecialCharTok{$}\NormalTok{AGEC\_binned }\OtherTok{\textless{}{-}} \FunctionTok{cut}\NormalTok{(}
\NormalTok{    data}\SpecialCharTok{$}\NormalTok{AGEC,}
    \AttributeTok{breaks =} \FunctionTok{rebuild\_bins}\NormalTok{(result\_agec, }\DecValTok{5}\NormalTok{)}
\NormalTok{    , }\AttributeTok{include.lowest =} \ConstantTok{TRUE}\NormalTok{, }\AttributeTok{ordered\_result =} \ConstantTok{TRUE}\NormalTok{)}
\NormalTok{n\_bins\_named\_list[[}\StringTok{"AGEC"}\NormalTok{]] }\OtherTok{\textless{}{-}} \DecValTok{5}
\end{Highlighting}
\end{Shaded}

\paragraph{Power Binning (POWER)}\label{power-binning-power}

\begin{Shaded}
\begin{Highlighting}[]
\NormalTok{baseline\_power }\OtherTok{\textless{}{-}} \FunctionTok{qbinning}\NormalTok{(data, }\AttributeTok{x =} \StringTok{"POWER"}\NormalTok{, }\AttributeTok{y =} \StringTok{"NCLAIMS"}\NormalTok{, }\AttributeTok{max\_bins =} \DecValTok{30}\NormalTok{)}
\NormalTok{result\_power }\OtherTok{\textless{}{-}} \FunctionTok{iv\_elbow}\NormalTok{(data, }\StringTok{"POWER"}\NormalTok{, }\StringTok{"NCLAIMS"}\NormalTok{, baseline\_power)}
\NormalTok{result\_named\_list[[}\StringTok{"POWER"}\NormalTok{]] }\OtherTok{\textless{}{-}}\NormalTok{ result\_power}
\FunctionTok{plot\_iv\_elbow}\NormalTok{(result\_power, }\AttributeTok{interactive =} \ConstantTok{TRUE}\NormalTok{)}
\end{Highlighting}
\end{Shaded}

\begin{center}\includegraphics{project_MD_files/figure-latex/power_-1} \end{center}

For the power of the car (\texttt{POWER}), we have stairs at 6, 8, and
13 bins. We'll compare these three steps with the max level of bins
(23).

\begin{Shaded}
\begin{Highlighting}[]
\NormalTok{plots\_power }\OtherTok{\textless{}{-}} \FunctionTok{list}\NormalTok{()}
\ControlFlowTok{for}\NormalTok{ (i }\ControlFlowTok{in} \FunctionTok{c}\NormalTok{(}\DecValTok{6}\NormalTok{, }\DecValTok{8}\NormalTok{, }\DecValTok{13}\NormalTok{, }\DecValTok{23}\NormalTok{)) \{}
\NormalTok{    temp\_bins\_power }\OtherTok{\textless{}{-}} \FunctionTok{rebuild\_bins}\NormalTok{(result\_power, i)}
\NormalTok{    woe\_temp\_power }\OtherTok{\textless{}{-}} \FunctionTok{woe\_calc}\NormalTok{(data, }\AttributeTok{x =} \StringTok{"POWER"}\NormalTok{, }\AttributeTok{y =} \StringTok{"NCLAIMS"}\NormalTok{, }\AttributeTok{bins =}\NormalTok{ temp\_bins\_power)}
\NormalTok{    plots\_power[[i]] }\OtherTok{\textless{}{-}} \FunctionTok{plot\_woe\_iv}\NormalTok{(woe\_temp\_power, }\AttributeTok{interactive =} \ConstantTok{FALSE}\NormalTok{)}
\NormalTok{\}}
\NormalTok{(plots\_power[[}\DecValTok{6}\NormalTok{]] }\SpecialCharTok{|}\NormalTok{ plots\_power[[}\DecValTok{8}\NormalTok{]]) }\SpecialCharTok{/}\NormalTok{ (plots\_power[[}\DecValTok{13}\NormalTok{]] }\SpecialCharTok{|}\NormalTok{ plots\_power[[}\DecValTok{23}\NormalTok{]])}
\end{Highlighting}
\end{Shaded}

\begin{center}\includegraphics{project_MD_files/figure-latex/rebuild_bins_power-1} \end{center}

The power constribution is harder to interpret. We have the low-power
cars that are safer (lower claims), and high power that are more
dangerous (higher claims). The middle range is not as clear. We notice
two informative risk pockets (bins=12). A quick web search indicates
these interpretations \footnote{The interpretations are the result of
  google IA's search and should be taken with caution. A true market
  research and assessment of the hidden factors is outside the scope of
  this assignment.}:

\begin{itemize}
\tightlist
\item
  Power between 44 and 47: this is the minimum legal and practical power
  for modern highway-legal cars. Below this range, cars are typically
  underpowered for regular driving conditions.
\item
  Power between 50 and 51: this is widely considered the mechanical
  sweet spot to maintain highway speed withouf sacrificing fuel economy
  in urban driving. It is typically a marker of entry-level city cars.
\end{itemize}

Based on this interpretation, the power could be a proxy for the type of
car and its urban usage. We will use 8 bins to capture this bump risk
pattern without overcomplicating the model.

\begin{Shaded}
\begin{Highlighting}[]
\NormalTok{data}\SpecialCharTok{$}\NormalTok{POWER\_binned }\OtherTok{\textless{}{-}} \FunctionTok{cut}\NormalTok{(}
\NormalTok{    data}\SpecialCharTok{$}\NormalTok{POWER,   }
    \AttributeTok{breaks =} \FunctionTok{rebuild\_bins}\NormalTok{(result\_power, }\DecValTok{8}\NormalTok{),}
    \AttributeTok{include.lowest =} \ConstantTok{TRUE}\NormalTok{, }\AttributeTok{ordered\_result =} \ConstantTok{TRUE}\NormalTok{)}
\NormalTok{n\_bins\_named\_list[[}\StringTok{"POWER"}\NormalTok{]] }\OtherTok{\textless{}{-}} \DecValTok{8}
\end{Highlighting}
\end{Shaded}

We will end this Data Transformation section here now that we have bins
for both categorical and numerical variables.

\subsection{Data Exploration}\label{data-exploration}

\subsubsection{Graphical Approach}\label{graphical-approach}

Based on all created bins, we can visualize the claim frequency across
different bins for each variable. This is what we already did in our
analysis of the relation between the age of the driver and claim
frequency. Remember that the MLE (Maximum Likelihood Estimator) for the
claim frequency based on the hypothesis that the claims follow a Poisson
distribution is the following:

\[NCLAIMS_{i} \sim \text{Poisson}(d_i\times \lambda) \text{, where } d_i \text{ is the exposure and } \lambda \text{ is the claim frequency rate.}\]
\[\hat{\lambda} = \frac{\sum_i NCLAIMS_i}{\sum_i d_i}\]
\[\hat{\lambda} \sim \text{Normal}\left(\hat{\lambda}, \frac{\hat{\lambda}}{\sum_i d_i}\right)\]

We find the variance as the inverse of the fisher information (itself
equal to \(-\ell''(\widehat{\lambda})\))
\[\text{Var}(\widehat{\lambda}) \approx \mathcal{I}(\widehat{\lambda})^{-1} = \frac{\widehat{\lambda}}{\sum_i d_i}.\]

Alternatively, with the gaussian approximation on \(NCLAIMS_i\), we
have:
\[NCLAIMS_i \sim \text{Normal}(d_i \times \lambda, d_i \times \lambda)\]
\[\hat{\lambda} = \frac{\sum_i NCLAIMS_i}{\sum_i d_i}\]
\[\hat{\lambda} \sim \text{Normal}\left(\frac{\sum_i d_i \lambda}{\sum_i d_i}, \frac{\sum_i d_i \lambda}{(\sum_i d_i)^2}\right)\]
\[\hat{\lambda} \sim \text{Normal}\left(\lambda, \frac{\lambda}{\sum_i d_i}\right)\]

Based on this mathematical understanding, we can build a main plot which
shows exposures, claims, and claim frequencies with their confidence
intervals (95 and 99\%) across different bins for each variable.

\begin{Shaded}
\begin{Highlighting}[]
\NormalTok{mean\_claim\_frequency }\OtherTok{\textless{}{-}} \FunctionTok{sum}\NormalTok{(data}\SpecialCharTok{$}\NormalTok{NCLAIMS) }\SpecialCharTok{/} \FunctionTok{sum}\NormalTok{(data}\SpecialCharTok{$}\NormalTok{EXP)}
\NormalTok{plot\_claim\_freq }\OtherTok{\textless{}{-}} \ControlFlowTok{function}\NormalTok{(data, variable) \{}
\NormalTok{    df\_grouped }\OtherTok{\textless{}{-}}\NormalTok{ data }\SpecialCharTok{\%\textgreater{}\%}
        \FunctionTok{mutate}\NormalTok{(}\AttributeTok{bin =}\NormalTok{ data[[variable]]) }\SpecialCharTok{\%\textgreater{}\%}
        \FunctionTok{group\_by}\NormalTok{(bin) }\SpecialCharTok{\%\textgreater{}\%} 
        \FunctionTok{summarise}\NormalTok{(}
            \AttributeTok{sum\_exp =} \FunctionTok{sum}\NormalTok{(EXP),}
            \AttributeTok{sum\_claims =} \FunctionTok{sum}\NormalTok{(NCLAIMS),}
            \AttributeTok{lambda =}\NormalTok{ sum\_claims }\SpecialCharTok{/}\NormalTok{ sum\_exp,}
            \AttributeTok{sd\_lambda =} \FunctionTok{sqrt}\NormalTok{(lambda }\SpecialCharTok{/}\NormalTok{ sum\_exp),}
            \AttributeTok{lambda\_95\_lower =}\NormalTok{ lambda }\SpecialCharTok{{-}} \FloatTok{1.96} \SpecialCharTok{*}\NormalTok{ sd\_lambda,}
            \AttributeTok{lambda\_95\_upper =}\NormalTok{ lambda }\SpecialCharTok{+} \FloatTok{1.96} \SpecialCharTok{*}\NormalTok{ sd\_lambda,}
            \AttributeTok{lambda\_99\_lower =}\NormalTok{ lambda }\SpecialCharTok{{-}} \FloatTok{2.58} \SpecialCharTok{*}\NormalTok{ sd\_lambda,}
            \AttributeTok{lambda\_99\_upper =}\NormalTok{ lambda }\SpecialCharTok{+} \FloatTok{2.58} \SpecialCharTok{*}\NormalTok{ sd\_lambda}
\NormalTok{        )}
\NormalTok{    df\_grouped}\SpecialCharTok{$}\NormalTok{sum\_no\_claims }\OtherTok{\textless{}{-}}\NormalTok{ df\_grouped}\SpecialCharTok{$}\NormalTok{sum\_exp }\SpecialCharTok{{-}}\NormalTok{ df\_grouped}\SpecialCharTok{$}\NormalTok{sum\_claims}
\NormalTok{    df\_counts }\OtherTok{\textless{}{-}} \FunctionTok{bind\_rows}\NormalTok{(}
        \FunctionTok{mutate}\NormalTok{(df\_grouped, }\AttributeTok{bin =}\NormalTok{ bin, }\AttributeTok{value =}\NormalTok{ sum\_claims, }\AttributeTok{metric =} \StringTok{"Claims"}\NormalTok{),}
        \FunctionTok{mutate}\NormalTok{(df\_grouped, }\AttributeTok{bin =}\NormalTok{ bin, }\AttributeTok{value =}\NormalTok{ sum\_no\_claims, }\AttributeTok{metric =} \StringTok{"No Claims"}\NormalTok{)}
\NormalTok{    )}
\NormalTok{    p\_counts }\OtherTok{\textless{}{-}} \FunctionTok{ggplot}\NormalTok{(df\_counts, }\FunctionTok{aes}\NormalTok{(}\AttributeTok{x =}\NormalTok{ bin, }\AttributeTok{y =}\NormalTok{ value, }\AttributeTok{fill =}\NormalTok{ metric)) }\SpecialCharTok{+}
        \FunctionTok{geom\_bar}\NormalTok{(}\AttributeTok{stat =} \StringTok{"identity"}\NormalTok{) }\SpecialCharTok{+} \FunctionTok{theme\_minimal}\NormalTok{() }\SpecialCharTok{+}
        \FunctionTok{scale\_fill\_manual}\NormalTok{(}\AttributeTok{values =} \FunctionTok{c}\NormalTok{(}\StringTok{"Claims"} \OtherTok{=} \StringTok{"firebrick"}\NormalTok{, }\StringTok{"No Claims"} \OtherTok{=} \StringTok{"lightgray"}\NormalTok{)) }\SpecialCharTok{+}
        \FunctionTok{labs}\NormalTok{(}\AttributeTok{x =}\NormalTok{ variable, }\AttributeTok{y =} \StringTok{"Count"}\NormalTok{, }\AttributeTok{fill =} \StringTok{"Metric"}\NormalTok{) }\SpecialCharTok{+}
        \FunctionTok{theme}\NormalTok{(}\AttributeTok{axis.text.x =} \FunctionTok{element\_text}\NormalTok{(}\AttributeTok{angle =} \DecValTok{45}\NormalTok{, }\AttributeTok{hjust =} \DecValTok{1}\NormalTok{)) }\SpecialCharTok{+}
        \FunctionTok{ggtitle}\NormalTok{(}\StringTok{"Exposure Split"}\NormalTok{)}
\NormalTok{    p\_freq }\OtherTok{\textless{}{-}} \FunctionTok{ggplot}\NormalTok{(df\_grouped, }\FunctionTok{aes}\NormalTok{(}\AttributeTok{x =}\NormalTok{ bin, }\AttributeTok{group =} \DecValTok{1}\NormalTok{)) }\SpecialCharTok{+}
        \FunctionTok{geom\_ribbon}\NormalTok{(}
            \FunctionTok{aes}\NormalTok{(}\AttributeTok{ymin =}\NormalTok{ lambda\_99\_lower, }\AttributeTok{ymax =}\NormalTok{ lambda\_99\_upper),}
            \AttributeTok{fill =} \StringTok{"gray"}\NormalTok{, }\AttributeTok{alpha =} \FloatTok{0.3}
\NormalTok{        ) }\SpecialCharTok{+}
        \FunctionTok{geom\_ribbon}\NormalTok{(}
            \FunctionTok{aes}\NormalTok{(}\AttributeTok{ymin =}\NormalTok{ lambda\_95\_lower, }\AttributeTok{ymax =}\NormalTok{ lambda\_95\_upper),}
            \AttributeTok{fill =} \StringTok{"lightblue"}\NormalTok{, }\AttributeTok{alpha =} \FloatTok{0.5}
\NormalTok{        ) }\SpecialCharTok{+}
        \FunctionTok{geom\_line}\NormalTok{(}\FunctionTok{aes}\NormalTok{(}\AttributeTok{y =}\NormalTok{ lambda), }\AttributeTok{color =} \StringTok{"black"}\NormalTok{) }\SpecialCharTok{+}
        \FunctionTok{geom\_point}\NormalTok{(}\FunctionTok{aes}\NormalTok{(}\AttributeTok{y =}\NormalTok{ lambda), }\AttributeTok{color =} \StringTok{"black"}\NormalTok{) }\SpecialCharTok{+}
        \FunctionTok{geom\_hline}\NormalTok{(}\AttributeTok{yintercept =}\NormalTok{ mean\_claim\_frequency, }\AttributeTok{color =} \StringTok{"red"}\NormalTok{, }\AttributeTok{linetype =} \StringTok{"dashed"}\NormalTok{) }\SpecialCharTok{+}
        \FunctionTok{theme\_minimal}\NormalTok{() }\SpecialCharTok{+}
        \FunctionTok{labs}\NormalTok{(}
            \AttributeTok{x =}\NormalTok{ variable,}
            \AttributeTok{y =} \StringTok{"Claim Frequency (Lambda)"}\NormalTok{,}
\NormalTok{        ) }\SpecialCharTok{+}
        \FunctionTok{theme}\NormalTok{(}\AttributeTok{axis.text.x =} \FunctionTok{element\_text}\NormalTok{(}\AttributeTok{angle =} \DecValTok{45}\NormalTok{, }\AttributeTok{hjust =} \DecValTok{1}\NormalTok{)) }\SpecialCharTok{+}
        \FunctionTok{ggtitle}\NormalTok{(}\StringTok{"Claim Frequency"}\NormalTok{)}
\NormalTok{    p\_all }\OtherTok{\textless{}{-}}\NormalTok{ p\_freq }\SpecialCharTok{/}\NormalTok{ p\_counts}
\NormalTok{    p\_all}
\NormalTok{\}}
\end{Highlighting}
\end{Shaded}

\begin{Shaded}
\begin{Highlighting}[]
\FunctionTok{plot\_claim\_freq}\NormalTok{(data, }\StringTok{"BM\_binned"}\NormalTok{)}
\end{Highlighting}
\end{Shaded}

\begin{center}\includegraphics{project_MD_files/figure-latex/test_bm_plot-1} \end{center}

For this first example on the Bonus-Malus binned variable, we see a
clear upward trend in the claim frequency as the Bonus-Malus level
increases. This is expected as risky behaviour is translated in higher
score on the BM scale. What is also interesting is that the claim
frequency is a lot more precise for the 0 to 1 bin as the exposure is
much higher compared to the higher BM levels. The confidence intervals
also enable us to quickly test hypothesis that the claim frequency for a
given bin is significantly different from the overall mean claim
frequency (in red) or from other bins within the same
variable\footnote{If the value is outside of the blue range, we can
  reject the hypothesis that the claim frequency for that bin is equal
  to the value with 95\% certainty. For the grey interval, we can reject
  the hypothesis with 99\% certainty.}.

The same approach can be used for all variables:

\paragraph{Categorical Variables}\label{categorical-variables-1}

\begin{Shaded}
\begin{Highlighting}[]
\ControlFlowTok{for}\NormalTok{ (cat\_var }\ControlFlowTok{in} \FunctionTok{c}\NormalTok{(}\StringTok{"COVERAGE"}\NormalTok{, }\StringTok{"FUEL"}\NormalTok{, }\StringTok{"USE"}\NormalTok{, }\StringTok{"FLEET"}\NormalTok{, }\StringTok{"SEX"}\NormalTok{)) \{}
    \FunctionTok{cat}\NormalTok{(}\FunctionTok{paste0}\NormalTok{(}\StringTok{"\textgreater{} "}\NormalTok{, cat\_var, }\StringTok{"}\SpecialCharTok{\textbackslash{}n}\StringTok{"}\NormalTok{))}
    \FunctionTok{plot}\NormalTok{(}\FunctionTok{plot\_claim\_freq}\NormalTok{(data, cat\_var))}
\NormalTok{\}}
\end{Highlighting}
\end{Shaded}

\begin{verbatim}
> COVERAGE
\end{verbatim}

\begin{center}\includegraphics{project_MD_files/figure-latex/show_cat-1} \end{center}

\begin{verbatim}
> FUEL
\end{verbatim}

\begin{center}\includegraphics{project_MD_files/figure-latex/show_cat-2} \end{center}

\begin{verbatim}
> USE
\end{verbatim}

\begin{center}\includegraphics{project_MD_files/figure-latex/show_cat-3} \end{center}

\begin{verbatim}
> FLEET
\end{verbatim}

\begin{center}\includegraphics{project_MD_files/figure-latex/show_cat-4} \end{center}

\begin{verbatim}
> SEX
\end{verbatim}

\begin{center}\includegraphics{project_MD_files/figure-latex/show_cat-5} \end{center}

Based on the graphs, here are the preliminary conclusions for univariate
analysis of categorical variables:

\begin{itemize}
\tightlist
\item
  \texttt{COVERAGE}
\end{itemize}

The claim frequency between partial and full omnium is not statistically
different. However both are different from third-party liability only.
It is interesting to see that the claim frequency is higher for these
who choose lesser coverage. This can be interpreted as a proxy for
risk-awareness, financial education, or even purchasing power. Based on
this exploration, we will merge the partial and full omnium to focus on
the significant difference.

\begin{Shaded}
\begin{Highlighting}[]
\NormalTok{data}\SpecialCharTok{$}\NormalTok{COVERAGE\_BIN }\OtherTok{\textless{}{-}} \FunctionTok{ifelse}\NormalTok{(data}\SpecialCharTok{$}\NormalTok{COVERAGE }\SpecialCharTok{==} \StringTok{"TPL"}\NormalTok{, }\StringTok{"TPL"}\NormalTok{, }\StringTok{"OMN"}\NormalTok{)}
\NormalTok{data}\SpecialCharTok{$}\NormalTok{COVERAGE\_BIN }\OtherTok{\textless{}{-}} \FunctionTok{as.factor}\NormalTok{(data}\SpecialCharTok{$}\NormalTok{COVERAGE\_BIN) }\SpecialCharTok{\%\textgreater{}\%} \FunctionTok{fct\_infreq}\NormalTok{()}
\FunctionTok{plot\_claim\_freq}\NormalTok{(data, }\StringTok{"COVERAGE\_BIN"}\NormalTok{)}
\end{Highlighting}
\end{Shaded}

\begin{center}\includegraphics{project_MD_files/figure-latex/adapt_coverage-1} \end{center}

\begin{itemize}
\tightlist
\item
  \texttt{FUEL}
\end{itemize}

Diesel vehicles have significantly higher claim frequency compared to
petrol vehicles. There is no obvious explanation for this statistical
difference.

\begin{itemize}
\tightlist
\item
  \texttt{USE}
\end{itemize}

Private and professional use show no statistical differences from one
another, nor from the mean. We will not include this variable in our
models to avoid noise. Note that the proportion of policies related to
work is very low, this has three consequences:

\begin{enumerate}
\def\labelenumi{\arabic{enumi}.}
\tightlist
\item
  The private use looks very similar to the average as it represents the
  majority of the policies.
\item
  The work use is to variable to be statistically significant due to its
  low representation.
\item
  The variable is too unbalanced to provide reliable insights for
  modeling purposes.
\end{enumerate}

\begin{itemize}
\tightlist
\item
  \texttt{FLEET}
\end{itemize}

Let's compare the graph for fleet to the one of \texttt{USE}. Here
again, we have a very unbalanced variable.

\begin{enumerate}
\def\labelenumi{\arabic{enumi}.}
\tightlist
\item
  The majority of policies are non-fleet, which makes the claim
  frequency for this category very similar to the overall average.
\item
  Despite its low representation, the fleet category is significantly
  less risky than the non-fleet category.
\item
  The variable has predictive potential despite the high volatility of
  the non-fleet category.
\end{enumerate}

\begin{itemize}
\tightlist
\item
  \texttt{SEX}
\end{itemize}

Similarly to the interpretation for the \texttt{FLEET}, there is a
significant difference in claim frequency even if one category is less
represented than the other. In this case, the unbalance is less strong,
the males are thus slightly further from the average and females are
significantly more prone to claims than males. This variable also has
predictive potential for modeling purposes.

\paragraph{Numerical Variables}\label{numerical-variables}

Let us now have a look at the binned numerical variables other than the
bonus-malus scale that we already interpreted.

\begin{Shaded}
\begin{Highlighting}[]
\ControlFlowTok{for}\NormalTok{ (cat\_var }\ControlFlowTok{in} \FunctionTok{c}\NormalTok{(}\StringTok{"AGEPH\_binned"}\NormalTok{, }\StringTok{"AGEC\_binned"}\NormalTok{, }\StringTok{"POWER\_binned"}\NormalTok{)) \{}
    \FunctionTok{cat}\NormalTok{(}\FunctionTok{paste0}\NormalTok{(}\StringTok{"\textgreater{} "}\NormalTok{, cat\_var, }\StringTok{"}\SpecialCharTok{\textbackslash{}n}\StringTok{"}\NormalTok{))}
    \FunctionTok{plot}\NormalTok{(}\FunctionTok{plot\_claim\_freq}\NormalTok{(data, cat\_var))}
\NormalTok{\}}
\end{Highlighting}
\end{Shaded}

\begin{verbatim}
> AGEPH_binned
\end{verbatim}

\begin{center}\includegraphics{project_MD_files/figure-latex/show_num-1} \end{center}

\begin{verbatim}
> AGEC_binned
\end{verbatim}

\begin{center}\includegraphics{project_MD_files/figure-latex/show_num-2} \end{center}

\begin{verbatim}
> POWER_binned
\end{verbatim}

\begin{center}\includegraphics{project_MD_files/figure-latex/show_num-3} \end{center}

\begin{itemize}
\tightlist
\item
  \texttt{AGEPH\_binned}
\end{itemize}

We see the risk diminution as the age of the policyholder increases as
we identified with our binning algorithm. However, we see that the last
three are not significantly different from one another in terms of claim
frequency. We will come back to this in the statistical approach section
and see with an anova test if the 7 bins are significantly better than a
reduced number of bins.

\begin{itemize}
\tightlist
\item
  \texttt{AGEC\_binned}
\end{itemize}

For the age of the car, we can confirm our previous observations that
brand new cars (0-1 years old) have more claims. We can also confirm the
lower risk for cars between 1 and 3 years old and above 15 years. The
remaining question is whether the bins 1 to 2 and 2 to 3 should be
merged in a single bin. We will revisit this question in the statistical
approach section to determine if merging these bins significantly
affects the model performance.

\begin{itemize}
\tightlist
\item
  \texttt{POWER\_binned}
\end{itemize}

Thanks to our binning algorithm, we correctly identified the low-risk
for power lower than 33 kilowatts. We also correctly identified the very
bump-shaped peak in risk for power 44 to 51 kilowatts. It is not clear
whether the last 3 bins should be merged or kept separate, as the risk
pattern is less pronounced in this range. We will also investigate this
in the statistical approach section.

\subsubsection{Statistical Approach}\label{statistical-approach}

We already addressed the statistical significance by visually checking
the confidence intervals around the claim frequencies. In this
statistical data exploration, we will add two anova layers onto that
visual inspection:

\paragraph{Global Anova Test:}\label{global-anova-test}

We will compare the deviance of the model with the variable as single
predictor to the null model with only the offset term. This will allow
us to test if the variable significantly improves the model fit:

\(H_0: \Delta D = D_0 - D_1 = 0\) \(H_1: \Delta D = D_0 - D_1 > 0\)

For our test statistic we have \(\Delta D \sim \chi^2_{df}\), where
\(df\) is the difference in the number of parameters between the null
and alternative models.

Our \(D_0\) represents the deviance of the null model with only the
offset term: \texttt{glm(NCLAIMS\ \textasciitilde{}\ offset(log(EXP)))}.
The \(D_1\) represents the deviance of the model with the variable as a
single predictor:
\texttt{glm(NCLAIMS\ \textasciitilde{}\ offset(log(EXP))\ +\ get(var))}.

\begin{Shaded}
\begin{Highlighting}[]
\NormalTok{m0 }\OtherTok{\textless{}{-}} \FunctionTok{glm}\NormalTok{(NCLAIMS }\SpecialCharTok{\textasciitilde{}} \FunctionTok{offset}\NormalTok{(}\FunctionTok{log}\NormalTok{(EXP)), }\AttributeTok{data =}\NormalTok{ data, }\AttributeTok{family =} \FunctionTok{poisson}\NormalTok{(}\AttributeTok{link =} \StringTok{"log"}\NormalTok{))}
\FunctionTok{cat}\NormalTok{(}\FunctionTok{sprintf}\NormalTok{(}\StringTok{"Because we use Poisson with its canonic link, note exp(Beta0) = \%f is equal to the empiric claim frequency = \%f"}\NormalTok{, }\FunctionTok{exp}\NormalTok{(}\FunctionTok{coef}\NormalTok{(m0)[}\DecValTok{1}\NormalTok{]), }\FunctionTok{sum}\NormalTok{(data}\SpecialCharTok{$}\NormalTok{NCLAIMS)}\SpecialCharTok{/}\FunctionTok{sum}\NormalTok{(data}\SpecialCharTok{$}\NormalTok{EXP)))}
\end{Highlighting}
\end{Shaded}

\begin{verbatim}
Because we use Poisson with its canonic link, note exp(Beta0) = 0.138624 is equal to the empiric claim frequency = 0.138624
\end{verbatim}

Let's now compare this null model to the models with each predictor
variable individually.

\begin{Shaded}
\begin{Highlighting}[]
\NormalTok{test\_level }\OtherTok{\textless{}{-}} \FloatTok{0.05}
\ControlFlowTok{for}\NormalTok{ (var }\ControlFlowTok{in} \FunctionTok{c}\NormalTok{(}\StringTok{"COVERAGE\_BIN"}\NormalTok{, }\StringTok{"FUEL"}\NormalTok{, }\StringTok{"USE"}\NormalTok{, }\StringTok{"FLEET"}\NormalTok{, }\StringTok{"SEX"}\NormalTok{, }\StringTok{"BM\_binned"}\NormalTok{, }\StringTok{"AGEPH\_binned"}\NormalTok{, }\StringTok{"AGEC\_binned"}\NormalTok{, }\StringTok{"POWER\_binned"}\NormalTok{)) \{}
\NormalTok{    m1 }\OtherTok{\textless{}{-}} \FunctionTok{glm}\NormalTok{(NCLAIMS }\SpecialCharTok{\textasciitilde{}} \FunctionTok{offset}\NormalTok{(}\FunctionTok{log}\NormalTok{(EXP)) }\SpecialCharTok{+} \FunctionTok{get}\NormalTok{(var), }\AttributeTok{data =}\NormalTok{ data, }\AttributeTok{family =} \FunctionTok{poisson}\NormalTok{(}\AttributeTok{link =} \StringTok{"log"}\NormalTok{))}
\NormalTok{    test\_result }\OtherTok{\textless{}{-}} \FunctionTok{anova}\NormalTok{(m0, m1, }\AttributeTok{test =} \StringTok{"Chisq"}\NormalTok{)}
    \FunctionTok{cat}\NormalTok{(}\FunctionTok{sprintf}\NormalTok{(}\StringTok{"Is variable \%s significant? \%s}\SpecialCharTok{\textbackslash{}n}\StringTok{"}\NormalTok{, var, }\FunctionTok{ifelse}\NormalTok{(test\_result}\SpecialCharTok{$}\StringTok{\textasciigrave{}}\AttributeTok{Pr(\textgreater{}Chi)}\StringTok{\textasciigrave{}}\NormalTok{[}\DecValTok{2}\NormalTok{] }\SpecialCharTok{\textless{}}\NormalTok{ test\_level, }\StringTok{"Yes"}\NormalTok{, }\StringTok{"No"}\NormalTok{)))}
\NormalTok{\}}
\end{Highlighting}
\end{Shaded}

\begin{verbatim}
Is variable COVERAGE_BIN significant? Yes
Is variable FUEL significant? Yes
Is variable USE significant? No
Is variable FLEET significant? Yes
Is variable SEX significant? Yes
Is variable BM_binned significant? Yes
Is variable AGEPH_binned significant? Yes
Is variable AGEC_binned significant? Yes
Is variable POWER_binned significant? Yes
\end{verbatim}

With this first global anova test, we confirm our graphical approach,
all variables should be kept except for \texttt{USE}.

\paragraph{Merged bins}\label{merged-bins}

Let's take the example of \texttt{COVERAGE} where we decided to merge
the two omnium levels. Can we justify this statistically with the ANOVA
test instead of simply with the confidence intervals?

\begin{Shaded}
\begin{Highlighting}[]
\NormalTok{m0 }\OtherTok{\textless{}{-}} \FunctionTok{glm}\NormalTok{(NCLAIMS }\SpecialCharTok{\textasciitilde{}} \FunctionTok{offset}\NormalTok{(}\FunctionTok{log}\NormalTok{(EXP)) }\SpecialCharTok{+}\NormalTok{ COVERAGE\_BIN, }\AttributeTok{data =}\NormalTok{ data, }\AttributeTok{family =} \FunctionTok{poisson}\NormalTok{(}\AttributeTok{link =} \StringTok{"log"}\NormalTok{))}
\NormalTok{m1 }\OtherTok{\textless{}{-}} \FunctionTok{glm}\NormalTok{(NCLAIMS }\SpecialCharTok{\textasciitilde{}} \FunctionTok{offset}\NormalTok{(}\FunctionTok{log}\NormalTok{(EXP)) }\SpecialCharTok{+}\NormalTok{ COVERAGE, }\AttributeTok{data =}\NormalTok{ data, }\AttributeTok{family =} \FunctionTok{poisson}\NormalTok{(}\AttributeTok{link =} \StringTok{"log"}\NormalTok{))}
\NormalTok{test\_result }\OtherTok{\textless{}{-}} \FunctionTok{anova}\NormalTok{(m0, m1, }\AttributeTok{test =} \StringTok{"Chisq"}\NormalTok{)}
\NormalTok{pval }\OtherTok{\textless{}{-}}\NormalTok{ test\_result}\SpecialCharTok{$}\StringTok{\textasciigrave{}}\AttributeTok{Pr(\textgreater{}Chi)}\StringTok{\textasciigrave{}}\NormalTok{[}\DecValTok{2}\NormalTok{]}
\FunctionTok{cat}\NormalTok{(}\FunctionTok{sprintf}\NormalTok{(}\StringTok{"Was it justified to merge the COVERAGE bins? \%s"}\NormalTok{, }\FunctionTok{ifelse}\NormalTok{(pval }\SpecialCharTok{\textless{}} \FloatTok{0.05}\NormalTok{, }\StringTok{"No"}\NormalTok{, }\StringTok{"Yes"}\NormalTok{)))}
\end{Highlighting}
\end{Shaded}

\begin{verbatim}
Was it justified to merge the COVERAGE bins? No
\end{verbatim}

\begin{Shaded}
\begin{Highlighting}[]
\FunctionTok{cat}\NormalTok{(}\FunctionTok{sprintf}\NormalTok{(}\StringTok{"Are you 99\%\% sure that it was not justified? \%s"}\NormalTok{, }\FunctionTok{ifelse}\NormalTok{(pval }\SpecialCharTok{\textless{}} \FloatTok{0.01}\NormalTok{, }\StringTok{"Yes"}\NormalTok{, }\StringTok{"No"}\NormalTok{)))}
\end{Highlighting}
\end{Shaded}

\begin{verbatim}
Are you 99% sure that it was not justified? No
\end{verbatim}

Because the p-value is 0.0469242, we can conclude that we reject the
hypothesis that the deviance was different between the merged and
unmerged COVERAGE bins at a 95\% confidence level but not at a 99\%
confidence level. It was therefore not a bad idea to merge the COVERAGE
bins based on our graphical analysis based on confidence intervals.

We can apply the same reasoning to clarify whether or not we should
merge:

\begin{itemize}
\tightlist
\item
  The last 3 bins of \texttt{AGEPH\_binned}
\end{itemize}

\begin{Shaded}
\begin{Highlighting}[]
\NormalTok{data}\SpecialCharTok{$}\NormalTok{AGEPH\_binned\_merged }\OtherTok{\textless{}{-}} \FunctionTok{fct\_collapse}\NormalTok{(}
\NormalTok{    data}\SpecialCharTok{$}\NormalTok{AGEPH\_binned, }
    \StringTok{"(66,95]"} \OtherTok{=} \FunctionTok{c}\NormalTok{(}\StringTok{"(66,68]"}\NormalTok{, }\StringTok{"(68,71]"}\NormalTok{, }\StringTok{"(71,95]"}\NormalTok{)}
\NormalTok{)}

\NormalTok{m0 }\OtherTok{\textless{}{-}} \FunctionTok{glm}\NormalTok{(NCLAIMS }\SpecialCharTok{\textasciitilde{}} \FunctionTok{offset}\NormalTok{(}\FunctionTok{log}\NormalTok{(EXP)) }\SpecialCharTok{+}\NormalTok{ AGEPH\_binned, }\AttributeTok{data =}\NormalTok{ data, }\AttributeTok{family =} \FunctionTok{poisson}\NormalTok{(}\AttributeTok{link =} \StringTok{"log"}\NormalTok{))}
\NormalTok{m1 }\OtherTok{\textless{}{-}} \FunctionTok{glm}\NormalTok{(NCLAIMS }\SpecialCharTok{\textasciitilde{}} \FunctionTok{offset}\NormalTok{(}\FunctionTok{log}\NormalTok{(EXP)) }\SpecialCharTok{+}\NormalTok{ AGEPH\_binned\_merged, }\AttributeTok{data =}\NormalTok{ data, }\AttributeTok{family =} \FunctionTok{poisson}\NormalTok{(}\AttributeTok{link =} \StringTok{"log"}\NormalTok{))}
\NormalTok{test\_result }\OtherTok{\textless{}{-}} \FunctionTok{anova}\NormalTok{(m0, m1, }\AttributeTok{test =} \StringTok{"Chisq"}\NormalTok{)}
\NormalTok{pval }\OtherTok{\textless{}{-}}\NormalTok{ test\_result}\SpecialCharTok{$}\StringTok{\textasciigrave{}}\AttributeTok{Pr(\textgreater{}Chi)}\StringTok{\textasciigrave{}}\NormalTok{[}\DecValTok{2}\NormalTok{]}

\ControlFlowTok{if}\NormalTok{(pval }\SpecialCharTok{\textless{}} \FloatTok{0.05}\NormalTok{) \{}
    \FunctionTok{cat}\NormalTok{(}\FunctionTok{sprintf}\NormalTok{(}\StringTok{"Merging the last 3 bins of AGEPH\_binned was not justified (p{-}value = \%f).}\SpecialCharTok{\textbackslash{}n}\StringTok{"}\NormalTok{, pval))}
\NormalTok{    data}\SpecialCharTok{$}\NormalTok{AGEPH\_binned\_merged }\OtherTok{\textless{}{-}} \ConstantTok{NULL}
\NormalTok{\} }\ControlFlowTok{else}\NormalTok{ \{}
    \FunctionTok{cat}\NormalTok{(}\FunctionTok{sprintf}\NormalTok{(}\StringTok{"Merging the last 3 bins of AGEPH\_binned was justified (p{-}value = \%f).}\SpecialCharTok{\textbackslash{}n}\StringTok{"}\NormalTok{, pval))}
    \FunctionTok{plot}\NormalTok{(}\FunctionTok{plot\_claim\_freq}\NormalTok{(data, }\StringTok{"AGEPH\_binned\_merged"}\NormalTok{))}
\NormalTok{\}}
\end{Highlighting}
\end{Shaded}

\begin{verbatim}
Merging the last 3 bins of AGEPH_binned was justified (p-value = 0.750730).
\end{verbatim}

\begin{center}\includegraphics{project_MD_files/figure-latex/merge_agephb-1} \end{center}

\begin{itemize}
\tightlist
\item
  The bins 1-2 and 2-3 of \texttt{AGEC\_binned}
\end{itemize}

\begin{Shaded}
\begin{Highlighting}[]
\NormalTok{data}\SpecialCharTok{$}\NormalTok{AGEC\_binned\_merged }\OtherTok{\textless{}{-}} \FunctionTok{fct\_collapse}\NormalTok{(}
\NormalTok{    data}\SpecialCharTok{$}\NormalTok{AGEC\_binned, }
    \StringTok{"(1,3]"} \OtherTok{=} \FunctionTok{c}\NormalTok{(}\StringTok{"(1,2]"}\NormalTok{, }\StringTok{"(2,3]"}\NormalTok{)}
\NormalTok{)}

\NormalTok{m0 }\OtherTok{\textless{}{-}} \FunctionTok{glm}\NormalTok{(NCLAIMS }\SpecialCharTok{\textasciitilde{}} \FunctionTok{offset}\NormalTok{(}\FunctionTok{log}\NormalTok{(EXP)) }\SpecialCharTok{+}\NormalTok{ AGEC\_binned, }\AttributeTok{data =}\NormalTok{ data, }\AttributeTok{family =} \FunctionTok{poisson}\NormalTok{(}\AttributeTok{link =} \StringTok{"log"}\NormalTok{))}
\NormalTok{m1 }\OtherTok{\textless{}{-}} \FunctionTok{glm}\NormalTok{(NCLAIMS }\SpecialCharTok{\textasciitilde{}} \FunctionTok{offset}\NormalTok{(}\FunctionTok{log}\NormalTok{(EXP)) }\SpecialCharTok{+}\NormalTok{ AGEC\_binned\_merged, }\AttributeTok{data =}\NormalTok{ data, }\AttributeTok{family =} \FunctionTok{poisson}\NormalTok{(}\AttributeTok{link =} \StringTok{"log"}\NormalTok{))}
\NormalTok{test\_result }\OtherTok{\textless{}{-}} \FunctionTok{anova}\NormalTok{(m0, m1, }\AttributeTok{test =} \StringTok{"Chisq"}\NormalTok{)}
\NormalTok{pval }\OtherTok{\textless{}{-}}\NormalTok{ test\_result}\SpecialCharTok{$}\StringTok{\textasciigrave{}}\AttributeTok{Pr(\textgreater{}Chi)}\StringTok{\textasciigrave{}}\NormalTok{[}\DecValTok{2}\NormalTok{]}

\ControlFlowTok{if}\NormalTok{(pval }\SpecialCharTok{\textless{}} \FloatTok{0.05}\NormalTok{) \{}
    \FunctionTok{cat}\NormalTok{(}\FunctionTok{sprintf}\NormalTok{(}\StringTok{"Merging the second and third bins of AGEC\_binned was not justified (p{-}value = \%f).}\SpecialCharTok{\textbackslash{}n}\StringTok{"}\NormalTok{, pval))}
\NormalTok{    data}\SpecialCharTok{$}\NormalTok{AGEC\_binned\_merged }\OtherTok{\textless{}{-}} \ConstantTok{NULL}
\NormalTok{\} }\ControlFlowTok{else}\NormalTok{ \{}
    \FunctionTok{cat}\NormalTok{(}\FunctionTok{sprintf}\NormalTok{(}\StringTok{"Merging the second and third bins of AGEC\_binned was justified (p{-}value = \%f).}\SpecialCharTok{\textbackslash{}n}\StringTok{"}\NormalTok{, pval))}
    \FunctionTok{plot}\NormalTok{(}\FunctionTok{plot\_claim\_freq}\NormalTok{(data, }\StringTok{"AGEC\_binned\_merged"}\NormalTok{))}
\NormalTok{\}}
\end{Highlighting}
\end{Shaded}

\begin{verbatim}
Merging the second and third bins of AGEC_binned was not justified (p-value = 0.036014).
\end{verbatim}

\begin{itemize}
\tightlist
\item
  The last 3 bins of \texttt{POWER\_binned}
\end{itemize}

\begin{Shaded}
\begin{Highlighting}[]
\NormalTok{data}\SpecialCharTok{$}\NormalTok{POWER\_binned\_merged }\OtherTok{\textless{}{-}} \FunctionTok{fct\_collapse}\NormalTok{(}
\NormalTok{    data}\SpecialCharTok{$}\NormalTok{POWER\_binned, }
    \StringTok{"(76,243]"} \OtherTok{=} \FunctionTok{c}\NormalTok{(}\StringTok{"(76,83]"}\NormalTok{, }\StringTok{"(83,99]"}\NormalTok{, }\StringTok{"(99,243]"}\NormalTok{)}
\NormalTok{)}
\NormalTok{m0 }\OtherTok{\textless{}{-}} \FunctionTok{glm}\NormalTok{(NCLAIMS }\SpecialCharTok{\textasciitilde{}} \FunctionTok{offset}\NormalTok{(}\FunctionTok{log}\NormalTok{(EXP)) }\SpecialCharTok{+}\NormalTok{ POWER\_binned, }\AttributeTok{data =}\NormalTok{ data, }\AttributeTok{family =} \FunctionTok{poisson}\NormalTok{(}\AttributeTok{link =} \StringTok{"log"}\NormalTok{))}
\NormalTok{m1 }\OtherTok{\textless{}{-}} \FunctionTok{glm}\NormalTok{(NCLAIMS }\SpecialCharTok{\textasciitilde{}} \FunctionTok{offset}\NormalTok{(}\FunctionTok{log}\NormalTok{(EXP)) }\SpecialCharTok{+}\NormalTok{ POWER\_binned\_merged, }\AttributeTok{data =}\NormalTok{ data, }\AttributeTok{family =} \FunctionTok{poisson}\NormalTok{(}\AttributeTok{link =} \StringTok{"log"}\NormalTok{))}
\NormalTok{test\_result }\OtherTok{\textless{}{-}} \FunctionTok{anova}\NormalTok{(m0, m1, }\AttributeTok{test =} \StringTok{"Chisq"}\NormalTok{)}
\NormalTok{pval }\OtherTok{\textless{}{-}}\NormalTok{ test\_result}\SpecialCharTok{$}\StringTok{\textasciigrave{}}\AttributeTok{Pr(\textgreater{}Chi)}\StringTok{\textasciigrave{}}\NormalTok{[}\DecValTok{2}\NormalTok{]}

\ControlFlowTok{if}\NormalTok{(pval }\SpecialCharTok{\textless{}} \FloatTok{0.05}\NormalTok{) \{}
    \FunctionTok{cat}\NormalTok{(}\FunctionTok{sprintf}\NormalTok{(}\StringTok{"Merging the last 3 bins of POWER\_binned was not justified (p{-}value = \%f).}\SpecialCharTok{\textbackslash{}n}\StringTok{"}\NormalTok{, pval))}
\NormalTok{    data}\SpecialCharTok{$}\NormalTok{POWER\_binned\_merged }\OtherTok{\textless{}{-}} \ConstantTok{NULL}
\NormalTok{\} }\ControlFlowTok{else}\NormalTok{ \{}
    \FunctionTok{cat}\NormalTok{(}\FunctionTok{sprintf}\NormalTok{(}\StringTok{"Merging the last 3 bins of POWER\_binned was justified (p{-}value = \%f).}\SpecialCharTok{\textbackslash{}n}\StringTok{"}\NormalTok{, pval))}
    \FunctionTok{plot}\NormalTok{(}\FunctionTok{plot\_claim\_freq}\NormalTok{(data, }\StringTok{"POWER\_binned\_merged"}\NormalTok{))}
\NormalTok{\}}
\end{Highlighting}
\end{Shaded}

\begin{verbatim}
Merging the last 3 bins of POWER_binned was justified (p-value = 0.307144).
\end{verbatim}

\begin{center}\includegraphics{project_MD_files/figure-latex/merge_power-1} \end{center}

We have now have thoroughly examined the data we had available exced for
the \texttt{TOWN} and \texttt{PC} variables. We will come back to them
in the feature engineering section.

\subsection{Feature Engineering}\label{feature-engineering}

Besides the existing features, we can create additional features to add
information to our design matrix and potentially improve our models.

\begin{itemize}
\item
  \textbf{Geospatial Features}: We will first explore the existing
  \texttt{PC} and \texttt{TOWN} variables`. Note that these are simply
  categorical representations of the longitude and latitude for
  different granularities of location. We will then explore different
  granularities and use the spatial information to enrich our feature
  set.
\item
  \textbf{Polynomial Features}: As we will see in the modelling section,
  adding polynomial features can be a good way to capture non-linear
  relationships in our GLM models. This is an intermediate step towards
  the smoothing functions of Generalized Additive Models (GAMs). We can
  test the second and third degree polynomial features for continuous
  variables and bin the statistically significant features to create new
  categorical variables for our GLM. Note that this is less
  straighforward to explain to a customer. It is best used in technical
  premium modeling but should likely be removed from the commercial
  premium models.
\item
  \textbf{Interaction Features}: All the analysis we have done so far
  has focused on individual features. The statistical modelling (GLM,
  GAM) will allow us to capture the controlled effects of the features.
  However, some interactions between features are not captured by linear
  terms alone. We can create interaction features to explicitly model
  these effects.
\end{itemize}

\subsubsection{Geospatial Features}\label{geospatial-features}

Location can be considered at multiple levels of granularity. In
Belgium, the \texttt{National\ Institute\ of\ Statistics} provides a
\texttt{NIS\ code} nomenclature that uniquely identifies each location
at various granularities, from most specific to most general:

\begin{itemize}
\tightlist
\item
  \texttt{Statistical\ Sector}:(\textasciitilde20,000 sectors) that are
  granular enough to capture neighborhood variations within a postal
  code. This is typically interesting for sociodemographic
  (socioeconomic indexes, population density, etc.) analysis.
\item
  \texttt{Postal\ Code\ (PC)}:(\textasciitilde600 codes) a more
  aggregated level, capturing aggregated characteristics. This is a more
  practical level for many analyses to avoid missing data issues and
  respecting privacy concerns on limited datasets.
\item
  \texttt{Administrative\ Arrondissement}:(\textasciitilde40
  arrondissements) this might be the less informative level for us as it
  is very aggregated but is not a political division yet.
\item
  \texttt{Province}:(\textasciitilde10 provinces) the political
  implications should be considered at this level as provinces have a
  certain degree of autonomy. Note that Brussels-Capital does not have
  this level of granularity as it is a region.
\item
  \texttt{Region}:(3 regions) the highest level of granularity, the
  biggest political and administrative decision-making occurs at this
  level. It captures a different kind of information compared to the
  more localized levels.
\end{itemize}

Choosing the best level of analysis is not an easy task. It first
depends on data availability and portfolio size. For smaller portfolios
or when data is sparse, higher-level aggregations like provinces will
allow for more robust statistical analysis. Conversely, for larger
portfolios with sufficient data, more granular levels like statistical
sectors or postal codes will help better differentiate risk pockets.

Besides, granularities will often capture completly different underlying
risk. Let's clarify this with a climate risk example \footnote{This is
  not in the scope of our analysis but it is an important topic worth a
  small introduction nowadays.} using different granularities.

\begin{itemize}
\item
  \texttt{Latitude} and \texttt{Longitude}: These \textbf{numerical}
  features simple indicate the point on a map. If we have a large enough
  geographical coverage, it will help us capture meteorological
  patterns. For \textbf{physical risks}, such as floods or storms, we
  understand that a northern country will face different risk (hail
  breaking the windhshield, floods damaging the cars parked outdoors)
  compared to a southern country where heatwaves can cause damage to
  tires and paintwork. In a single country such as Belgium, it won't be
  very helpful as such but might capture other factors such as traffic
  density, urbanization, or quality of infrastructure. The
  \texttt{Statistical\ Sector} or \texttt{Postal\ Code} can also act as
  proxies for these underlying factors.
\item
  \texttt{Regions}: The \textbf{categorical} features will capture
  administrative and political divisions, which can be important for
  understanding regional policies and regulations. When we look at
  \textbf{transition risk}, this becomes particularly relevant. For
  example, when the Brussels-Capital Region restricted the access to old
  diesel vehicles, this became stranded assets for owners of such
  vehicles within the region. Regions can take other decisions that
  impact risk exposure, such as subsidising electric vehicles, promoting
  public transportation, or offering advantageous schemes for company
  cars. This will play a role in the behavioral patterns and risk
  exposure of vehicle owners within the region.
\end{itemize}

We should not ``choose'' between granularities as they capture different
underlying information. However, \textbf{multicollinearity} is a clear
risk when working with various granularities of a same feature which can
lead to unstable model with hard-to-interpret individual feature
effects. In our dataset, we already have the postcode (583 levels) and
town (578 levels) which seem to capture the same underlying geographical
information and the \texttt{LAT} and \texttt{LONG} features.

To avoid redundancy, we will drop the \texttt{TOWN} feature and keep the
\texttt{PC} which is part of the official NIS nomenclature. \texttt{LAT}
and \texttt{LONG} will also be kept as numerical input for the smoothing
functions we will use in the GAM models. To capture political and
administrative underlying factors, we will add the \texttt{Regions}
feature as a categorical variable. Finally, we will fetch the
demographic density for each postcode which will serve as a proxy for
urbanization and traffic exposure. We expect higher frequency of claims
and lower severity in areas with higher demographic density, as traffic
congestion may lead to more minor accidents rather than severe ones.
Conversely, we expect lower frequency and higher severity in more rural
areas. It will be interesting to control the predictive power of
postcode with demographic density to see if it has predictive power
beyond what is captured by urbanization and traffic exposure.

\paragraph{Dataset Enrichment}\label{dataset-enrichment}

Before the enrichment, let us drop the \texttt{TOWN} feature as it is
redundant but not perfecly aligned with the postcode, we'll replace it
with postcode labels.

\begin{Shaded}
\begin{Highlighting}[]
\NormalTok{data }\OtherTok{\textless{}{-}}\NormalTok{ data }\SpecialCharTok{\%\textgreater{}\%} \FunctionTok{select}\NormalTok{(}\SpecialCharTok{{-}}\NormalTok{TOWN)}
\end{Highlighting}
\end{Shaded}

There are two main options for enriching our dataset with geographical
information:

\textbf{Geospatial Enrichment}

\begin{itemize}
\tightlist
\item
  Geocode the \texttt{LAT} and \texttt{LONG} features to obtain precise
  geographical coordinates. We will use the Belgian Lambert 2008 (3812)
  projection\footnote{Coordinate Reference System are identifiers to
    explain how the projection from a 3D sphere to a 2D plane is
    performed. The Belgian Lambert 2008 (3812) projection is commonly
    used for visualizing geographical data in Belgium.} to be in line
  with the NIS standards.
\item
  Download the
  \href{https://statbel.fgov.be/en/open-data/statistical-sectors-2025}{shapefiles
  for Belgian administrative boundaries}.
\item
  Perform spatial matching between the geocoded points and the
  administrative boundaries to enrich the dataset with region, postcode,
  and statistical sector information.
\end{itemize}

\textbf{Code Hierarchy Enrichment}

\begin{itemize}
\tightlist
\item
  We can take a shortcut by finding matching tables that link postcodes
  to refnis, and use the code to assign regions, avoiding the geocoding
  and spatial matching steps.
\item
  Then we merge these tables with our dataset based on the postcode to
  enrich it with the corresponding region.
\end{itemize}

After this step, we can look for national statistics of population
density for 1997 (year of the dataset) in each municipality and merge
them with our dataset to further enrich it.

\begin{Shaded}
\begin{Highlighting}[]
\NormalTok{conv }\OtherTok{\textless{}{-}} \FunctionTok{read.csv}\NormalTok{(}\StringTok{"../data/postcode\_refnis\_conversion.csv"}\NormalTok{, }\AttributeTok{sep =} \StringTok{";"}\NormalTok{) }\SpecialCharTok{\%\textgreater{}\%} \FunctionTok{data.frame}\NormalTok{()}
\NormalTok{refnis }\OtherTok{\textless{}{-}} \FunctionTok{read.csv}\NormalTok{(}\StringTok{"../data/refnis\_matching.csv"}\NormalTok{, }\AttributeTok{sep =} \StringTok{";"}\NormalTok{) }\SpecialCharTok{\%\textgreater{}\%} \FunctionTok{data.frame}\NormalTok{()}
\NormalTok{pop }\OtherTok{\textless{}{-}} \FunctionTok{read.csv}\NormalTok{(}\StringTok{"../data/pop\_1997.csv"}\NormalTok{, }\AttributeTok{sep =} \StringTok{";"}\NormalTok{) }\SpecialCharTok{\%\textgreater{}\%} \FunctionTok{data.frame}\NormalTok{() }\SpecialCharTok{\%\textgreater{}\%}
    \FunctionTok{mutate}\NormalTok{(}\AttributeTok{NIS.code =} \FunctionTok{as.integer}\NormalTok{(NIS.code)) }\SpecialCharTok{\%\textgreater{}\%}
    \FunctionTok{filter}\NormalTok{(}\SpecialCharTok{!}\FunctionTok{is.na}\NormalTok{(NIS.code) }\SpecialCharTok{\&}\NormalTok{ NIS.code }\SpecialCharTok{!=} \StringTok{""}\NormalTok{)}

\NormalTok{dup\_conv }\OtherTok{\textless{}{-}}\NormalTok{ conv[}\FunctionTok{duplicated}\NormalTok{(conv}\SpecialCharTok{$}\NormalTok{Postal.code), ]}
\ControlFlowTok{if}\NormalTok{(}\FunctionTok{nrow}\NormalTok{(dup\_conv) }\SpecialCharTok{\textgreater{}} \DecValTok{0}\NormalTok{) \{}
    \FunctionTok{warning}\NormalTok{(}\StringTok{"There are duplicate postal codes in the conversion table."}\NormalTok{)}
\NormalTok{    conv }\SpecialCharTok{\%\textgreater{}\%} \FunctionTok{filter}\NormalTok{(Postal.code }\SpecialCharTok{\%in\%}\NormalTok{ dup\_conv}\SpecialCharTok{$}\NormalTok{Postal.code) }\SpecialCharTok{\%\textgreater{}\%} \FunctionTok{data.frame}\NormalTok{() }\SpecialCharTok{\%\textgreater{}\%}
    \FunctionTok{kable}\NormalTok{(}\AttributeTok{booktabs =} \ConstantTok{TRUE}\NormalTok{, }\AttributeTok{width =} \DecValTok{100}\NormalTok{, }\AttributeTok{caption =} \StringTok{"Duplicate Postal Codes in Conversion Table"}\NormalTok{)}
\NormalTok{\}}
\end{Highlighting}
\end{Shaded}

\begin{longtable}[]{@{}rrll@{}}
\caption{Duplicate Postal Codes in Conversion Table}\tabularnewline
\toprule\noalign{}
Postal.code & Refnis.code & Gemeentenaam & Nom.commune \\
\midrule\noalign{}
\endfirsthead
\toprule\noalign{}
Postal.code & Refnis.code & Gemeentenaam & Nom.commune \\
\midrule\noalign{}
\endhead
\bottomrule\noalign{}
\endlastfoot
1040 & 21004 & Brussel & Bruxelles \\
1040 & 21005 & Etterbeek & Etterbeek \\
1050 & 21004 & Brussel & Bruxelles \\
1050 & 21009 & Elsene & Ixelles \\
1804 & 23088 & Vilvoorde & Vilvorde \\
1804 & 23096 & Zemst & Zemst \\
\end{longtable}

\begin{Shaded}
\begin{Highlighting}[]
\NormalTok{dup\_refnis }\OtherTok{\textless{}{-}}\NormalTok{ refnis[}\FunctionTok{duplicated}\NormalTok{(refnis}\SpecialCharTok{$}\NormalTok{NIS\_MUN), ]}
\ControlFlowTok{if}\NormalTok{(}\FunctionTok{nrow}\NormalTok{(dup\_refnis) }\SpecialCharTok{\textgreater{}} \DecValTok{0}\NormalTok{) \{}
    \FunctionTok{warning}\NormalTok{(}\StringTok{"There are duplicate NIS\_MUN codes in the refnis table."}\NormalTok{)}
\NormalTok{    refnis }\SpecialCharTok{\%\textgreater{}\%} \FunctionTok{filter}\NormalTok{(NIS\_MUN }\SpecialCharTok{\%in\%}\NormalTok{ dup\_refnis}\SpecialCharTok{$}\NormalTok{NIS\_MUN) }\SpecialCharTok{\%\textgreater{}\%} \FunctionTok{data.frame}\NormalTok{() }\SpecialCharTok{\%\textgreater{}\%}
    \FunctionTok{kable}\NormalTok{(}\AttributeTok{booktabs =} \ConstantTok{TRUE}\NormalTok{, }\AttributeTok{width =} \DecValTok{100}\NormalTok{, }\AttributeTok{caption =} \StringTok{"Duplicate NIS\_MUN Codes in Refnis Table"}\NormalTok{)}
\NormalTok{\}}

\NormalTok{dup\_pop }\OtherTok{\textless{}{-}}\NormalTok{ pop[}\FunctionTok{duplicated}\NormalTok{(pop}\SpecialCharTok{$}\NormalTok{NIS.code), ]}
\ControlFlowTok{if}\NormalTok{(}\FunctionTok{nrow}\NormalTok{(dup\_pop) }\SpecialCharTok{\textgreater{}} \DecValTok{0}\NormalTok{) \{}
    \FunctionTok{warning}\NormalTok{(}\StringTok{"There are duplicate NIS codes in the population table."}\NormalTok{)}
\NormalTok{    pop }\SpecialCharTok{\%\textgreater{}\%} \FunctionTok{filter}\NormalTok{(NIS.code }\SpecialCharTok{\%in\%}\NormalTok{ dup\_pop}\SpecialCharTok{$}\NormalTok{NIS.code) }\SpecialCharTok{\%\textgreater{}\%} \FunctionTok{data.frame}\NormalTok{() }\SpecialCharTok{\%\textgreater{}\%}
    \FunctionTok{kable}\NormalTok{(}\AttributeTok{booktabs =} \ConstantTok{TRUE}\NormalTok{, }\AttributeTok{width =} \DecValTok{100}\NormalTok{, }\AttributeTok{caption =} \StringTok{"Duplicate NIS Codes in Population Table"}\NormalTok{)}
\NormalTok{\}}
\end{Highlighting}
\end{Shaded}

Before proceeding with the matching, we need to make sure that our
matching keys are unique. We see that in the conversion table, there are
duplicate postal codes, 1040 and 1050 can also be matched to Bruxelles
but the other matched provide more insight (Etterbeek and Ixelles). For
the postcode 1804, Vilvoorde is the corresponding municipality with some
overlap with nearby Zemst which we will drop.

\begin{Shaded}
\begin{Highlighting}[]
\NormalTok{conv\_matching }\OtherTok{\textless{}{-}}\NormalTok{ conv }\SpecialCharTok{\%\textgreater{}\%} \FunctionTok{filter}\NormalTok{(}\SpecialCharTok{!}\NormalTok{(}
\NormalTok{    (Postal.code }\SpecialCharTok{==} \DecValTok{1040} \SpecialCharTok{\&}\NormalTok{ Nom.commune }\SpecialCharTok{==} \StringTok{"Bruxelles"}\NormalTok{) }\SpecialCharTok{|}
\NormalTok{    (Postal.code }\SpecialCharTok{==} \DecValTok{1050} \SpecialCharTok{\&}\NormalTok{ Nom.commune }\SpecialCharTok{==} \StringTok{"Bruxelles"}\NormalTok{) }\SpecialCharTok{|}
\NormalTok{    (Postal.code }\SpecialCharTok{==} \DecValTok{1804} \SpecialCharTok{\&}\NormalTok{ Nom.commune }\SpecialCharTok{==} \StringTok{"Zemst"}\NormalTok{)}
\NormalTok{))}
\ControlFlowTok{if}\NormalTok{(}\FunctionTok{any}\NormalTok{(}\FunctionTok{duplicated}\NormalTok{(conv\_matching}\SpecialCharTok{$}\NormalTok{Postal.code))) \{}
    \FunctionTok{stop}\NormalTok{(}\StringTok{"There are still duplicate postal codes in the filtered conversion table."}\NormalTok{)}
\NormalTok{\}}

\NormalTok{matched }\OtherTok{\textless{}{-}}\NormalTok{ data }\SpecialCharTok{\%\textgreater{}\%} \FunctionTok{select}\NormalTok{(}\FunctionTok{c}\NormalTok{(}\StringTok{"ID"}\NormalTok{, }\StringTok{"PC"}\NormalTok{)) }\SpecialCharTok{\%\textgreater{}\%}
    \FunctionTok{mutate}\NormalTok{(}\AttributeTok{PC =} \FunctionTok{as.integer}\NormalTok{(}\FunctionTok{as.character}\NormalTok{(PC))) }\SpecialCharTok{\%\textgreater{}\%}
    \FunctionTok{left\_join}\NormalTok{(}\FunctionTok{select}\NormalTok{(conv\_matching, }\FunctionTok{c}\NormalTok{(}\StringTok{"Postal.code"}\NormalTok{,}\StringTok{"Refnis.code"}\NormalTok{)), }\AttributeTok{by =} \FunctionTok{c}\NormalTok{(}\StringTok{"PC"} \OtherTok{=} \StringTok{"Postal.code"}\NormalTok{)) }\SpecialCharTok{\%\textgreater{}\%}
    \FunctionTok{left\_join}\NormalTok{(}\FunctionTok{select}\NormalTok{(refnis, }\FunctionTok{c}\NormalTok{(}\StringTok{"NIS\_MUN"}\NormalTok{, }\StringTok{"MUNICIPALITY"}\NormalTok{, }\StringTok{"REGION"}\NormalTok{)), }\AttributeTok{by =} \FunctionTok{c}\NormalTok{(}\StringTok{"Refnis.code"} \OtherTok{=} \StringTok{"NIS\_MUN"}\NormalTok{)) }\SpecialCharTok{\%\textgreater{}\%}
    \FunctionTok{left\_join}\NormalTok{(}\FunctionTok{select}\NormalTok{(pop, }\FunctionTok{c}\NormalTok{(}\StringTok{"NIS.code"}\NormalTok{, }\StringTok{"Totaal"}\NormalTok{)), }\AttributeTok{by =} \FunctionTok{c}\NormalTok{(}\StringTok{"Refnis.code"} \OtherTok{=} \StringTok{"NIS.code"}\NormalTok{)) }\SpecialCharTok{\%\textgreater{}\%}
    \FunctionTok{rename}\NormalTok{(}
        \AttributeTok{REFNIS =}\NormalTok{ Refnis.code,}
        \AttributeTok{MUNICIPALITY =}\NormalTok{ MUNICIPALITY,}
        \AttributeTok{REGION =}\NormalTok{ REGION,}
        \AttributeTok{POPULATION =}\NormalTok{ Totaal}
\NormalTok{    )}

\ControlFlowTok{if}\NormalTok{(}\FunctionTok{any}\NormalTok{(}\FunctionTok{duplicated}\NormalTok{(matched}\SpecialCharTok{$}\NormalTok{ID))) }\FunctionTok{stop}\NormalTok{(}\StringTok{"There are duplicate IDs in the matched dataset."}\NormalTok{)}
\ControlFlowTok{if}\NormalTok{(}\FunctionTok{nrow}\NormalTok{(matched) }\SpecialCharTok{!=} \FunctionTok{nrow}\NormalTok{(data)) }\FunctionTok{stop}\NormalTok{(}\StringTok{"The number of matched rows does not match the original data."}\NormalTok{)}

\NormalTok{data\_enriched }\OtherTok{\textless{}{-}}\NormalTok{ data }\SpecialCharTok{\%\textgreater{}\%}
    \FunctionTok{left\_join}\NormalTok{(}\FunctionTok{select}\NormalTok{(matched, }\FunctionTok{c}\NormalTok{(}\StringTok{"ID"}\NormalTok{, }\StringTok{"REFNIS"}\NormalTok{, }\StringTok{"MUNICIPALITY"}\NormalTok{, }\StringTok{"REGION"}\NormalTok{, }\StringTok{"POPULATION"}\NormalTok{)), }\AttributeTok{by =} \StringTok{"ID"}\NormalTok{) }
\NormalTok{data\_enriched }\OtherTok{\textless{}{-}}\NormalTok{ data\_enriched }\SpecialCharTok{\%\textgreater{}\%}
    \FunctionTok{mutate}\NormalTok{(}
        \AttributeTok{REFNIS =} \FunctionTok{factor}\NormalTok{(REFNIS) }\SpecialCharTok{\%\textgreater{}\%} \FunctionTok{fct\_infreq}\NormalTok{(),}
        \AttributeTok{MUNICIPALITY =} \FunctionTok{factor}\NormalTok{(MUNICIPALITY) }\SpecialCharTok{\%\textgreater{}\%} \FunctionTok{fct\_infreq}\NormalTok{(),}
        \AttributeTok{REGION =} \FunctionTok{factor}\NormalTok{(REGION) }\SpecialCharTok{\%\textgreater{}\%} \FunctionTok{fct\_infreq}\NormalTok{(),}
        \AttributeTok{POPULATION =} \FunctionTok{as.numeric}\NormalTok{(}\FunctionTok{gsub}\NormalTok{(}\StringTok{"}\SpecialCharTok{\textbackslash{}\textbackslash{}}\StringTok{."}\NormalTok{, }\StringTok{""}\NormalTok{, POPULATION)),}
        \AttributeTok{POPULATION =} \FunctionTok{ifelse}\NormalTok{(}\FunctionTok{is.na}\NormalTok{(POPULATION), }\FunctionTok{mean}\NormalTok{(POPULATION, }\AttributeTok{na.rm =} \ConstantTok{TRUE}\NormalTok{), POPULATION)}
\NormalTok{    )}

\ControlFlowTok{if}\NormalTok{(}\FunctionTok{nrow}\NormalTok{(data) }\SpecialCharTok{!=} \FunctionTok{nrow}\NormalTok{(data\_enriched)) }\FunctionTok{stop}\NormalTok{(}\StringTok{"The number of rows in the enriched data does not match the original data."}\NormalTok{)}
\ControlFlowTok{if}\NormalTok{(}\FunctionTok{any}\NormalTok{(}\FunctionTok{is.na}\NormalTok{(data\_enriched}\SpecialCharTok{$}\NormalTok{ID))) }\FunctionTok{stop}\NormalTok{(}\StringTok{"There are missing IDs in the enriched data."}\NormalTok{)}

\FunctionTok{summary}\NormalTok{(data\_enriched }\SpecialCharTok{\%\textgreater{}\%} \FunctionTok{select}\NormalTok{(}\FunctionTok{c}\NormalTok{(}\StringTok{"PC"}\NormalTok{, }\StringTok{"REFNIS"}\NormalTok{, }\StringTok{"MUNICIPALITY"}\NormalTok{, }\StringTok{"REGION"}\NormalTok{, }\StringTok{"POPULATION"}\NormalTok{))) }\SpecialCharTok{\%\textgreater{}\%}
\FunctionTok{kable}\NormalTok{(}\AttributeTok{booktabs =} \ConstantTok{TRUE}\NormalTok{, }\AttributeTok{width =} \DecValTok{100}\NormalTok{, }\AttributeTok{caption =} \StringTok{"Summary of Enriched Data"}\NormalTok{)}
\end{Highlighting}
\end{Shaded}

\begin{longtable}[]{@{}
  >{\raggedright\arraybackslash}p{(\linewidth - 10\tabcolsep) * \real{0.0294}}
  >{\raggedright\arraybackslash}p{(\linewidth - 10\tabcolsep) * \real{0.1471}}
  >{\raggedright\arraybackslash}p{(\linewidth - 10\tabcolsep) * \real{0.1471}}
  >{\raggedright\arraybackslash}p{(\linewidth - 10\tabcolsep) * \real{0.1863}}
  >{\raggedright\arraybackslash}p{(\linewidth - 10\tabcolsep) * \real{0.3431}}
  >{\raggedright\arraybackslash}p{(\linewidth - 10\tabcolsep) * \real{0.1471}}@{}}
\caption{Summary of Enriched Data}\tabularnewline
\toprule\noalign{}
\begin{minipage}[b]{\linewidth}\raggedright
\end{minipage} & \begin{minipage}[b]{\linewidth}\raggedright
PC
\end{minipage} & \begin{minipage}[b]{\linewidth}\raggedright
REFNIS
\end{minipage} & \begin{minipage}[b]{\linewidth}\raggedright
MUNICIPALITY
\end{minipage} & \begin{minipage}[b]{\linewidth}\raggedright
REGION
\end{minipage} & \begin{minipage}[b]{\linewidth}\raggedright
POPULATION
\end{minipage} \\
\midrule\noalign{}
\endfirsthead
\toprule\noalign{}
\begin{minipage}[b]{\linewidth}\raggedright
\end{minipage} & \begin{minipage}[b]{\linewidth}\raggedright
PC
\end{minipage} & \begin{minipage}[b]{\linewidth}\raggedright
REFNIS
\end{minipage} & \begin{minipage}[b]{\linewidth}\raggedright
MUNICIPALITY
\end{minipage} & \begin{minipage}[b]{\linewidth}\raggedright
REGION
\end{minipage} & \begin{minipage}[b]{\linewidth}\raggedright
POPULATION
\end{minipage} \\
\midrule\noalign{}
\endhead
\bottomrule\noalign{}
\endlastfoot
& 6000 : 3992 & 52011 : 3992 & Charleroi : 3992 & RÉGION WALLONNE :62973
& Min. : 995 \\
& 2000 : 3716 & 11002 : 3837 & Anvers : 3837 & RÉGION FLAMANDE :59750 &
1st Qu.: 13222 \\
& 9000 : 2770 & 44021 : 2770 & Gand : 2770 & RÉGION DE
BRUXELLES-CAPITALE: 7862 & Median : 24586 \\
& 4000 : 2115 & 62063 : 2115 & Liège : 2115 & NA & Mean : 56488 \\
& 7100 : 1806 & 58001 : 1806 & La Louvière: 1806 & NA & 3rd Qu.:
56488 \\
& 7000 : 1649 & 53053 : 1649 & Mons : 1649 & NA & Max. :453030 \\
& (Other):114537 & (Other):114416 & (Other) :114416 & NA & NA \\
\end{longtable}

Now that we have the enriched data, we can bin the population density
into categories and assess the significance of municipality, region, and
population density. Note that we have 0 missing values in the population
column. We'll use mean imputation to fill these missing values before
binning.

\paragraph{Population Density}\label{population-density}

\begin{Shaded}
\begin{Highlighting}[]
\NormalTok{baseline\_pop }\OtherTok{\textless{}{-}} \FunctionTok{qbinning}\NormalTok{(data\_enriched, }\StringTok{"POPULATION"}\NormalTok{, }\StringTok{"NCLAIMS"}\NormalTok{, }\AttributeTok{max\_bins =} \DecValTok{30}\NormalTok{)}
\NormalTok{result\_pop }\OtherTok{\textless{}{-}} \FunctionTok{iv\_elbow}\NormalTok{(data\_enriched, }\StringTok{"POPULATION"}\NormalTok{, }\StringTok{"NCLAIMS"}\NormalTok{, baseline\_pop)}
\NormalTok{result\_named\_list[[}\StringTok{"POPULATION"}\NormalTok{]] }\OtherTok{\textless{}{-}}\NormalTok{ result\_pop}
\FunctionTok{plot\_iv\_elbow}\NormalTok{(result\_pop)}
\end{Highlighting}
\end{Shaded}

\begin{center}\includegraphics{project_MD_files/figure-latex/binning_population-1} \end{center}

Based on the elbow plot, we can look at the WoE and IV for 5, 6, 12 and
20 bins:

\begin{Shaded}
\begin{Highlighting}[]
\NormalTok{plots\_pop }\OtherTok{\textless{}{-}} \FunctionTok{list}\NormalTok{()}
\ControlFlowTok{for}\NormalTok{ (i }\ControlFlowTok{in} \FunctionTok{c}\NormalTok{(}\DecValTok{5}\NormalTok{, }\DecValTok{6}\NormalTok{, }\DecValTok{12}\NormalTok{, }\DecValTok{20}\NormalTok{)) \{}
\NormalTok{    temp\_bins\_pop }\OtherTok{\textless{}{-}} \FunctionTok{rebuild\_bins}\NormalTok{(result\_pop, i)}
\NormalTok{    woe\_temp\_pop }\OtherTok{\textless{}{-}} \FunctionTok{woe\_calc}\NormalTok{(data\_enriched, }\AttributeTok{x =} \StringTok{"POPULATION"}\NormalTok{, }\AttributeTok{y =} \StringTok{"NCLAIMS"}\NormalTok{, }\AttributeTok{bins =}\NormalTok{ temp\_bins\_pop)}
\NormalTok{    plots\_pop[[i]] }\OtherTok{\textless{}{-}} \FunctionTok{plot\_woe\_iv}\NormalTok{(woe\_temp\_pop, }\AttributeTok{interactive =} \ConstantTok{FALSE}\NormalTok{)}
\NormalTok{\}}
\NormalTok{(plots\_pop[[}\DecValTok{5}\NormalTok{]] }\SpecialCharTok{|}\NormalTok{ plots\_pop[[}\DecValTok{6}\NormalTok{]]) }\SpecialCharTok{/}\NormalTok{ (plots\_pop[[}\DecValTok{12}\NormalTok{]] }\SpecialCharTok{|}\NormalTok{ plots\_pop[[}\DecValTok{20}\NormalTok{]])}
\end{Highlighting}
\end{Shaded}

\begin{center}\includegraphics{project_MD_files/figure-latex/woe_population-1} \end{center}

We see a linear diminution of the WoE as the population increases,
indicating that higher population areas tend to have higher claim
frequencies as we expected. The linear trend seems sufficiently captured
with 5 bins. We'll create the new binned variable with 5 levels.

\begin{Shaded}
\begin{Highlighting}[]
\NormalTok{data\_enriched}\SpecialCharTok{$}\NormalTok{POPULATION\_binned }\OtherTok{\textless{}{-}} \FunctionTok{cut}\NormalTok{(}
\NormalTok{    data\_enriched}\SpecialCharTok{$}\NormalTok{POPULATION,}
    \AttributeTok{breaks =} \FunctionTok{rebuild\_bins}\NormalTok{(result\_pop, }\DecValTok{5}\NormalTok{),}
    \AttributeTok{include.lowest =} \ConstantTok{TRUE}\NormalTok{,}
    \AttributeTok{ordered\_result =} \ConstantTok{TRUE}
\NormalTok{)}
\NormalTok{n\_bins\_named\_list[[}\StringTok{"POPULATION"}\NormalTok{]] }\OtherTok{\textless{}{-}} \DecValTok{5}
\NormalTok{m0 }\OtherTok{\textless{}{-}} \FunctionTok{glm}\NormalTok{(NCLAIMS }\SpecialCharTok{\textasciitilde{}} \FunctionTok{offset}\NormalTok{(}\FunctionTok{log}\NormalTok{(EXP)), }\AttributeTok{data =}\NormalTok{ data\_enriched, }\AttributeTok{family =} \FunctionTok{poisson}\NormalTok{(}\AttributeTok{link =} \StringTok{"log"}\NormalTok{))}
\NormalTok{m1 }\OtherTok{\textless{}{-}} \FunctionTok{glm}\NormalTok{(NCLAIMS }\SpecialCharTok{\textasciitilde{}}\NormalTok{ POPULATION\_binned }\SpecialCharTok{+} \FunctionTok{offset}\NormalTok{(}\FunctionTok{log}\NormalTok{(EXP)), }\AttributeTok{data =}\NormalTok{ data\_enriched, }\AttributeTok{family =} \FunctionTok{poisson}\NormalTok{(}\AttributeTok{link =} \StringTok{"log"}\NormalTok{))}
\NormalTok{chisq\_result }\OtherTok{\textless{}{-}} \FunctionTok{anova}\NormalTok{(m0, m1, }\AttributeTok{test =} \StringTok{"Chisq"}\NormalTok{)}
\FunctionTok{cat}\NormalTok{(}\FunctionTok{sprintf}\NormalTok{(}\StringTok{"Is the effect of POPULATION\_binned significant at 95\%\% confidence? \%s"}\NormalTok{, }\FunctionTok{ifelse}\NormalTok{(chisq\_result}\SpecialCharTok{$}\StringTok{\textasciigrave{}}\AttributeTok{Pr(\textgreater{}Chi)}\StringTok{\textasciigrave{}}\NormalTok{[}\DecValTok{2}\NormalTok{] }\SpecialCharTok{\textless{}} \FloatTok{0.05}\NormalTok{, }\StringTok{"Yes"}\NormalTok{, }\StringTok{"No"}\NormalTok{)))}
\end{Highlighting}
\end{Shaded}

\begin{verbatim}
Is the effect of POPULATION_binned significant at 95% confidence? Yes
\end{verbatim}

\begin{Shaded}
\begin{Highlighting}[]
\FunctionTok{plot\_claim\_freq}\NormalTok{(data\_enriched, }\StringTok{"POPULATION\_binned"}\NormalTok{)}
\end{Highlighting}
\end{Shaded}

\begin{center}\includegraphics{project_MD_files/figure-latex/binned_population-1} \end{center}

The variable is significant at 95\% confidence, however we see that the
3 largest population bins have very similar claim frequencies, we can
try merging them to get have a simple three-level binning (Low, Medium,
High Density).

\begin{Shaded}
\begin{Highlighting}[]
\NormalTok{initial\_bins\_merged }\OtherTok{\textless{}{-}}\NormalTok{ data\_enriched}\SpecialCharTok{$}\NormalTok{POPULATION\_binned}

\NormalTok{three\_bins\_merged }\OtherTok{\textless{}{-}} \FunctionTok{fct\_collapse}\NormalTok{(}
\NormalTok{    data\_enriched}\SpecialCharTok{$}\NormalTok{POPULATION\_binned,}
    \StringTok{"Low Density"} \OtherTok{=} \FunctionTok{c}\NormalTok{(}\StringTok{"[995,7.68e+03]"}\NormalTok{),}
    \StringTok{"Medium Density"} \OtherTok{=} \FunctionTok{c}\NormalTok{(}\StringTok{"(7.68e+03,4.3e+04]"}\NormalTok{),}
    \StringTok{"High Density"} \OtherTok{=} \FunctionTok{c}\NormalTok{(}\StringTok{"(4.3e+04,1.34e+05]"}\NormalTok{, }\StringTok{"(1.34e+05,2.05e+05]"}\NormalTok{, }\StringTok{"(2.05e+05,4.53e+05]"}\NormalTok{)}
\NormalTok{)}

\NormalTok{m1 }\OtherTok{\textless{}{-}} \FunctionTok{glm}\NormalTok{(NCLAIMS }\SpecialCharTok{\textasciitilde{}} \FunctionTok{offset}\NormalTok{(}\FunctionTok{log}\NormalTok{(EXP)) }\SpecialCharTok{+}\NormalTok{ initial\_bins\_merged, }\AttributeTok{data =}\NormalTok{ data\_enriched, }\AttributeTok{family =} \FunctionTok{poisson}\NormalTok{(}\AttributeTok{link =} \StringTok{"log"}\NormalTok{))}
\NormalTok{m0 }\OtherTok{\textless{}{-}} \FunctionTok{glm}\NormalTok{(NCLAIMS }\SpecialCharTok{\textasciitilde{}} \FunctionTok{offset}\NormalTok{(}\FunctionTok{log}\NormalTok{(EXP)) }\SpecialCharTok{+}\NormalTok{ three\_bins\_merged, }\AttributeTok{data =}\NormalTok{ data\_enriched, }\AttributeTok{family =} \FunctionTok{poisson}\NormalTok{(}\AttributeTok{link =} \StringTok{"log"}\NormalTok{))}
\NormalTok{test\_pop }\OtherTok{\textless{}{-}} \FunctionTok{anova}\NormalTok{(m0, m1, }\AttributeTok{test =} \StringTok{"Chisq"}\NormalTok{)}
\FunctionTok{cat}\NormalTok{(}\FunctionTok{sprintf}\NormalTok{(}\StringTok{"Are we loosing significant information by merging the population bins? \%s (p{-}value: \%.4f)"}\NormalTok{, }\FunctionTok{ifelse}\NormalTok{(test\_pop}\SpecialCharTok{$}\StringTok{\textasciigrave{}}\AttributeTok{Pr(\textgreater{}Chi)}\StringTok{\textasciigrave{}}\NormalTok{[}\DecValTok{2}\NormalTok{] }\SpecialCharTok{\textless{}} \FloatTok{0.05}\NormalTok{, }\StringTok{"Yes"}\NormalTok{, }\StringTok{"No"}\NormalTok{), test\_pop}\SpecialCharTok{$}\StringTok{\textasciigrave{}}\AttributeTok{Pr(\textgreater{}Chi)}\StringTok{\textasciigrave{}}\NormalTok{[}\DecValTok{2}\NormalTok{]))}
\end{Highlighting}
\end{Shaded}

\begin{verbatim}
Are we loosing significant information by merging the population bins? Yes (p-value: 0.0020)
\end{verbatim}

\begin{Shaded}
\begin{Highlighting}[]
\NormalTok{four\_bins\_merged }\OtherTok{\textless{}{-}} \FunctionTok{fct\_collapse}\NormalTok{(}
\NormalTok{    data\_enriched}\SpecialCharTok{$}\NormalTok{POPULATION\_binned,}
    \StringTok{"(1.34e+05,4.53e+05]"} \OtherTok{=} \FunctionTok{c}\NormalTok{(}\StringTok{"(1.34e+05,2.05e+05]"}\NormalTok{, }\StringTok{"(2.05e+05,4.53e+05]"}\NormalTok{)}
\NormalTok{)}

\NormalTok{m0 }\OtherTok{\textless{}{-}} \FunctionTok{glm}\NormalTok{(NCLAIMS }\SpecialCharTok{\textasciitilde{}} \FunctionTok{offset}\NormalTok{(}\FunctionTok{log}\NormalTok{(EXP)) }\SpecialCharTok{+}\NormalTok{ four\_bins\_merged, }\AttributeTok{data =}\NormalTok{ data\_enriched, }\AttributeTok{family =} \FunctionTok{poisson}\NormalTok{(}\AttributeTok{link =} \StringTok{"log"}\NormalTok{))}
\NormalTok{test\_pop\_four }\OtherTok{\textless{}{-}} \FunctionTok{anova}\NormalTok{(m0, m1, }\AttributeTok{test =} \StringTok{"Chisq"}\NormalTok{)}
\FunctionTok{cat}\NormalTok{(}\FunctionTok{sprintf}\NormalTok{(}\StringTok{"Are we loosing significant information by merging the last two population bins? \%s (p{-}value: \%.4f)"}\NormalTok{, }\FunctionTok{ifelse}\NormalTok{(test\_pop\_four}\SpecialCharTok{$}\StringTok{\textasciigrave{}}\AttributeTok{Pr(\textgreater{}Chi)}\StringTok{\textasciigrave{}}\NormalTok{[}\DecValTok{2}\NormalTok{] }\SpecialCharTok{\textless{}} \FloatTok{0.05}\NormalTok{, }\StringTok{"Yes"}\NormalTok{, }\StringTok{"No"}\NormalTok{), test\_pop\_four}\SpecialCharTok{$}\StringTok{\textasciigrave{}}\AttributeTok{Pr(\textgreater{}Chi)}\StringTok{\textasciigrave{}}\NormalTok{[}\DecValTok{2}\NormalTok{]))}
\end{Highlighting}
\end{Shaded}

\begin{verbatim}
Are we loosing significant information by merging the last two population bins? No (p-value: 0.5786)
\end{verbatim}

Based on the result of the anova test, it seems that merging the bins
into three levels does not significantly increase the residual deviance.
We'll therefore adopt the three-level binning for the population
variable which is a lot easier to interpret and work with.

\begin{Shaded}
\begin{Highlighting}[]
\NormalTok{data\_enriched}\SpecialCharTok{$}\NormalTok{POPULATION\_binned }\OtherTok{\textless{}{-}}\NormalTok{ four\_bins\_merged}
\FunctionTok{plot\_claim\_freq}\NormalTok{(data\_enriched, }\StringTok{"POPULATION\_binned"}\NormalTok{)}
\end{Highlighting}
\end{Shaded}

\begin{center}\includegraphics{project_MD_files/figure-latex/three_bins_plot-1} \end{center}

We get a linear trend in claim frequency with respect to population
density bins (ordinal variable).

\paragraph{Regions}\label{regions}

We will check the significance and claim frequency across regions

\begin{Shaded}
\begin{Highlighting}[]
\NormalTok{m0 }\OtherTok{\textless{}{-}} \FunctionTok{glm}\NormalTok{(NCLAIMS }\SpecialCharTok{\textasciitilde{}} \FunctionTok{offset}\NormalTok{(}\FunctionTok{log}\NormalTok{(EXP)), }\AttributeTok{data =}\NormalTok{ data\_enriched, }\AttributeTok{family =} \FunctionTok{poisson}\NormalTok{(}\AttributeTok{link =} \StringTok{"log"}\NormalTok{))}
\NormalTok{m1 }\OtherTok{\textless{}{-}} \FunctionTok{glm}\NormalTok{(NCLAIMS }\SpecialCharTok{\textasciitilde{}} \FunctionTok{offset}\NormalTok{(}\FunctionTok{log}\NormalTok{(EXP)) }\SpecialCharTok{+}\NormalTok{ REGION, }\AttributeTok{data =}\NormalTok{ data\_enriched, }\AttributeTok{family =} \FunctionTok{poisson}\NormalTok{(}\AttributeTok{link =} \StringTok{"log"}\NormalTok{))}
\NormalTok{    test\_result }\OtherTok{\textless{}{-}} \FunctionTok{anova}\NormalTok{(m0, m1, }\AttributeTok{test =} \StringTok{"Chisq"}\NormalTok{)}
    \FunctionTok{cat}\NormalTok{(}\FunctionTok{sprintf}\NormalTok{(}\StringTok{"Is variable \%s significant? \%s}\SpecialCharTok{\textbackslash{}n}\StringTok{"}\NormalTok{, }\StringTok{"REGION"}\NormalTok{, }\FunctionTok{ifelse}\NormalTok{(test\_result}\SpecialCharTok{$}\StringTok{\textasciigrave{}}\AttributeTok{Pr(\textgreater{}Chi)}\StringTok{\textasciigrave{}}\NormalTok{[}\DecValTok{2}\NormalTok{] }\SpecialCharTok{\textless{}}\NormalTok{ test\_level, }\StringTok{"Yes"}\NormalTok{, }\StringTok{"No"}\NormalTok{)))}
\end{Highlighting}
\end{Shaded}

\begin{verbatim}
Is variable REGION significant? Yes
\end{verbatim}

\begin{Shaded}
\begin{Highlighting}[]
\FunctionTok{plot\_claim\_freq}\NormalTok{(data\_enriched, }\StringTok{"REGION"}\NormalTok{)}
\end{Highlighting}
\end{Shaded}

\begin{center}\includegraphics{project_MD_files/figure-latex/municipality_region_significance-1} \end{center}

It is clear that the Brussels region has a distinct claim frequency
pattern compared to the other regions. We'll test if it makes sense to
keep the Walloon and Flemish regions separate as well or not.

\begin{Shaded}
\begin{Highlighting}[]
\NormalTok{merged\_regions }\OtherTok{\textless{}{-}} \FunctionTok{fct\_collapse}\NormalTok{(}
\NormalTok{    data\_enriched}\SpecialCharTok{$}\NormalTok{REGION,}
    \StringTok{"Brussels"} \OtherTok{=} \FunctionTok{c}\NormalTok{(}\StringTok{"RÉGION DE BRUXELLES{-}CAPITALE"}\NormalTok{),}
    \StringTok{"Others"} \OtherTok{=} \FunctionTok{c}\NormalTok{(}\StringTok{"RÉGION WALLONNE"}\NormalTok{, }\StringTok{"RÉGION FLAMANDE"}\NormalTok{),}
\NormalTok{)}
\NormalTok{data\_enriched}\SpecialCharTok{$}\NormalTok{REGION\_merged }\OtherTok{\textless{}{-}}\NormalTok{ merged\_regions}
\NormalTok{m0 }\OtherTok{\textless{}{-}} \FunctionTok{glm}\NormalTok{(NCLAIMS }\SpecialCharTok{\textasciitilde{}} \FunctionTok{offset}\NormalTok{(}\FunctionTok{log}\NormalTok{(EXP)) }\SpecialCharTok{+}\NormalTok{ REGION, }\AttributeTok{data =}\NormalTok{ data\_enriched, }\AttributeTok{family =} \FunctionTok{poisson}\NormalTok{(}\AttributeTok{link =} \StringTok{"log"}\NormalTok{))}
\NormalTok{m1 }\OtherTok{\textless{}{-}} \FunctionTok{glm}\NormalTok{(NCLAIMS }\SpecialCharTok{\textasciitilde{}} \FunctionTok{offset}\NormalTok{(}\FunctionTok{log}\NormalTok{(EXP)) }\SpecialCharTok{+}\NormalTok{ REGION\_merged, }\AttributeTok{data =}\NormalTok{ data\_enriched, }\AttributeTok{family =} \FunctionTok{poisson}\NormalTok{(}\AttributeTok{link =} \StringTok{"log"}\NormalTok{))}
\NormalTok{test\_result }\OtherTok{\textless{}{-}} \FunctionTok{anova}\NormalTok{(m0, m1, }\AttributeTok{test =} \StringTok{"Chisq"}\NormalTok{)}
\FunctionTok{cat}\NormalTok{(}\FunctionTok{sprintf}\NormalTok{(}\StringTok{"Was it justified to merge the REGION bins? \%s"}\NormalTok{, }\FunctionTok{ifelse}\NormalTok{(test\_result}\SpecialCharTok{$}\StringTok{\textasciigrave{}}\AttributeTok{Pr(\textgreater{}Chi)}\StringTok{\textasciigrave{}}\NormalTok{[}\DecValTok{2}\NormalTok{] }\SpecialCharTok{\textless{}}\NormalTok{ test\_level, }\StringTok{"No"}\NormalTok{, }\StringTok{"Yes"}\NormalTok{)))}
\end{Highlighting}
\end{Shaded}

\begin{verbatim}
Was it justified to merge the REGION bins? Yes
\end{verbatim}

\begin{Shaded}
\begin{Highlighting}[]
\FunctionTok{plot\_claim\_freq}\NormalTok{(data\_enriched, }\StringTok{"REGION\_merged"}\NormalTok{)}
\end{Highlighting}
\end{Shaded}

\begin{center}\includegraphics{project_MD_files/figure-latex/region_significance-1} \end{center}

We see that despite a much lower exposure in Brussels, the claim
frequency is significantly higher from the other regions. This confirms
our hypothesis that urban areas, such as Brussels, tend to have a higher
claim frequency compared to less densely populated regions. We can
quickly check our second hypothesis that despite higher frequencies, the
severity is lower in urban areas like Brussels.

\begin{Shaded}
\begin{Highlighting}[]
\NormalTok{urb\_rur\_comp }\OtherTok{\textless{}{-}}\NormalTok{ data\_enriched }\SpecialCharTok{\%\textgreater{}\%}
    \FunctionTok{group\_by}\NormalTok{(REGION\_merged) }\SpecialCharTok{\%\textgreater{}\%}
    \FunctionTok{summarise}\NormalTok{(}
        \AttributeTok{total\_claims =} \FunctionTok{sum}\NormalTok{(NCLAIMS),}
        \AttributeTok{total\_exposure =} \FunctionTok{sum}\NormalTok{(EXP),}
        \AttributeTok{total\_claims =} \FunctionTok{sum}\NormalTok{(NCLAIMS),}
        \AttributeTok{total\_claim\_amount =} \FunctionTok{sum}\NormalTok{(AMOUNT),}
        \AttributeTok{frequency =}\NormalTok{ total\_claims }\SpecialCharTok{/}\NormalTok{ total\_exposure,}
        \AttributeTok{severity =}\NormalTok{ total\_claim\_amount }\SpecialCharTok{/}\NormalTok{ total\_claims,}
        \AttributeTok{expected\_claims =}\NormalTok{ frequency }\SpecialCharTok{*}\NormalTok{ severity}
\NormalTok{    ) }\SpecialCharTok{\%\textgreater{}\%}
    \FunctionTok{arrange}\NormalTok{(}\FunctionTok{desc}\NormalTok{(frequency)) }
\NormalTok{urb\_rur\_comp }\SpecialCharTok{\%\textgreater{}\%}
    \FunctionTok{kable}\NormalTok{(}
        \AttributeTok{caption =} \StringTok{"Expected claims by merged region}\SpecialCharTok{\textbackslash{}n}\StringTok{Split by frequency and severity"}\NormalTok{,}
        \AttributeTok{align =} \StringTok{"l"}\NormalTok{, }\AttributeTok{booktabs =} \ConstantTok{TRUE}
\NormalTok{    )}
\end{Highlighting}
\end{Shaded}

\begin{longtable}[]{@{}
  >{\raggedright\arraybackslash}p{(\linewidth - 12\tabcolsep) * \real{0.1458}}
  >{\raggedright\arraybackslash}p{(\linewidth - 12\tabcolsep) * \real{0.1354}}
  >{\raggedright\arraybackslash}p{(\linewidth - 12\tabcolsep) * \real{0.1562}}
  >{\raggedright\arraybackslash}p{(\linewidth - 12\tabcolsep) * \real{0.1979}}
  >{\raggedright\arraybackslash}p{(\linewidth - 12\tabcolsep) * \real{0.1042}}
  >{\raggedright\arraybackslash}p{(\linewidth - 12\tabcolsep) * \real{0.0938}}
  >{\raggedright\arraybackslash}p{(\linewidth - 12\tabcolsep) * \real{0.1667}}@{}}
\caption{Expected claims by merged region Split by frequency and
severity}\tabularnewline
\toprule\noalign{}
\begin{minipage}[b]{\linewidth}\raggedright
REGION\_merged
\end{minipage} & \begin{minipage}[b]{\linewidth}\raggedright
total\_claims
\end{minipage} & \begin{minipage}[b]{\linewidth}\raggedright
total\_exposure
\end{minipage} & \begin{minipage}[b]{\linewidth}\raggedright
total\_claim\_amount
\end{minipage} & \begin{minipage}[b]{\linewidth}\raggedright
frequency
\end{minipage} & \begin{minipage}[b]{\linewidth}\raggedright
severity
\end{minipage} & \begin{minipage}[b]{\linewidth}\raggedright
expected\_claims
\end{minipage} \\
\midrule\noalign{}
\endfirsthead
\toprule\noalign{}
\begin{minipage}[b]{\linewidth}\raggedright
REGION\_merged
\end{minipage} & \begin{minipage}[b]{\linewidth}\raggedright
total\_claims
\end{minipage} & \begin{minipage}[b]{\linewidth}\raggedright
total\_exposure
\end{minipage} & \begin{minipage}[b]{\linewidth}\raggedright
total\_claim\_amount
\end{minipage} & \begin{minipage}[b]{\linewidth}\raggedright
frequency
\end{minipage} & \begin{minipage}[b]{\linewidth}\raggedright
severity
\end{minipage} & \begin{minipage}[b]{\linewidth}\raggedright
expected\_claims
\end{minipage} \\
\midrule\noalign{}
\endhead
\bottomrule\noalign{}
\endlastfoot
Brussels & 1477 & 6482.844 & 1776783 & 0.2278321 & 1202.968 &
274.0746 \\
Others & 14632 & 109723.584 & 25498069 & 0.1333533 & 1742.624 &
232.3846 \\
\end{longtable}

With this quick descriptive analysis, we can confirm that in our sample:

\begin{itemize}
\tightlist
\item
  The claim frequency is higher in urban areas like Brussels compared to
  other regions (factor 1.7084852)
\item
  The claim severity is lower in urban areas like Brussels compared to
  other regions (factor 0.6903198)
\item
  The frequency effect outweighs the severity effect in determining the
  expected claims between urban and other regions (factor 1.1794012)
\end{itemize}

We expect to see similar effects for other granularities, such as
municipalities, where urban areas may exhibit higher claim frequencies
but lower severities compared to less densely populated areas.

\paragraph{Municipalities}\label{municipalities}

There are far too many municipalities to use it as a simple categorical
variable in our models. We can thus focus on the n largest
municipalities by exposure and group the remaining ones into an
``Other'' category.

\begin{Shaded}
\begin{Highlighting}[]
\NormalTok{data\_enriched }\SpecialCharTok{\%\textgreater{}\%}
    \FunctionTok{group\_by}\NormalTok{(MUNICIPALITY) }\SpecialCharTok{\%\textgreater{}\%}
    \FunctionTok{summarise}\NormalTok{(}
        \AttributeTok{total\_exposure =} \FunctionTok{sum}\NormalTok{(EXP)}
\NormalTok{    ) }\SpecialCharTok{\%\textgreater{}\%}
    \FunctionTok{arrange}\NormalTok{(}\FunctionTok{desc}\NormalTok{(total\_exposure)) }\SpecialCharTok{\%\textgreater{}\%} \FunctionTok{head}\NormalTok{(}\AttributeTok{n=}\DecValTok{50}\NormalTok{) }\SpecialCharTok{\%\textgreater{}\%}
    \FunctionTok{ggplot}\NormalTok{(}\FunctionTok{aes}\NormalTok{(}\AttributeTok{x =} \FunctionTok{reorder}\NormalTok{(MUNICIPALITY, }\SpecialCharTok{{-}}\NormalTok{total\_exposure), }\AttributeTok{y =}\NormalTok{ total\_exposure)) }\SpecialCharTok{+}
    \FunctionTok{geom\_bar}\NormalTok{(}\AttributeTok{stat =} \StringTok{"identity"}\NormalTok{) }\SpecialCharTok{+}
    \FunctionTok{labs}\NormalTok{(}\AttributeTok{x =} \StringTok{"Municipality"}\NormalTok{, }\AttributeTok{y =} \StringTok{"Total Exposure"}\NormalTok{) }\SpecialCharTok{+}
    \FunctionTok{theme}\NormalTok{(}\AttributeTok{axis.text.x =} \FunctionTok{element\_text}\NormalTok{(}\AttributeTok{angle =} \DecValTok{45}\NormalTok{, }\AttributeTok{hjust =} \DecValTok{1}\NormalTok{))}
\end{Highlighting}
\end{Shaded}

\begin{center}\includegraphics{project_MD_files/figure-latex/municipalities_aggregation-1} \end{center}

We can focus on the 6 largest municipalities by exposure before the
exposure distribution flattens out.

\begin{Shaded}
\begin{Highlighting}[]
\NormalTok{exp\_breakpoint }\OtherTok{\textless{}{-}}\NormalTok{ data\_enriched }\SpecialCharTok{\%\textgreater{}\%}
    \FunctionTok{group\_by}\NormalTok{(MUNICIPALITY) }\SpecialCharTok{\%\textgreater{}\%}
    \FunctionTok{summarise}\NormalTok{(}
        \AttributeTok{total\_exposure =} \FunctionTok{sum}\NormalTok{(EXP)}
\NormalTok{    ) }\SpecialCharTok{\%\textgreater{}\%}
    \FunctionTok{arrange}\NormalTok{(}\FunctionTok{desc}\NormalTok{(total\_exposure)) }\SpecialCharTok{\%\textgreater{}\%}
    \FunctionTok{head}\NormalTok{(}\AttributeTok{n=}\DecValTok{6}\NormalTok{) }
\NormalTok{data\_enriched}\SpecialCharTok{$}\NormalTok{MUNICIPALITY\_binned }\OtherTok{\textless{}{-}} \FunctionTok{ifelse}\NormalTok{(data\_enriched}\SpecialCharTok{$}\NormalTok{MUNICIPALITY }\SpecialCharTok{\%in\%}\NormalTok{ exp\_breakpoint}\SpecialCharTok{$}\NormalTok{MUNICIPALITY, }\FunctionTok{as.character}\NormalTok{(data\_enriched}\SpecialCharTok{$}\NormalTok{MUNICIPALITY), }\StringTok{"Other"}\NormalTok{)}
\NormalTok{data\_enriched}\SpecialCharTok{$}\NormalTok{MUNICIPALITY\_binned }\OtherTok{\textless{}{-}} \FunctionTok{factor}\NormalTok{(data\_enriched}\SpecialCharTok{$}\NormalTok{MUNICIPALITY\_binned) }\SpecialCharTok{\%\textgreater{}\%} \FunctionTok{fct\_infreq}\NormalTok{()}

\FunctionTok{plot\_claim\_freq}\NormalTok{(data\_enriched, }\StringTok{"MUNICIPALITY\_binned"}\NormalTok{)}
\end{Highlighting}
\end{Shaded}

\begin{center}\includegraphics{project_MD_files/figure-latex/top_municipalities-1} \end{center}

\begin{Shaded}
\begin{Highlighting}[]
\NormalTok{m0 }\OtherTok{\textless{}{-}} \FunctionTok{glm}\NormalTok{(NCLAIMS }\SpecialCharTok{\textasciitilde{}} \FunctionTok{offset}\NormalTok{(}\FunctionTok{log}\NormalTok{(EXP)), }\AttributeTok{data =}\NormalTok{ data\_enriched, }\AttributeTok{family =} \FunctionTok{poisson}\NormalTok{(}\AttributeTok{link =} \StringTok{"log"}\NormalTok{))}
\NormalTok{m1 }\OtherTok{\textless{}{-}} \FunctionTok{glm}\NormalTok{(NCLAIMS }\SpecialCharTok{\textasciitilde{}}\NormalTok{ MUNICIPALITY\_binned }\SpecialCharTok{+} \FunctionTok{offset}\NormalTok{(}\FunctionTok{log}\NormalTok{(EXP)), }\AttributeTok{data =}\NormalTok{ data\_enriched, }\AttributeTok{family =} \FunctionTok{poisson}\NormalTok{(}\AttributeTok{link =} \StringTok{"log"}\NormalTok{))}
\NormalTok{results }\OtherTok{\textless{}{-}} \FunctionTok{anova}\NormalTok{(m0, m1, }\AttributeTok{test =} \StringTok{"Chisq"}\NormalTok{)}
\FunctionTok{cat}\NormalTok{(}\FunctionTok{sprintf}\NormalTok{(}\StringTok{"Is variable \%s significant? \%s}\SpecialCharTok{\textbackslash{}n}\StringTok{"}\NormalTok{, }\StringTok{"MUNICIPALITY\_binned"}\NormalTok{, }\FunctionTok{ifelse}\NormalTok{(results}\SpecialCharTok{$}\StringTok{\textasciigrave{}}\AttributeTok{Pr(\textgreater{}Chi)}\StringTok{\textasciigrave{}}\NormalTok{[}\DecValTok{2}\NormalTok{] }\SpecialCharTok{\textless{}}\NormalTok{ test\_level, }\StringTok{"Yes"}\NormalTok{, }\StringTok{"No"}\NormalTok{)))}
\end{Highlighting}
\end{Shaded}

\begin{verbatim}
Is variable MUNICIPALITY_binned significant? Yes
\end{verbatim}

We now have a binned version of the MUNICIPALITY variable which is
significant and usable in our model. Of course, the exposure is limited
even in the bigger municipalities and that causes some variability in
the frequency estimates. However, for the sake of interpretation, we'll
keep the cities separated even though Charleroi and Antwerp are similar
in terms of frequency and so are Gand and Liège. We seem to capture a
similar phenomenon with cities having a higher claim frequency compared
to smaller municipalities even if the effect is less clear than for the
regions.

\begin{Shaded}
\begin{Highlighting}[]
\NormalTok{data\_enriched }\SpecialCharTok{\%\textgreater{}\%}
    \FunctionTok{group\_by}\NormalTok{(MUNICIPALITY\_binned) }\SpecialCharTok{\%\textgreater{}\%}
    \FunctionTok{summarise}\NormalTok{(}
        \AttributeTok{total\_claims =} \FunctionTok{sum}\NormalTok{(NCLAIMS),}
        \AttributeTok{total\_exposure =} \FunctionTok{sum}\NormalTok{(EXP),}
        \AttributeTok{total\_claim\_amount =} \FunctionTok{sum}\NormalTok{(AMOUNT),}
        \AttributeTok{frequency =}\NormalTok{ total\_claims }\SpecialCharTok{/}\NormalTok{ total\_exposure,}
        \AttributeTok{severity =}\NormalTok{ total\_claim\_amount }\SpecialCharTok{/}\NormalTok{ total\_claims,}
        \AttributeTok{expected\_claims =}\NormalTok{ frequency }\SpecialCharTok{*}\NormalTok{ severity}
\NormalTok{    ) }\SpecialCharTok{\%\textgreater{}\%}
    \FunctionTok{arrange}\NormalTok{(}\FunctionTok{desc}\NormalTok{(frequency)) }\SpecialCharTok{\%\textgreater{}\%}
    \FunctionTok{kable}\NormalTok{(}
        \AttributeTok{caption =} \StringTok{"Expected claims by municipality}\SpecialCharTok{\textbackslash{}n}\StringTok{Split by frequency and severity"}\NormalTok{,,}
        \AttributeTok{align =} \StringTok{"l"}\NormalTok{, }\AttributeTok{booktabs =} \ConstantTok{TRUE}
\NormalTok{    )}
\end{Highlighting}
\end{Shaded}

\begin{longtable}[]{@{}
  >{\raggedright\arraybackslash}p{(\linewidth - 12\tabcolsep) * \real{0.1942}}
  >{\raggedright\arraybackslash}p{(\linewidth - 12\tabcolsep) * \real{0.1262}}
  >{\raggedright\arraybackslash}p{(\linewidth - 12\tabcolsep) * \real{0.1456}}
  >{\raggedright\arraybackslash}p{(\linewidth - 12\tabcolsep) * \real{0.1845}}
  >{\raggedright\arraybackslash}p{(\linewidth - 12\tabcolsep) * \real{0.0971}}
  >{\raggedright\arraybackslash}p{(\linewidth - 12\tabcolsep) * \real{0.0971}}
  >{\raggedright\arraybackslash}p{(\linewidth - 12\tabcolsep) * \real{0.1553}}@{}}
\caption{Expected claims by municipality Split by frequency and
severity}\tabularnewline
\toprule\noalign{}
\begin{minipage}[b]{\linewidth}\raggedright
MUNICIPALITY\_binned
\end{minipage} & \begin{minipage}[b]{\linewidth}\raggedright
total\_claims
\end{minipage} & \begin{minipage}[b]{\linewidth}\raggedright
total\_exposure
\end{minipage} & \begin{minipage}[b]{\linewidth}\raggedright
total\_claim\_amount
\end{minipage} & \begin{minipage}[b]{\linewidth}\raggedright
frequency
\end{minipage} & \begin{minipage}[b]{\linewidth}\raggedright
severity
\end{minipage} & \begin{minipage}[b]{\linewidth}\raggedright
expected\_claims
\end{minipage} \\
\midrule\noalign{}
\endfirsthead
\toprule\noalign{}
\begin{minipage}[b]{\linewidth}\raggedright
MUNICIPALITY\_binned
\end{minipage} & \begin{minipage}[b]{\linewidth}\raggedright
total\_claims
\end{minipage} & \begin{minipage}[b]{\linewidth}\raggedright
total\_exposure
\end{minipage} & \begin{minipage}[b]{\linewidth}\raggedright
total\_claim\_amount
\end{minipage} & \begin{minipage}[b]{\linewidth}\raggedright
frequency
\end{minipage} & \begin{minipage}[b]{\linewidth}\raggedright
severity
\end{minipage} & \begin{minipage}[b]{\linewidth}\raggedright
expected\_claims
\end{minipage} \\
\midrule\noalign{}
\endhead
\bottomrule\noalign{}
\endlastfoot
Gand & 474 & 2484.419 & 756248.0 & 0.1907891 & 1595.4599 & 304.3963 \\
Liège & 329 & 1838.449 & 311331.2 & 0.1789552 & 946.2954 & 169.3445 \\
Charleroi & 565 & 3590.329 & 2681716.8 & 0.1573672 & 4746.4015 &
746.9279 \\
Mons & 222 & 1449.710 & 266246.7 & 0.1531341 & 1199.3094 & 183.6552 \\
Anvers & 521 & 3405.019 & 637351.8 & 0.1530094 & 1223.3240 & 187.1801 \\
La Louvière & 227 & 1658.945 & 319355.9 & 0.1368339 & 1406.8543 &
192.5054 \\
Other & 13771 & 101779.556 & 22302602.0 & 0.1353022 & 1619.5340 &
219.1265 \\
\end{longtable}

Even if some cities follow the high frequency, low severity pattern, we
do find outliers such as Charleroi where the severity is unusually high
or Liège where it falls unusually low. We then seem to capture something
more than the simple urban-rural divide which justifies keeping the
\texttt{MUNICIPALITY\_binned} variable in our model.

\subsubsection{Polynomial Features}\label{polynomial-features}

For the five numerical variables, we will consider the polynomial
features up to the third degree to capture potential non-linear
relationships with the response variable. This is an intermediate step
towards the GAM models where we will explicitly model non-linear
effects. Note that: - We will use Pearson correlation (continuous
variables) instead of Spearman correlation because the rank-rank
approach won't help us differentiate between the polynomial features of
the same positive variable (ranks are maintained). - We will look at the
p-values of the correlation tests instead of the correlation
coefficients themselves to assess the significance of the relationships.
- For the correlation with average claim amount (\texttt{AVG}), we will
only consider the claims where \texttt{NCLAIMS} is greater than zero. -
Orthogonal polynomial features will be considered separately to avoid
multicollinearity issues inherent in raw polynomial features. - To
compute orthogonal polynomial features, we will filter out the NA values
from the data.

\begin{Shaded}
\begin{Highlighting}[]
\NormalTok{numerical\_continuous }\OtherTok{\textless{}{-}} \FunctionTok{c}\NormalTok{(}\StringTok{"BM"}\NormalTok{, }\StringTok{"AGEPH"}\NormalTok{, }\StringTok{"AGEC"}\NormalTok{, }\StringTok{"POWER"}\NormalTok{, }\StringTok{"POPULATION"}\NormalTok{)}

\NormalTok{polynomial\_features }\OtherTok{\textless{}{-}} \FunctionTok{lapply}\NormalTok{(numerical\_continuous, }\ControlFlowTok{function}\NormalTok{(var) \{}
    \FunctionTok{poly}\NormalTok{(data\_enriched[[var]], }\AttributeTok{degree =} \DecValTok{3}\NormalTok{, }\AttributeTok{raw =} \ConstantTok{TRUE}\NormalTok{)}
\NormalTok{\})}
\FunctionTok{names}\NormalTok{(polynomial\_features) }\OtherTok{\textless{}{-}}\NormalTok{ numerical\_continuous}
\CommentTok{\# Initialize data frame with proper structure}
\NormalTok{correlation\_matrix }\OtherTok{\textless{}{-}} \FunctionTok{data.frame}\NormalTok{(}
    \FunctionTok{matrix}\NormalTok{(}\ConstantTok{NA}\NormalTok{, }\AttributeTok{nrow =} \FunctionTok{length}\NormalTok{(numerical\_continuous), }\AttributeTok{ncol =} \DecValTok{6}\NormalTok{)}
\NormalTok{)}
\FunctionTok{rownames}\NormalTok{(correlation\_matrix) }\OtherTok{\textless{}{-}}\NormalTok{ numerical\_continuous}
\FunctionTok{colnames}\NormalTok{(correlation\_matrix) }\OtherTok{\textless{}{-}} \FunctionTok{c}\NormalTok{(}
    \StringTok{"Claims Poly1"}\NormalTok{, }\StringTok{"Claims Poly2"}\NormalTok{, }\StringTok{"Claims Poly3"}\NormalTok{,}
    \StringTok{"Avg Poly1"}\NormalTok{, }\StringTok{"Avg Poly2"}\NormalTok{, }\StringTok{"Avg Poly3"}\NormalTok{)}

\CommentTok{\# Fill the matrix}
\ControlFlowTok{for}\NormalTok{ (var }\ControlFlowTok{in}\NormalTok{ numerical\_continuous) \{}
\NormalTok{    poly\_values }\OtherTok{\textless{}{-}}\NormalTok{ polynomial\_features[[var]]}
    \ControlFlowTok{for}\NormalTok{ (j }\ControlFlowTok{in} \DecValTok{1}\SpecialCharTok{:}\DecValTok{3}\NormalTok{) \{}
\NormalTok{        correlation\_matrix[var, j] }\OtherTok{\textless{}{-}} \FunctionTok{cor.test}\NormalTok{(data\_enriched[[}\StringTok{"NCLAIMS"}\NormalTok{]], poly\_values[, j], }\AttributeTok{method =} \StringTok{"pearson"}\NormalTok{)}\SpecialCharTok{$}\NormalTok{p.value}
\NormalTok{        claims\_only\_avg }\OtherTok{\textless{}{-}}\NormalTok{ data\_enriched[[}\StringTok{"AVG"}\NormalTok{]][data\_enriched[[}\StringTok{"NCLAIMS"}\NormalTok{]] }\SpecialCharTok{\textgreater{}} \DecValTok{0}\NormalTok{]}
\NormalTok{        claims\_only\_poly\_values }\OtherTok{\textless{}{-}}\NormalTok{ poly\_values[data\_enriched[[}\StringTok{"NCLAIMS"}\NormalTok{]] }\SpecialCharTok{\textgreater{}} \DecValTok{0}\NormalTok{, j]}
\NormalTok{        correlation\_matrix[var, j }\SpecialCharTok{+} \DecValTok{3}\NormalTok{] }\OtherTok{\textless{}{-}} \FunctionTok{cor.test}\NormalTok{(claims\_only\_avg, claims\_only\_poly\_values, }\AttributeTok{method =} \StringTok{"pearson"}\NormalTok{)}\SpecialCharTok{$}\NormalTok{p.value}
\NormalTok{    \}}
\NormalTok{\}}

\FunctionTok{pheatmap}\NormalTok{(correlation\_matrix,}
    \AttributeTok{main =} \StringTok{"Correlation P{-}Values Heatmap {-} Polynomial Features}\SpecialCharTok{\textbackslash{}n}\StringTok{(Green = Significant p\textless{}0.05, Red = Non{-}significant)"}\NormalTok{,}
    \AttributeTok{display\_numbers =} \ConstantTok{TRUE}\NormalTok{,}
    \AttributeTok{number\_format =} \StringTok{"\%.4f"}\NormalTok{,}
    \AttributeTok{breaks =} \FunctionTok{c}\NormalTok{(}\DecValTok{0}\NormalTok{, }\FloatTok{0.01}\NormalTok{, }\FloatTok{0.05}\NormalTok{, }\DecValTok{1}\NormalTok{),}
    \AttributeTok{color =} \FunctionTok{c}\NormalTok{(}\StringTok{"lightgreen"}\NormalTok{, }\StringTok{"white"}\NormalTok{, }\StringTok{"lightcoral"}\NormalTok{),}
    \AttributeTok{legend =} \ConstantTok{FALSE}\NormalTok{,}
    \AttributeTok{cluster\_rows =} \ConstantTok{FALSE}\NormalTok{,}
    \AttributeTok{cluster\_cols =} \ConstantTok{FALSE}
\NormalTok{)}
\end{Highlighting}
\end{Shaded}

\begin{center}\includegraphics{project_MD_files/figure-latex/polynomial_features-1} \end{center}

We first notice that \texttt{AGEC}is not significantly correlated with
\texttt{NCLAIMS} when considered with the first and second-degree
polynomial features. As we saw before, the non-linear effect is captures
in our binning. An alternative approach could have been to use the
third-degree polynomial. We would still have the issue of fixed effects
being imposed by the polynomial terms, which might not fully capture the
underlying non-linear relationship. This is a good example of why
binning can sometimes be more flexible in capturing non-linear
relationships compared to fixed polynomial terms.

We notice that the correllations between the polynomial features of a
same variable are similar, this is foreseeable given that higher-degree
polynomials are derived from the lower-degree ones. We are here exposed
to multicollinearity issues. A solution is to check the
\textbf{orthogonal polynomial features} instead of the raw ones, as they
are uncorrelated by construction. This also enables us to differentiate
the individual contributions of each polynomial term more clearly in the
modeling process (linear, quadratic, cubic).

\begin{Shaded}
\begin{Highlighting}[]
\NormalTok{no\_na }\OtherTok{\textless{}{-}} \FunctionTok{na.omit}\NormalTok{(data\_enriched)}
\NormalTok{orthogonal\_polynomial\_features }\OtherTok{\textless{}{-}} \FunctionTok{lapply}\NormalTok{(numerical\_continuous, }\ControlFlowTok{function}\NormalTok{(var) \{}
    \FunctionTok{poly}\NormalTok{(no\_na[[var]], }\AttributeTok{degree =} \DecValTok{3}\NormalTok{, }\AttributeTok{raw =} \ConstantTok{FALSE}\NormalTok{)}
\NormalTok{\})}
\FunctionTok{names}\NormalTok{(orthogonal\_polynomial\_features) }\OtherTok{\textless{}{-}}\NormalTok{ numerical\_continuous}
\NormalTok{orth\_correlation\_matrix }\OtherTok{\textless{}{-}} \FunctionTok{data.frame}\NormalTok{(}
    \FunctionTok{matrix}\NormalTok{(}\ConstantTok{NA}\NormalTok{, }\AttributeTok{nrow =} \FunctionTok{length}\NormalTok{(numerical\_continuous), }\AttributeTok{ncol =} \DecValTok{6}\NormalTok{)}
\NormalTok{)}
\FunctionTok{rownames}\NormalTok{(orth\_correlation\_matrix) }\OtherTok{\textless{}{-}}\NormalTok{ numerical\_continuous}
\FunctionTok{colnames}\NormalTok{(orth\_correlation\_matrix) }\OtherTok{\textless{}{-}} \FunctionTok{c}\NormalTok{(}
    \StringTok{"OrthClaims Poly1"}\NormalTok{, }\StringTok{"OrthClaims Poly2"}\NormalTok{, }\StringTok{"OrthClaims Poly3"}\NormalTok{,}
    \StringTok{"OrthAvg Poly1"}\NormalTok{, }\StringTok{"OrthAvg Poly2"}\NormalTok{, }\StringTok{"OrthAvg Poly3"}\NormalTok{)}

\CommentTok{\# Fill the matrix}
\ControlFlowTok{for}\NormalTok{ (var }\ControlFlowTok{in}\NormalTok{ numerical\_continuous) \{}
\NormalTok{    poly\_values }\OtherTok{\textless{}{-}}\NormalTok{ orthogonal\_polynomial\_features[[var]]}
    \ControlFlowTok{for}\NormalTok{ (j }\ControlFlowTok{in} \DecValTok{1}\SpecialCharTok{:}\DecValTok{3}\NormalTok{) \{}
\NormalTok{        orth\_correlation\_matrix[var, j] }\OtherTok{\textless{}{-}} \FunctionTok{cor.test}\NormalTok{(no\_na[[}\StringTok{"NCLAIMS"}\NormalTok{]], poly\_values[, j], }\AttributeTok{method =} \StringTok{"pearson"}\NormalTok{)}\SpecialCharTok{$}\NormalTok{p.value}
\NormalTok{        claims\_only\_avg }\OtherTok{\textless{}{-}}\NormalTok{ no\_na[[}\StringTok{"AVG"}\NormalTok{]][no\_na[[}\StringTok{"NCLAIMS"}\NormalTok{]] }\SpecialCharTok{\textgreater{}} \DecValTok{0}\NormalTok{]}
\NormalTok{        claims\_only\_poly\_values }\OtherTok{\textless{}{-}}\NormalTok{ poly\_values[no\_na[[}\StringTok{"NCLAIMS"}\NormalTok{]] }\SpecialCharTok{\textgreater{}} \DecValTok{0}\NormalTok{, j]}
\NormalTok{        orth\_correlation\_matrix[var, j }\SpecialCharTok{+} \DecValTok{3}\NormalTok{] }\OtherTok{\textless{}{-}} \FunctionTok{cor.test}\NormalTok{(claims\_only\_avg, claims\_only\_poly\_values, }\AttributeTok{method =} \StringTok{"pearson"}\NormalTok{)}\SpecialCharTok{$}\NormalTok{p.value}
\NormalTok{    \}}
\NormalTok{\}}

\FunctionTok{pheatmap}\NormalTok{(orth\_correlation\_matrix,}
    \AttributeTok{main =} \StringTok{"Correlation P{-}Values Heatmap {-} Orthogonal Polynomial Features}\SpecialCharTok{\textbackslash{}n}\StringTok{(Green = Significant p\textless{}0.05, Red = Non{-}significant)"}\NormalTok{,}
    \AttributeTok{display\_numbers =} \ConstantTok{TRUE}\NormalTok{,}
    \AttributeTok{number\_format =} \StringTok{"\%.4f"}\NormalTok{,}
    \AttributeTok{breaks =} \FunctionTok{c}\NormalTok{(}\DecValTok{0}\NormalTok{, }\FloatTok{0.01}\NormalTok{, }\FloatTok{0.05}\NormalTok{, }\DecValTok{1}\NormalTok{),}
    \AttributeTok{color =} \FunctionTok{c}\NormalTok{(}\StringTok{"lightgreen"}\NormalTok{, }\StringTok{"white"}\NormalTok{, }\StringTok{"lightcoral"}\NormalTok{),}
    \AttributeTok{legend =} \ConstantTok{FALSE}\NormalTok{,}
    \AttributeTok{cluster\_rows =} \ConstantTok{FALSE}\NormalTok{,}
    \AttributeTok{cluster\_cols =} \ConstantTok{FALSE}
\NormalTok{) }
\end{Highlighting}
\end{Shaded}

\begin{center}\includegraphics{project_MD_files/figure-latex/orthogonal_polynomial_features-1} \end{center}

From this orthogonal polynomial feature exploration we can conclude
that:

\begin{itemize}
\tightlist
\item
  Polynomial features are more significantly correlated with the number
  of claims (\texttt{NCLAIMS}) than with the average claim amount
  (\texttt{AVG}).
\item
  Orthogonal polynomial features help assess whether each polynomial
  term contributes significantly to the relationship with the target
  variables, without being confounded by multicollinearity.
\item
  There is no clear rule to know if first, second, and/or third-degree
  polynomial terms are significant. The WoE-based binning approach
  provides a more flexible and interpretable way to capture non-linear
  effects.
\end{itemize}

In conclusion, polynomials are a way to capture non-linear
relationships. However, the incremental effect of these polynomial terms
is still fixed in a GLM. We therefore prefer working with binned
features rather than relying solely on polynomial terms. We will
nevertheless see if the smoothing functions of the GAM models confirm
the conclusions drawn from our exploratory analysis.

\subsubsection{Interaction Features}\label{interaction-features}

Interactions are usually added as a result of business knowledge. In our
case, we know from a theorical background that the age of the
policyholder interacts with the power of the car. We can add this
interaction effect and test its significance in our models. Instead of
manually testing all possible combinations. We'll use a decision tree
approach which nature is to look for interactions and splits that best
explain the variability in the target variable. With a shallow decision
tree (depth of 3), we can identify the most relevant interactions
without overfitting the data. Watch out the decision trees are not
robusts and can be sensitive to small changes in the data. We'll do a
10-fold cross-validation to mitigate this issue.

\begin{Shaded}
\begin{Highlighting}[]
\NormalTok{tree\_model }\OtherTok{\textless{}{-}} \FunctionTok{rpart}\NormalTok{(}
\NormalTok{    NCLAIMS }\SpecialCharTok{\textasciitilde{}}\NormalTok{ COVERAGE\_BIN }\SpecialCharTok{+}\NormalTok{ FUEL }\SpecialCharTok{+}\NormalTok{ FLEET }\SpecialCharTok{+}\NormalTok{ SEX }\SpecialCharTok{+} \CommentTok{\# Categorical features}
\NormalTok{    BM\_binned }\SpecialCharTok{+}\NormalTok{ AGEPH\_binned\_merged }\SpecialCharTok{+}\NormalTok{ AGEC\_binned }\SpecialCharTok{+}\NormalTok{ POWER\_binned\_merged }\SpecialCharTok{+} \CommentTok{\# Binned numerical features}
\NormalTok{    POPULATION\_binned }\SpecialCharTok{+}\NormalTok{ REGION\_merged }\SpecialCharTok{+}\NormalTok{ MUNICIPALITY\_binned, }\CommentTok{\# Enriched features}
    \AttributeTok{data =}\NormalTok{ data\_enriched,}
    \AttributeTok{method =} \StringTok{"poisson"}\NormalTok{,}
    \AttributeTok{control =} \FunctionTok{rpart.control}\NormalTok{(}\AttributeTok{maxdepth =} \DecValTok{3}\NormalTok{, }\AttributeTok{xval =} \DecValTok{10}\NormalTok{, }\AttributeTok{cp =} \DecValTok{0}\NormalTok{)}
\NormalTok{)}
\FunctionTok{rpart.plot}\NormalTok{(tree\_model)}
\end{Highlighting}
\end{Shaded}

\begin{center}\includegraphics{project_MD_files/figure-latex/interaction-tree-1} \end{center}

\begin{Shaded}
\begin{Highlighting}[]
\NormalTok{importance\_df }\OtherTok{\textless{}{-}} \FunctionTok{data.frame}\NormalTok{(}
    \AttributeTok{feature =} \FunctionTok{names}\NormalTok{(tree\_model}\SpecialCharTok{$}\NormalTok{variable.importance),}
    \AttributeTok{importance =}\NormalTok{ tree\_model}\SpecialCharTok{$}\NormalTok{variable.importance }\SpecialCharTok{\%\textgreater{}\%} \StringTok{\textasciigrave{}}\AttributeTok{/}\StringTok{\textasciigrave{}}\NormalTok{(}\FunctionTok{sum}\NormalTok{(.))}
\NormalTok{)}
\NormalTok{importance\_df }\OtherTok{\textless{}{-}}\NormalTok{ importance\_df }\SpecialCharTok{\%\textgreater{}\%} \FunctionTok{arrange}\NormalTok{(}\FunctionTok{desc}\NormalTok{(importance))}

\FunctionTok{ggplot}\NormalTok{(importance\_df, }\FunctionTok{aes}\NormalTok{(}\AttributeTok{x =} \FunctionTok{reorder}\NormalTok{(feature, importance), }\AttributeTok{y =}\NormalTok{ importance)) }\SpecialCharTok{+}
    \FunctionTok{geom\_bar}\NormalTok{(}\AttributeTok{stat =} \StringTok{"identity"}\NormalTok{, }\AttributeTok{fill =} \StringTok{"steelblue"}\NormalTok{) }\SpecialCharTok{+}
    \FunctionTok{coord\_flip}\NormalTok{() }\SpecialCharTok{+}
    \FunctionTok{labs}\NormalTok{(}\AttributeTok{title =} \StringTok{"Feature Importance in Decision Tree"}\NormalTok{,}
            \AttributeTok{x =} \StringTok{"Feature"}\NormalTok{,}
            \AttributeTok{y =} \StringTok{"Importance Score"}\NormalTok{) }\SpecialCharTok{+}
    \FunctionTok{theme\_minimal}\NormalTok{()}
\end{Highlighting}
\end{Shaded}

\begin{center}\includegraphics{project_MD_files/figure-latex/interaction-tree-2} \end{center}

Based on this machine learning approach, we can identify both specific
interactions and the most important features that influence the target
variable. When adding interaction terms to our models, we can thus
consider:

\textbf{Specific Interactions:} representing the two first splits
identified by the decision tree.

\begin{itemize}
\tightlist
\item
  \texttt{BM\_binned.{[}0,1{]}} x \texttt{FUEL.gasoline}
\item
  \texttt{BM\_binned.{[}0,1{]}} x \texttt{!FUEL.gasoline}
\item
  \texttt{!BM\_binned.{[}0,1{]}} x \texttt{REGION\_merged.Others}: will
  not add it as it is captured below already
\item
  \texttt{!BM\_binned.{[}0,1{]}} x \texttt{!REGION\_merged.Others}: will
  not add it as it is captured below already
\end{itemize}

\begin{Shaded}
\begin{Highlighting}[]
\NormalTok{data\_enriched }\OtherTok{\textless{}{-}}\NormalTok{ data\_enriched }\SpecialCharTok{\%\textgreater{}\%}
    \FunctionTok{mutate}\NormalTok{(}
        \AttributeTok{BM\_01\_FUEL\_gasoline =}\NormalTok{ BM\_binned }\SpecialCharTok{==} \StringTok{"[0,1]"} \SpecialCharTok{\&}\NormalTok{ FUEL }\SpecialCharTok{==} \StringTok{"gasoline"}\NormalTok{,}
        \AttributeTok{BM\_01\_FUEL\_not\_gasoline =}\NormalTok{ BM\_binned }\SpecialCharTok{==} \StringTok{"[0,1]"} \SpecialCharTok{\&}\NormalTok{ FUEL }\SpecialCharTok{!=} \StringTok{"gasoline"}
\NormalTok{    )}
\end{Highlighting}
\end{Shaded}

\textbf{Variable Interactions:} representing the 3 most important
features identified by the decision tree.

\begin{itemize}
\tightlist
\item
  \texttt{BM\_binned} x \texttt{AGEPH\_binned\_merged}
\item
  \texttt{REGION\_merged} x \texttt{POWER\_binned\_merged}
\item
  \texttt{BM\_binned} x \texttt{REGION\_merged}
\end{itemize}

\section{Modelling}\label{modelling}

Now that we have explored, transformed, and enriched our data. We are
ready to proceed with the modelling phase. We have a rich set of
features to gradually incorporate into our models to improve predictive
performance. We will also clean all intermediary objects that won't be
used for modelling.

\begin{Shaded}
\begin{Highlighting}[]
\CommentTok{\# Clean intermediary objects not used for modelling}
\NormalTok{functions }\OtherTok{\textless{}{-}} \FunctionTok{ls}\NormalTok{()[}\FunctionTok{sapply}\NormalTok{(}\FunctionTok{ls}\NormalTok{(), }\ControlFlowTok{function}\NormalTok{(x) }\FunctionTok{is.function}\NormalTok{(}\FunctionTok{get}\NormalTok{(x)))]}
\FunctionTok{rm}\NormalTok{(}\AttributeTok{list =} \FunctionTok{setdiff}\NormalTok{(}\FunctionTok{ls}\NormalTok{(), }\FunctionTok{c}\NormalTok{(}
    \StringTok{"data"}\NormalTok{, }\StringTok{"data\_enriched"}\NormalTok{, }\StringTok{"holdout"}\NormalTok{, }\StringTok{"mtpl\_be"}\NormalTok{,}
    \StringTok{"conv\_matching"}\NormalTok{, }\StringTok{"refnis"}\NormalTok{, }\StringTok{"pop"}\NormalTok{, }
    \StringTok{"result\_named\_list"}\NormalTok{, }\StringTok{"n\_bins\_named\_list"}\NormalTok{,}
\NormalTok{    functions)))}
\FunctionTok{cat}\NormalTok{(}\FunctionTok{sprintf}\NormalTok{(}\StringTok{"Remaining objects and functions: \%s"}\NormalTok{, }\FunctionTok{paste}\NormalTok{(}\FunctionTok{ls}\NormalTok{(), }\AttributeTok{collapse =} \StringTok{", "}\NormalTok{)))}
\end{Highlighting}
\end{Shaded}

\begin{verbatim}
Remaining objects and functions: assert_ref_level, conv_matching, data, data_enriched, holdout, iv_elbow, mtpl_be, n_bins_named_list, plot_claim_freq, plot_iv_elbow, plot_woe_iv, pop, qbinning, rebuild_bins, refnis, result_named_list, woe_calc
\end{verbatim}

\begin{Shaded}
\begin{Highlighting}[]
\NormalTok{original\_categorical\_features }\OtherTok{\textless{}{-}} \FunctionTok{c}\NormalTok{(}\StringTok{"COVERAGE\_BIN"}\NormalTok{, }\StringTok{"FUEL"}\NormalTok{, }\StringTok{"FLEET"}\NormalTok{, }\StringTok{"SEX"}\NormalTok{)}
\NormalTok{continuous\_features }\OtherTok{\textless{}{-}} \FunctionTok{c}\NormalTok{(}\StringTok{"BM"}\NormalTok{, }\StringTok{"AGEPH"}\NormalTok{, }\StringTok{"AGEC"}\NormalTok{, }\StringTok{"POWER"}\NormalTok{)}
\NormalTok{binned\_continuous\_features }\OtherTok{\textless{}{-}} \FunctionTok{c}\NormalTok{(}\StringTok{"BM\_binned"}\NormalTok{, }\StringTok{"AGEPH\_binned\_merged"}\NormalTok{, }\StringTok{"AGEC\_binned"}\NormalTok{, }\StringTok{"POWER\_binned\_merged"}\NormalTok{)}
\NormalTok{enriched\_features }\OtherTok{\textless{}{-}} \FunctionTok{c}\NormalTok{(}\StringTok{"REGION\_merged"}\NormalTok{, }\StringTok{"MUNICIPALITY\_binned"}\NormalTok{, }\StringTok{"POPULATION"}\NormalTok{)}
\NormalTok{binned\_enriched\_features }\OtherTok{\textless{}{-}} \FunctionTok{c}\NormalTok{(}\StringTok{"REGION\_merged"}\NormalTok{, }\StringTok{"MUNICIPALITY\_binned"}\NormalTok{, }\StringTok{"POPULATION\_binned"}\NormalTok{)}
\NormalTok{interaction\_features }\OtherTok{\textless{}{-}} \FunctionTok{c}\NormalTok{(}\StringTok{"BM\_01\_FUEL\_gasoline"}\NormalTok{, }\StringTok{"BM\_01\_FUEL\_not\_gasoline"}\NormalTok{, }\StringTok{"BM\_binned:AGEPH\_binned\_merged"}\NormalTok{, }\StringTok{"REGION\_merged:AGEPH\_binned\_merged"}\NormalTok{, }\StringTok{"BM\_binned:REGION\_merged"}\NormalTok{)}

\NormalTok{levels\_named\_list }\OtherTok{\textless{}{-}} \FunctionTok{list}\NormalTok{()}
\ControlFlowTok{for}\NormalTok{ (col }\ControlFlowTok{in} \FunctionTok{c}\NormalTok{(original\_categorical\_features, binned\_continuous\_features, binned\_enriched\_features)) \{}
    \ControlFlowTok{if}\NormalTok{ (}\SpecialCharTok{!}\NormalTok{col }\SpecialCharTok{\%in\%} \FunctionTok{colnames}\NormalTok{(data\_enriched)) \{}\FunctionTok{stop}\NormalTok{(}\FunctionTok{sprintf}\NormalTok{(}\StringTok{"Column \%s is missing from data\_enriched"}\NormalTok{, col))\}}
\NormalTok{    levels\_named\_list[[col]] }\OtherTok{\textless{}{-}} \FunctionTok{levels}\NormalTok{(data\_enriched[[col]])}
\NormalTok{\}}
\end{Highlighting}
\end{Shaded}

\textbf{Original Categorical Features:} Note that \texttt{USE} was
excluded because it was not statistically significant in preliminary
analyses. And \texttt{COVERAGE} was transformed into
\texttt{COVERAGE\_BIN} after merging the partial and full omnium levels
that did not significantly add information.

\begin{itemize}
\tightlist
\item
  \texttt{COVERAGE\_BIN}
\item
  \texttt{FUEL}
\item
  \texttt{FLEET}
\item
  \texttt{SEX}
\end{itemize}

\textbf{Continuous Features:} The continuous features as such will be
included in the GAM models only. We binned and merged them into
categorical variables for the GLM modelling.

\begin{itemize}
\tightlist
\item
  \texttt{BM} -\textgreater{} \texttt{BM\_binned}
\item
  \texttt{AGEPH} -\textgreater{} \texttt{AGEPH\_binned\_merged}
\item
  \texttt{AGEC} -\textgreater{} \texttt{AGEC\_binned}
\item
  \texttt{POWER} -\textgreater{} \texttt{POWER\_binned\_merged}
\end{itemize}

\textbf{Enriched Features:} Based on geospation information, we have
created, binned and merged both categorical and continuous features to
capture urban/rural distinctions and municipality-level characteristics.

\begin{itemize}
\tightlist
\item
  \texttt{REGION\_merged}
\item
  \texttt{MUNICIPALITY\_merged}
\item
  \texttt{POPULATION} -\textgreater{} \texttt{POPULATION\_binned}
\end{itemize}

\textbf{Interaction Features:} We have also created specific and
variable interaction features based on expert knowledge and the decision
tree analysis.

\begin{itemize}
\tightlist
\item
  \texttt{POWER\_binned\_merged} x \texttt{AGEPH\_binned\_merged}:
  expert knowledge
\item
  \texttt{BM\_01\_FUEL\_gasoline}: specific interaction based on
  decision tree analysis
\item
  \texttt{BM\_01\_FUEL\_not\_gasoline}: specific interaction based on
  decision tree analysis
\item
  \texttt{BM\_binned} x \texttt{AGEPH\_binned\_merged}: variable
  interaction based on feature importance
\item
  \texttt{REGION\_merged} x \texttt{AGEPH\_binned\_merged}: variable
  interaction based on feature importance
\item
  \texttt{BM\_binned} x \texttt{REGION\_merged}: variable interaction
  based on feature importance
\end{itemize}

\textbf{Polynomial Features:} We also considered raw and orthogonal
polynomial features for the continuous variables to capture potential
non-linear relationships in the GAM models. We decided not to include
them in the GLM models for the sake of model simplicity and
interpretability. These polynomial effects will be naturally captured in
the GAM models through the smooth functions. We will then be able to see
whether the GAM results confirm our exploratory analysis.

\subsection{Validation Framework}\label{validation-framework}

Remember that at the start of the Data Preparation chapter, we kept a
holdout set separate from the training data for model validation
purposes. We will now create a workflow including:

\begin{itemize}
\tightlist
\item
  \textbf{Data Preparation}: enriching the dataset and transforming
  features as described in the previous sections.
\item
  \textbf{Model Training}: fitting the selected models to the prepared
  training data with a 5-fold cross-validation.
\item
  \textbf{Model Validation}: evaluating the performance of the trained
  models on the holdout set.
\end{itemize}

\subsubsection{Data Preparation
Workflow}\label{data-preparation-workflow}

\begin{Shaded}
\begin{Highlighting}[]
\CommentTok{\# Data enrichment}

\NormalTok{data\_enrichment }\OtherTok{\textless{}{-}} \ControlFlowTok{function}\NormalTok{(input\_data) \{}
    \ControlFlowTok{if}\NormalTok{ (}\SpecialCharTok{!}\NormalTok{(}\StringTok{"PC"} \SpecialCharTok{\%in\%} \FunctionTok{colnames}\NormalTok{(input\_data))) \{}\FunctionTok{stop}\NormalTok{(}\StringTok{"PC column is missing from raw\_data"}\NormalTok{)\}}
    \ControlFlowTok{if}\NormalTok{ (}\SpecialCharTok{!}\NormalTok{(}\StringTok{"ID"} \SpecialCharTok{\%in\%} \FunctionTok{colnames}\NormalTok{(input\_data))) \{}\FunctionTok{stop}\NormalTok{(}\StringTok{"ID column is missing from raw\_data"}\NormalTok{)\}}
    \ControlFlowTok{if}\NormalTok{ (}\SpecialCharTok{!}\FunctionTok{exists}\NormalTok{(}\StringTok{"conv\_matching"}\NormalTok{)) \{}\FunctionTok{stop}\NormalTok{(}\StringTok{"conv\_matching object is missing"}\NormalTok{)\}}
    \ControlFlowTok{if}\NormalTok{ (}\SpecialCharTok{!}\FunctionTok{exists}\NormalTok{(}\StringTok{"refnis"}\NormalTok{)) \{}\FunctionTok{stop}\NormalTok{(}\StringTok{"refnis object is missing"}\NormalTok{)\}}
    \ControlFlowTok{if}\NormalTok{ (}\SpecialCharTok{!}\FunctionTok{exists}\NormalTok{(}\StringTok{"pop"}\NormalTok{)) \{}\FunctionTok{stop}\NormalTok{(}\StringTok{"pop object is missing"}\NormalTok{)\}}
\NormalTok{    output\_data }\OtherTok{\textless{}{-}}\NormalTok{ input\_data }\SpecialCharTok{\%\textgreater{}\%}
        \FunctionTok{mutate}\NormalTok{(}\AttributeTok{TEMP\_PC =} \FunctionTok{as.integer}\NormalTok{(PC)) }\SpecialCharTok{\%\textgreater{}\%}
        \FunctionTok{left\_join}\NormalTok{(}\FunctionTok{select}\NormalTok{(conv\_matching, }\FunctionTok{c}\NormalTok{(}\StringTok{"Postal.code"}\NormalTok{,}\StringTok{"Refnis.code"}\NormalTok{)), }\AttributeTok{by =} \FunctionTok{c}\NormalTok{(}\StringTok{"TEMP\_PC"} \OtherTok{=} \StringTok{"Postal.code"}\NormalTok{)) }\SpecialCharTok{\%\textgreater{}\%}
        \FunctionTok{left\_join}\NormalTok{(}\FunctionTok{select}\NormalTok{(refnis, }\FunctionTok{c}\NormalTok{(}\StringTok{"NIS\_MUN"}\NormalTok{, }\StringTok{"MUNICIPALITY"}\NormalTok{, }\StringTok{"REGION"}\NormalTok{)), }\AttributeTok{by =} \FunctionTok{c}\NormalTok{(}\StringTok{"Refnis.code"} \OtherTok{=} \StringTok{"NIS\_MUN"}\NormalTok{)) }\SpecialCharTok{\%\textgreater{}\%}
        \FunctionTok{left\_join}\NormalTok{(}\FunctionTok{select}\NormalTok{(pop, }\FunctionTok{c}\NormalTok{(}\StringTok{"NIS.code"}\NormalTok{, }\StringTok{"Totaal"}\NormalTok{)), }\AttributeTok{by =} \FunctionTok{c}\NormalTok{(}\StringTok{"Refnis.code"} \OtherTok{=} \StringTok{"NIS.code"}\NormalTok{)) }\SpecialCharTok{\%\textgreater{}\%}
        \FunctionTok{rename}\NormalTok{(}
            \AttributeTok{REFNIS =}\NormalTok{ Refnis.code,}
            \AttributeTok{MUNICIPALITY =}\NormalTok{ MUNICIPALITY,}
            \AttributeTok{REGION =}\NormalTok{ REGION,}
            \AttributeTok{POPULATION =}\NormalTok{ Totaal}
\NormalTok{        ) }\SpecialCharTok{\%\textgreater{}\%}
        \FunctionTok{select}\NormalTok{(}\SpecialCharTok{{-}}\NormalTok{TEMP\_PC) }\SpecialCharTok{\%\textgreater{}\%}
        \FunctionTok{mutate}\NormalTok{(}
            \AttributeTok{POPULATION =} \FunctionTok{as.numeric}\NormalTok{(}\FunctionTok{gsub}\NormalTok{(}\StringTok{"}\SpecialCharTok{\textbackslash{}\textbackslash{}}\StringTok{."}\NormalTok{, }\StringTok{""}\NormalTok{, POPULATION)),}
\NormalTok{        ) }\SpecialCharTok{\%\textgreater{}\%}
        \FunctionTok{mutate}\NormalTok{(}
            \AttributeTok{POPULATION =} \FunctionTok{ifelse}\NormalTok{(}\FunctionTok{is.na}\NormalTok{(POPULATION), }\FunctionTok{mean}\NormalTok{(POPULATION, }\AttributeTok{na.rm =} \ConstantTok{TRUE}\NormalTok{), POPULATION)}
\NormalTok{        )}
    \ControlFlowTok{if}\NormalTok{ (}\FunctionTok{any}\NormalTok{(}\FunctionTok{duplicated}\NormalTok{(output\_data}\SpecialCharTok{$}\NormalTok{ID))) \{}\FunctionTok{stop}\NormalTok{(}\StringTok{"Duplicated IDs found in enriched\_data"}\NormalTok{)\}}
    \ControlFlowTok{if}\NormalTok{ (}\FunctionTok{nrow}\NormalTok{(output\_data) }\SpecialCharTok{!=} \FunctionTok{nrow}\NormalTok{(input\_data)) \{}\FunctionTok{stop}\NormalTok{(}\StringTok{"Row count mismatch between enriched\_data and raw\_data"}\NormalTok{)\}    }
    \FunctionTok{return}\NormalTok{(output\_data)}
\NormalTok{\}}

\CommentTok{\# Step 2: Data Preparation}

\NormalTok{data\_preparation\_original\_categorical }\OtherTok{\textless{}{-}} \ControlFlowTok{function}\NormalTok{(input\_data) \{}
    \CommentTok{\# Convert original categorical features to factors and validate their levels}
    \ControlFlowTok{if}\NormalTok{ (}\SpecialCharTok{!}\FunctionTok{exists}\NormalTok{(}\StringTok{"original\_categorical\_features"}\NormalTok{)) \{}\FunctionTok{stop}\NormalTok{(}\StringTok{"original\_categorical\_features object is missing"}\NormalTok{)\}}
    \ControlFlowTok{if}\NormalTok{ (}\SpecialCharTok{!}\FunctionTok{exists}\NormalTok{(}\StringTok{"levels\_named\_list"}\NormalTok{)) \{}\FunctionTok{stop}\NormalTok{(}\StringTok{"levels\_named\_list object is missing"}\NormalTok{)\}}
\NormalTok{    temp\_data }\OtherTok{\textless{}{-}}\NormalTok{ input\_data}
    \ControlFlowTok{for}\NormalTok{ (col }\ControlFlowTok{in}\NormalTok{ original\_categorical\_features) \{}
\NormalTok{        base\_col }\OtherTok{\textless{}{-}}\NormalTok{ col }\SpecialCharTok{\%\textgreater{}\%} \FunctionTok{gsub}\NormalTok{(}\StringTok{"\_binned"}\NormalTok{, }\StringTok{""}\NormalTok{, .) }\SpecialCharTok{\%\textgreater{}\%} \FunctionTok{gsub}\NormalTok{(}\StringTok{"\_merged"}\NormalTok{, }\StringTok{""}\NormalTok{, .) }\SpecialCharTok{\%\textgreater{}\%} \FunctionTok{gsub}\NormalTok{(}\StringTok{"\_BIN"}\NormalTok{, }\StringTok{""}\NormalTok{, .)}
        \ControlFlowTok{if}\NormalTok{ (}\SpecialCharTok{!}\NormalTok{base\_col }\SpecialCharTok{\%in\%} \FunctionTok{colnames}\NormalTok{(temp\_data)) \{}\FunctionTok{stop}\NormalTok{(}\FunctionTok{sprintf}\NormalTok{(}\StringTok{"Column \%s is missing from input\_data"}\NormalTok{, base\_col))\}}
\NormalTok{        temp\_data[[base\_col]] }\OtherTok{\textless{}{-}} \FunctionTok{as.factor}\NormalTok{(temp\_data[[base\_col]])}
\NormalTok{    \}}
\NormalTok{    temp\_data}\SpecialCharTok{$}\NormalTok{COVERAGE\_BIN }\OtherTok{\textless{}{-}}\NormalTok{ temp\_data}\SpecialCharTok{$}\NormalTok{COVERAGE }\SpecialCharTok{\%\textgreater{}\%} \FunctionTok{fct\_collapse}\NormalTok{(}
        \StringTok{"TPL"} \OtherTok{=} \FunctionTok{c}\NormalTok{(}\StringTok{"TPL"}\NormalTok{),}
        \StringTok{"OMN"} \OtherTok{=} \FunctionTok{c}\NormalTok{(}\StringTok{"PO"}\NormalTok{, }\StringTok{"FO"}\NormalTok{)}
\NormalTok{    )}
    \ControlFlowTok{for}\NormalTok{ (col }\ControlFlowTok{in}\NormalTok{ original\_categorical\_features) \{}
\NormalTok{        temp\_data[[col]] }\OtherTok{\textless{}{-}} \FunctionTok{factor}\NormalTok{(temp\_data[[col]], }\AttributeTok{levels =}\NormalTok{ levels\_named\_list[[col]])}
\NormalTok{    \}}
\NormalTok{    output\_data }\OtherTok{\textless{}{-}}\NormalTok{ temp\_data}
    \FunctionTok{return}\NormalTok{(output\_data)}
\NormalTok{\}}

\NormalTok{data\_preparation\_binned\_continuous }\OtherTok{\textless{}{-}} \ControlFlowTok{function}\NormalTok{(input\_data) \{}
    \ControlFlowTok{if}\NormalTok{ (}\SpecialCharTok{!}\FunctionTok{exists}\NormalTok{(}\StringTok{"result\_named\_list"}\NormalTok{)) \{}\FunctionTok{stop}\NormalTok{(}\StringTok{"result\_named\_list object is missing"}\NormalTok{)\}}
    \ControlFlowTok{if}\NormalTok{ (}\SpecialCharTok{!}\FunctionTok{exists}\NormalTok{(}\StringTok{"n\_bins\_named\_list"}\NormalTok{)) \{}\FunctionTok{stop}\NormalTok{(}\StringTok{"n\_bins\_named\_list object is missing"}\NormalTok{)\}}
    \ControlFlowTok{if}\NormalTok{ (}\SpecialCharTok{!}\FunctionTok{exists}\NormalTok{(}\StringTok{"binned\_continuous\_features"}\NormalTok{)) \{}\FunctionTok{stop}\NormalTok{(}\StringTok{"binned\_continuous\_features object is missing"}\NormalTok{)\}}
\NormalTok{    temp\_data }\OtherTok{\textless{}{-}}\NormalTok{ input\_data}
    \ControlFlowTok{for}\NormalTok{ (col }\ControlFlowTok{in}\NormalTok{ binned\_continuous\_features) \{}
\NormalTok{        base\_col }\OtherTok{\textless{}{-}}\NormalTok{ col }\SpecialCharTok{\%\textgreater{}\%} \FunctionTok{gsub}\NormalTok{(}\StringTok{"\_binned"}\NormalTok{, }\StringTok{""}\NormalTok{, .) }\SpecialCharTok{\%\textgreater{}\%} \FunctionTok{gsub}\NormalTok{(}\StringTok{"\_merged"}\NormalTok{, }\StringTok{""}\NormalTok{, .)}
        \ControlFlowTok{if}\NormalTok{ (}\SpecialCharTok{!}\NormalTok{base\_col }\SpecialCharTok{\%in\%} \FunctionTok{colnames}\NormalTok{(temp\_data)) \{}\FunctionTok{stop}\NormalTok{(}\FunctionTok{sprintf}\NormalTok{(}\StringTok{"Column \%s is missing from input\_data"}\NormalTok{, base\_col))\}}
\NormalTok{        result }\OtherTok{\textless{}{-}}\NormalTok{ result\_named\_list[[base\_col]]}
\NormalTok{        n\_bins }\OtherTok{\textless{}{-}}\NormalTok{ n\_bins\_named\_list[[base\_col]]}
\NormalTok{        bins }\OtherTok{\textless{}{-}} \FunctionTok{rebuild\_bins}\NormalTok{(result, n\_bins)}
\NormalTok{        temp\_data[[}\FunctionTok{paste0}\NormalTok{(base\_col, }\StringTok{"\_binned"}\NormalTok{)]] }\OtherTok{\textless{}{-}} \FunctionTok{cut}\NormalTok{(temp\_data[[base\_col]], }\AttributeTok{breaks =}\NormalTok{ bins, }\AttributeTok{ordered\_result =} \ConstantTok{TRUE}\NormalTok{, }\AttributeTok{include.lowest =} \ConstantTok{TRUE}\NormalTok{)}
        \ControlFlowTok{if}\NormalTok{ (}\FunctionTok{length}\NormalTok{(}\FunctionTok{unique}\NormalTok{(temp\_data[[}\FunctionTok{paste0}\NormalTok{(base\_col, }\StringTok{"\_binned"}\NormalTok{)]])) }\SpecialCharTok{!=}\NormalTok{ n\_bins) \{}\FunctionTok{stop}\NormalTok{(}\FunctionTok{sprintf}\NormalTok{(}\StringTok{"Binning mismatch for column \%s"}\NormalTok{, base\_col))\}}
\NormalTok{    \}}
\NormalTok{    temp\_data}\SpecialCharTok{$}\NormalTok{AGEPH\_binned\_merged }\OtherTok{\textless{}{-}}\NormalTok{ temp\_data}\SpecialCharTok{$}\NormalTok{AGEPH\_binned }\SpecialCharTok{\%\textgreater{}\%} \FunctionTok{fct\_collapse}\NormalTok{(}
        \StringTok{"(66,95]"} \OtherTok{=} \FunctionTok{c}\NormalTok{(}\StringTok{"(66,68]"}\NormalTok{, }\StringTok{"(68,71]"}\NormalTok{, }\StringTok{"(71,95]"}\NormalTok{))}
\NormalTok{    temp\_data}\SpecialCharTok{$}\NormalTok{POWER\_binned\_merged }\OtherTok{\textless{}{-}}\NormalTok{ temp\_data}\SpecialCharTok{$}\NormalTok{POWER\_binned }\SpecialCharTok{\%\textgreater{}\%} \FunctionTok{fct\_collapse}\NormalTok{(}
        \StringTok{"(76,243]"} \OtherTok{=} \FunctionTok{c}\NormalTok{(}\StringTok{"(76,83]"}\NormalTok{, }\StringTok{"(83,99]"}\NormalTok{, }\StringTok{"(99,243]"}\NormalTok{)}
\NormalTok{    )}
    \ControlFlowTok{for}\NormalTok{ (col }\ControlFlowTok{in}\NormalTok{ binned\_continuous\_features) \{}
\NormalTok{        temp\_data[[col]] }\OtherTok{\textless{}{-}} \FunctionTok{factor}\NormalTok{(temp\_data[[col]], }\AttributeTok{levels =}\NormalTok{ levels\_named\_list[[col]], }\AttributeTok{ordered =} \ConstantTok{FALSE}\NormalTok{)}
\NormalTok{    \}}
\NormalTok{    output\_data }\OtherTok{\textless{}{-}}\NormalTok{ temp\_data}
    \FunctionTok{return}\NormalTok{(output\_data)}
\NormalTok{\}}

\NormalTok{data\_preparation\_enriched }\OtherTok{\textless{}{-}} \ControlFlowTok{function}\NormalTok{(input\_data) \{}
\NormalTok{    temp\_data }\OtherTok{\textless{}{-}}\NormalTok{ input\_data}
    \ControlFlowTok{if}\NormalTok{ (}\SpecialCharTok{!}\StringTok{"POPULATION"} \SpecialCharTok{\%in\%} \FunctionTok{colnames}\NormalTok{(temp\_data)) \{}\FunctionTok{stop}\NormalTok{(}\StringTok{"Column POPULATION is missing from input\_data"}\NormalTok{)\}}
    \ControlFlowTok{if}\NormalTok{ (}\SpecialCharTok{!}\StringTok{"REGION"} \SpecialCharTok{\%in\%} \FunctionTok{colnames}\NormalTok{(temp\_data)) \{}\FunctionTok{stop}\NormalTok{(}\StringTok{"Column REGION is missing from input\_data"}\NormalTok{)\}}
    \ControlFlowTok{if}\NormalTok{ (}\SpecialCharTok{!}\StringTok{"MUNICIPALITY"} \SpecialCharTok{\%in\%} \FunctionTok{colnames}\NormalTok{(temp\_data)) \{}\FunctionTok{stop}\NormalTok{(}\StringTok{"Column MUNICIPALITY is missing from input\_data"}\NormalTok{)\}}
\NormalTok{    temp\_data}\SpecialCharTok{$}\NormalTok{POPULATION\_binned }\OtherTok{\textless{}{-}} \FunctionTok{cut}\NormalTok{(temp\_data}\SpecialCharTok{$}\NormalTok{POPULATION, }\AttributeTok{breaks =} \FunctionTok{rebuild\_bins}\NormalTok{(}
\NormalTok{        result\_named\_list[[}\StringTok{"POPULATION"}\NormalTok{]], n\_bins\_named\_list[[}\StringTok{"POPULATION"}\NormalTok{]]}
\NormalTok{    ), }\AttributeTok{ordered\_result =} \ConstantTok{TRUE}\NormalTok{, }\AttributeTok{include.lowest =} \ConstantTok{TRUE}\NormalTok{) }\SpecialCharTok{\%\textgreater{}\%}
    \FunctionTok{fct\_collapse}\NormalTok{(}
        \StringTok{"(1.34e+05,4.53e+05]"} \OtherTok{=} \FunctionTok{c}\NormalTok{(}\StringTok{"(1.34e+05,2.05e+05]"}\NormalTok{, }\StringTok{"(2.05e+05,4.53e+05]"}\NormalTok{)}
\NormalTok{    )}
\NormalTok{    temp\_data}\SpecialCharTok{$}\NormalTok{REGION\_merged }\OtherTok{\textless{}{-}}\NormalTok{ temp\_data}\SpecialCharTok{$}\NormalTok{REGION }\SpecialCharTok{\%\textgreater{}\%} \FunctionTok{as.factor}\NormalTok{() }\SpecialCharTok{\%\textgreater{}\%}
        \FunctionTok{fct\_collapse}\NormalTok{(}
            \StringTok{"Brussels"} \OtherTok{=} \FunctionTok{c}\NormalTok{(}\StringTok{"RÉGION DE BRUXELLES{-}CAPITALE"}\NormalTok{),}
            \StringTok{"Others"} \OtherTok{=} \FunctionTok{c}\NormalTok{(}\StringTok{"RÉGION WALLONNE"}\NormalTok{, }\StringTok{"RÉGION FLAMANDE"}\NormalTok{)}
\NormalTok{        )}
\NormalTok{    temp\_data}\SpecialCharTok{$}\NormalTok{MUNICIPALITY\_binned }\OtherTok{\textless{}{-}} \FunctionTok{ifelse}\NormalTok{(temp\_data}\SpecialCharTok{$}\NormalTok{MUNICIPALITY }\SpecialCharTok{\%in\%}\NormalTok{ levels\_named\_list[[}\StringTok{"MUNICIPALITY\_binned"}\NormalTok{]], temp\_data}\SpecialCharTok{$}\NormalTok{MUNICIPALITY, }\StringTok{"Other"}\NormalTok{) }\SpecialCharTok{\%\textgreater{}\%}
        \FunctionTok{as.factor}\NormalTok{()}
    \ControlFlowTok{for}\NormalTok{ (col }\ControlFlowTok{in}\NormalTok{ binned\_enriched\_features) \{}
\NormalTok{        temp\_data[[col]] }\OtherTok{\textless{}{-}} \FunctionTok{factor}\NormalTok{(temp\_data[[col]], }\AttributeTok{levels =}\NormalTok{ levels\_named\_list[[col]], }\AttributeTok{ordered =} \ConstantTok{FALSE}\NormalTok{)}
\NormalTok{    \}}
\NormalTok{    output\_data }\OtherTok{\textless{}{-}}\NormalTok{ temp\_data}
    \FunctionTok{return}\NormalTok{(output\_data)}
\NormalTok{\}}

\NormalTok{data\_preparation\_interaction }\OtherTok{\textless{}{-}} \ControlFlowTok{function}\NormalTok{(input\_data) \{}
    \ControlFlowTok{if}\NormalTok{ (}\SpecialCharTok{!}\StringTok{"BM\_binned"} \SpecialCharTok{\%in\%} \FunctionTok{colnames}\NormalTok{(input\_data)) \{}\FunctionTok{stop}\NormalTok{(}\StringTok{"Column BM\_binned is missing from input\_data"}\NormalTok{)\}}
    \ControlFlowTok{if}\NormalTok{ (}\SpecialCharTok{!}\StringTok{"FUEL"} \SpecialCharTok{\%in\%} \FunctionTok{colnames}\NormalTok{(input\_data)) \{}\FunctionTok{stop}\NormalTok{(}\StringTok{"Column FUEL is missing from input\_data"}\NormalTok{)\}}
\NormalTok{    temp\_data }\OtherTok{\textless{}{-}}\NormalTok{ input\_data }\SpecialCharTok{\%\textgreater{}\%}
        \FunctionTok{mutate}\NormalTok{(}
            \AttributeTok{BM\_01\_FUEL\_gasoline =}\NormalTok{ BM\_binned }\SpecialCharTok{==} \StringTok{"[0,1]"} \SpecialCharTok{\&}\NormalTok{ FUEL }\SpecialCharTok{==} \StringTok{"gasoline"}\NormalTok{,}
            \AttributeTok{BM\_01\_FUEL\_not\_gasoline =}\NormalTok{ BM\_binned }\SpecialCharTok{==} \StringTok{"[0,1]"} \SpecialCharTok{\&}\NormalTok{ FUEL }\SpecialCharTok{!=} \StringTok{"gasoline"}
\NormalTok{        )}
\NormalTok{    output\_data }\OtherTok{\textless{}{-}}\NormalTok{ temp\_data}
    \FunctionTok{return}\NormalTok{(output\_data)}
\NormalTok{\}}

\NormalTok{data\_preparation\_workflow }\OtherTok{\textless{}{-}} \ControlFlowTok{function}\NormalTok{(input\_data) \{}
\NormalTok{    temp\_data }\OtherTok{\textless{}{-}}\NormalTok{ input\_data }\SpecialCharTok{\%\textgreater{}\%}
\NormalTok{        data\_enrichment }\SpecialCharTok{\%\textgreater{}\%}
\NormalTok{        data\_preparation\_original\_categorical }\SpecialCharTok{\%\textgreater{}\%}
\NormalTok{        data\_preparation\_binned\_continuous }\SpecialCharTok{\%\textgreater{}\%}
\NormalTok{        data\_preparation\_enriched }\SpecialCharTok{\%\textgreater{}\%}
\NormalTok{        data\_preparation\_interaction}
\NormalTok{    output\_data }\OtherTok{\textless{}{-}}\NormalTok{ temp\_data}
    \FunctionTok{return}\NormalTok{(output\_data)}
\NormalTok{\}}
\end{Highlighting}
\end{Shaded}

\begin{Shaded}
\begin{Highlighting}[]
\NormalTok{test\_data }\OtherTok{\textless{}{-}}\NormalTok{ mtpl\_be }\SpecialCharTok{\%\textgreater{}\%}
    \FunctionTok{filter}\NormalTok{(ID }\SpecialCharTok{\%in\%}\NormalTok{ data}\SpecialCharTok{$}\NormalTok{ID) }\SpecialCharTok{\%\textgreater{}\%}
    \FunctionTok{data\_preparation\_workflow}\NormalTok{()}

\CommentTok{\# Set ordered factors for binned continuous features to unordered}
\ControlFlowTok{for}\NormalTok{ (col }\ControlFlowTok{in} \FunctionTok{names}\NormalTok{(levels\_named\_list)) \{}
\NormalTok{    data\_enriched[[col]] }\OtherTok{\textless{}{-}} \FunctionTok{factor}\NormalTok{(data\_enriched[[col]], }\AttributeTok{levels =}\NormalTok{ levels\_named\_list[[col]], }\AttributeTok{ordered =} \ConstantTok{FALSE}\NormalTok{)}
\NormalTok{\}}

\ControlFlowTok{for}\NormalTok{ (col }\ControlFlowTok{in} \FunctionTok{c}\NormalTok{(original\_categorical\_features, binned\_continuous\_features, enriched\_features)) \{}
    \ControlFlowTok{if}\NormalTok{ (}\SpecialCharTok{!}\NormalTok{col }\SpecialCharTok{\%in\%} \FunctionTok{colnames}\NormalTok{(test\_data)) \{}
        \FunctionTok{stop}\NormalTok{(}\FunctionTok{paste}\NormalTok{(}\StringTok{"Column"}\NormalTok{, col, }\StringTok{"is missing from test\_data"}\NormalTok{))}
\NormalTok{    \}}
    \ControlFlowTok{if}\NormalTok{ (}\SpecialCharTok{!}\FunctionTok{identical}\NormalTok{(test\_data[[col]], data\_enriched[[col]])) \{}
        \FunctionTok{stop}\NormalTok{(}\FunctionTok{paste}\NormalTok{(}\StringTok{"Column"}\NormalTok{, col, }\StringTok{"does not match the enriched data"}\NormalTok{))}
\NormalTok{    \}}
\NormalTok{\}}

\FunctionTok{cat}\NormalTok{(}\FunctionTok{sprintf}\NormalTok{(}\StringTok{"All columns in test\_data match the enriched data.}\SpecialCharTok{\textbackslash{}n}\StringTok{(\%s)"}\NormalTok{, }\FunctionTok{paste}\NormalTok{(}\FunctionTok{c}\NormalTok{(original\_categorical\_features, binned\_continuous\_features, enriched\_features), }\AttributeTok{collapse =} \StringTok{", "}\NormalTok{)))}
\end{Highlighting}
\end{Shaded}

\begin{verbatim}
All columns in test_data match the enriched data.
(COVERAGE_BIN, FUEL, FLEET, SEX, BM_binned, AGEPH_binned_merged, AGEC_binned, POWER_binned_merged, REGION_merged, MUNICIPALITY_binned, POPULATION)
\end{verbatim}

\subsubsection{Cross-Validation}\label{cross-validation}

To better understand cross-validation, let's illustrate it with a simple
manual example. We create 5 data sets (folds) for cross-validation. Each
fold will be used once as a validation set while the remaining folds
form the training set.

\begin{Shaded}
\begin{Highlighting}[]
\NormalTok{cv\_split }\OtherTok{\textless{}{-}} \ControlFlowTok{function}\NormalTok{(input\_data, kfold, strata) \{}
\NormalTok{    folds }\OtherTok{\textless{}{-}}\NormalTok{ input\_data }\SpecialCharTok{\%\textgreater{}\%} \FunctionTok{vfold\_cv}\NormalTok{(}\AttributeTok{strata =}\NormalTok{ strata, }\AttributeTok{v =}\NormalTok{ kfold)}
    \FunctionTok{return}\NormalTok{(folds)}
\NormalTok{\}}
\NormalTok{base\_formula }\OtherTok{\textless{}{-}} \StringTok{"NCLAIMS \textasciitilde{} offset(log(EXP))"}
\NormalTok{folds }\OtherTok{\textless{}{-}}\NormalTok{ mtpl\_be }\SpecialCharTok{\%\textgreater{}\%}
    \FunctionTok{filter}\NormalTok{(ID }\SpecialCharTok{\%in\%}\NormalTok{ data}\SpecialCharTok{$}\NormalTok{ID) }\SpecialCharTok{\%\textgreater{}\%}
    \FunctionTok{data\_preparation\_workflow}\NormalTok{() }\SpecialCharTok{\%\textgreater{}\%}
    \FunctionTok{cv\_split}\NormalTok{(}\AttributeTok{kfold =} \DecValTok{5}\NormalTok{, }\AttributeTok{strata =} \StringTok{"NCLAIMS"}\NormalTok{)}
\NormalTok{models }\OtherTok{\textless{}{-}} \FunctionTok{vector}\NormalTok{(}\StringTok{"list"}\NormalTok{, }\FunctionTok{length}\NormalTok{(folds}\SpecialCharTok{$}\NormalTok{splits))}
\NormalTok{poisson\_log\_loss\_list }\OtherTok{\textless{}{-}} \FunctionTok{list}\NormalTok{()}
\NormalTok{beta0 }\OtherTok{\textless{}{-}} \FunctionTok{list}\NormalTok{()}
\ControlFlowTok{for}\NormalTok{ (i }\ControlFlowTok{in} \FunctionTok{seq\_along}\NormalTok{(folds}\SpecialCharTok{$}\NormalTok{splits)) \{}
\NormalTok{    train\_data }\OtherTok{\textless{}{-}} \FunctionTok{analysis}\NormalTok{(folds}\SpecialCharTok{$}\NormalTok{splits[[i]])}
\NormalTok{    val\_data }\OtherTok{\textless{}{-}} \FunctionTok{assessment}\NormalTok{(folds}\SpecialCharTok{$}\NormalTok{splits[[i]])}
\NormalTok{    models[[i]] }\OtherTok{\textless{}{-}} \FunctionTok{glm}\NormalTok{(}\AttributeTok{formula =} \FunctionTok{as.formula}\NormalTok{(base\_formula), }\AttributeTok{data =}\NormalTok{ train\_data, }\AttributeTok{family =} \FunctionTok{poisson}\NormalTok{(}\AttributeTok{link =} \StringTok{"log"}\NormalTok{))}
\NormalTok{    log\_loss }\OtherTok{\textless{}{-}} \SpecialCharTok{{-}}\FunctionTok{mean}\NormalTok{(}\FunctionTok{dpois}\NormalTok{(val\_data}\SpecialCharTok{$}\NormalTok{NCLAIMS, }\AttributeTok{lambda =} \FunctionTok{predict}\NormalTok{(models[[i]], }\AttributeTok{newdata =}\NormalTok{ val\_data, }\AttributeTok{type =} \StringTok{"response"}\NormalTok{), }\AttributeTok{log =} \ConstantTok{TRUE}\NormalTok{))}
\NormalTok{    poisson\_log\_loss\_list[[i]] }\OtherTok{\textless{}{-}}\NormalTok{ log\_loss}
\NormalTok{    beta0[[i]] }\OtherTok{\textless{}{-}} \FunctionTok{coef}\NormalTok{(models[[i]])[}\DecValTok{1}\NormalTok{]}
\NormalTok{\}}
\NormalTok{m0 }\OtherTok{\textless{}{-}} \FunctionTok{glm}\NormalTok{(}\AttributeTok{formula =} \FunctionTok{as.formula}\NormalTok{(base\_formula), }\AttributeTok{data =}\NormalTok{ mtpl\_be }\SpecialCharTok{\%\textgreater{}\%} \FunctionTok{filter}\NormalTok{(ID }\SpecialCharTok{\%in\%}\NormalTok{ data}\SpecialCharTok{$}\NormalTok{ID) }\SpecialCharTok{\%\textgreater{}\%} \FunctionTok{data\_preparation\_workflow}\NormalTok{(), }\AttributeTok{family =} \FunctionTok{poisson}\NormalTok{(}\AttributeTok{link =} \StringTok{"log"}\NormalTok{))}
\NormalTok{beta0\_full }\OtherTok{\textless{}{-}} \FunctionTok{coef}\NormalTok{(m0)[}\DecValTok{1}\NormalTok{]}
\NormalTok{log\_loss\_full }\OtherTok{\textless{}{-}} \SpecialCharTok{{-}}\FunctionTok{mean}\NormalTok{(}\FunctionTok{dpois}\NormalTok{(holdout}\SpecialCharTok{$}\NormalTok{NCLAIMS, }\AttributeTok{lambda =} \FunctionTok{predict}\NormalTok{(m0, }\AttributeTok{newdata =}\NormalTok{ holdout, }\AttributeTok{type =} \StringTok{"response"}\NormalTok{), }\AttributeTok{log =} \ConstantTok{TRUE}\NormalTok{))}
\FunctionTok{cat}\NormalTok{(}\FunctionTok{sprintf}\NormalTok{(}\StringTok{"For the simple GLM model:}\SpecialCharTok{\textbackslash{}n}\StringTok{  Log loss ({-}log{-}likelihood/n): \%s}\SpecialCharTok{\textbackslash{}n}\StringTok{  Intercept (beta0): \%s}\SpecialCharTok{\textbackslash{}n}\StringTok{  95\%\% CI for beta0: [\%s, \%s]}\SpecialCharTok{\textbackslash{}n}\StringTok{"}\NormalTok{, log\_loss\_full, beta0\_full, beta0\_full }\SpecialCharTok{{-}} \FloatTok{1.96} \SpecialCharTok{*} \FunctionTok{summary}\NormalTok{(m0)}\SpecialCharTok{$}\NormalTok{coefficients[}\DecValTok{1}\NormalTok{,}\DecValTok{2}\NormalTok{], beta0\_full }\SpecialCharTok{+} \FloatTok{1.96} \SpecialCharTok{*} \FunctionTok{summary}\NormalTok{(m0)}\SpecialCharTok{$}\NormalTok{coefficients[}\DecValTok{1}\NormalTok{,}\DecValTok{2}\NormalTok{]))}
\end{Highlighting}
\end{Shaded}

\begin{verbatim}
For the simple GLM model:
  Log loss (-log-likelihood/n): 0.397287337307078
  Intercept (beta0): -1.97599003410665
  95% CI for beta0: [-1.9914326792921, -1.96054738892119]
\end{verbatim}

\begin{Shaded}
\begin{Highlighting}[]
\FunctionTok{cat}\NormalTok{(}\FunctionTok{sprintf}\NormalTok{(}\StringTok{"Cross{-}validated log loss: min = \%s, mean = \%s, max = \%s}\SpecialCharTok{\textbackslash{}n}\StringTok{"}\NormalTok{, }\FunctionTok{min}\NormalTok{(}\FunctionTok{unlist}\NormalTok{(poisson\_log\_loss\_list)), }\FunctionTok{mean}\NormalTok{(}\FunctionTok{unlist}\NormalTok{(poisson\_log\_loss\_list)), }\FunctionTok{max}\NormalTok{(}\FunctionTok{unlist}\NormalTok{(poisson\_log\_loss\_list))))}
\end{Highlighting}
\end{Shaded}

\begin{verbatim}
Cross-validated log loss: min = 0.384639629129922, mean = 0.389544236202061, max = 0.39372190147375
\end{verbatim}

\begin{Shaded}
\begin{Highlighting}[]
\FunctionTok{cat}\NormalTok{(}\FunctionTok{sprintf}\NormalTok{(}\StringTok{"Cross{-}validated beta0: min = \%s, mean = \%s, max = \%s}\SpecialCharTok{\textbackslash{}n}\StringTok{"}\NormalTok{, }\FunctionTok{min}\NormalTok{(}\FunctionTok{unlist}\NormalTok{(beta0)), }\FunctionTok{mean}\NormalTok{(}\FunctionTok{unlist}\NormalTok{(beta0)), }\FunctionTok{max}\NormalTok{(}\FunctionTok{unlist}\NormalTok{(beta0))))}
\end{Highlighting}
\end{Shaded}

\begin{verbatim}
Cross-validated beta0: min = -1.9793873830536, mean = -1.97599273053119, max = -1.97336821528614
\end{verbatim}

This illustrates how cross-validation enables enables to check the
robustness and variability of the model's performance and parameter
estimates across different subsets of the data. Comparing the
cross-validated mean log loss, AIC or BIC across various models will
yield a more robust choice for model selection. Now that the intuition
is set, we can create our workflow including cross-validation.

\subsubsection{Performance Metrics}\label{performance-metrics}

To compare the performance of different models, we can consider the
following cross-validated metrics:

\begin{itemize}
\tightlist
\item
  Log Loss: Measures predictive accuracy by evaluating the negative
  log-likelihood (scaled by the number of observations) of the observed
  outcomes given the predicted probabilities.
\item
  AIC (Akaike Information Criterion): Evaluates model quality by
  balancing goodness of fit and model complexity. Lower values indicate
  better models.
\item
  BIC (Bayesian Information Criterion): Similar to AIC, but with a
  stronger penalty for model complexity. Lower values indicate better
  models.
\end{itemize}

Note that the \texttt{Log-\ Loss}, \texttt{AIC} and \texttt{BIC} are all
based on the likelihood of the model, here are how they are computed:

\begin{itemize}
\tightlist
\item
  Log Loss:
  \texttt{-mean(dpois(y,\ lambda\ =\ predicted\_lambda,\ log\ =\ TRUE))}
  or \texttt{-logLik(model)/n} where \texttt{n} is the number of
  observations.
\item
  AIC: \texttt{2\ *\ k\ -\ 2\ *\ logLik(model)} where \texttt{k} is the
  number of parameters.
\item
  BIC: \texttt{log(n)\ *\ k\ -\ 2\ *\ logLik(model)} where \texttt{n} is
  the number of observations and \texttt{k} is the number of parameters.
\end{itemize}

We see that \texttt{Log\ loss}simply does not account for model
complexity, unlike \texttt{AIC} and \texttt{BIC} which include penalties
of 2 * k and 11.5566361\footnote{The term \texttt{log(nrow(data)*(4/5))}
  comes from the fact that we are using cross-validation with 5 folds,
  so each training set contains 4/5 of the data.} * k respectively.
Because the penalty is so much higher, due to the large size of our
portfolio, \texttt{BIC} tends to favor simpler models compared to
\texttt{AIC}.

The models will be gradually complexified by adding covariates. For each
step, we will compare the final model with the previous step based on:

\begin{itemize}
\tightlist
\item
  Scaled Residual Deviance: Measures the goodness of fit of a model
  relative to a saturated model, scaled by the residual deviance of the
  null model.
\item
  Anova Test: Compares nested models by evaluating whether the addition
  of new covariates significantly improves the model fit.
\end{itemize}

For the scaled residual deviance, here are the maximum and minimum
deviance values:

\begin{itemize}
\tightlist
\item
  Null Deviance:
  \(D_0 = 2 \times \left( logLik(SaturatedModel) - logLik(NullModel) \right)\)
\item
  Residual Deviance:
  \(D_m = 2 \times \left( logLik(SaturatedModel) - logLik(FittedModel) \right)\)
\item
  Scaled Residual Deviance: \(D_s = \frac{D_m}{D_0}\)
\end{itemize}

\begin{Shaded}
\begin{Highlighting}[]
\NormalTok{train\_data }\OtherTok{\textless{}{-}}\NormalTok{ mtpl\_be }\SpecialCharTok{\%\textgreater{}\%}
    \FunctionTok{filter}\NormalTok{(ID }\SpecialCharTok{\%in\%}\NormalTok{ data}\SpecialCharTok{$}\NormalTok{ID) }\SpecialCharTok{\%\textgreater{}\%}
    \FunctionTok{data\_preparation\_workflow}\NormalTok{()}
\NormalTok{m\_null }\OtherTok{\textless{}{-}} \FunctionTok{glm}\NormalTok{(NCLAIMS }\SpecialCharTok{\textasciitilde{}} \FunctionTok{offset}\NormalTok{(}\FunctionTok{log}\NormalTok{(EXP)), }\AttributeTok{family =}\NormalTok{ poisson, }\AttributeTok{data =}\NormalTok{ train\_data)}
\NormalTok{m\_ex }\OtherTok{\textless{}{-}} \FunctionTok{glm}\NormalTok{(NCLAIMS }\SpecialCharTok{\textasciitilde{}} \FunctionTok{offset}\NormalTok{(}\FunctionTok{log}\NormalTok{(EXP)) }\SpecialCharTok{+}\NormalTok{ BM\_binned, }\AttributeTok{family =}\NormalTok{ poisson, }\AttributeTok{data =}\NormalTok{ train\_data)}
\FunctionTok{cat}\NormalTok{(}\FunctionTok{sprintf}\NormalTok{(}
    \StringTok{"Scaled Residual Deviance: \%.2f\%\% (\%.0f / \%.0f)"}\NormalTok{, }
    \FunctionTok{deviance}\NormalTok{(m\_ex) }\SpecialCharTok{/} \FunctionTok{deviance}\NormalTok{(m\_null)}\SpecialCharTok{*}\DecValTok{100}\NormalTok{,}
    \FunctionTok{deviance}\NormalTok{(m\_ex),}
    \FunctionTok{deviance}\NormalTok{(m\_null)}
\NormalTok{))}
\end{Highlighting}
\end{Shaded}

\begin{verbatim}
Scaled Residual Deviance: 98.28% (70460 / 71692)
\end{verbatim}

\begin{Shaded}
\begin{Highlighting}[]
\NormalTok{ll\_sat }\OtherTok{\textless{}{-}} \FunctionTok{deviance}\NormalTok{(m\_null)}\SpecialCharTok{/}\DecValTok{2} \SpecialCharTok{+} \FunctionTok{logLik}\NormalTok{(m\_null)}
\NormalTok{calc\_dm }\OtherTok{\textless{}{-}} \DecValTok{2} \SpecialCharTok{*}\NormalTok{ (ll\_sat }\SpecialCharTok{{-}} \FunctionTok{logLik}\NormalTok{(m\_ex))}
\FunctionTok{cat}\NormalTok{(}\FunctionTok{sprintf}\NormalTok{(}
    \StringTok{"Calculated D\_m = \%.2f and Model Deviance = \%.0f"}\NormalTok{, }
\NormalTok{    calc\_dm,}
    \FunctionTok{deviance}\NormalTok{(m\_ex)}
\NormalTok{))}
\end{Highlighting}
\end{Shaded}

\begin{verbatim}
Calculated D_m = 70460.30 and Model Deviance = 70460
\end{verbatim}

\begin{Shaded}
\begin{Highlighting}[]
\FunctionTok{cat}\NormalTok{(}\FunctionTok{sprintf}\NormalTok{(}\StringTok{"We thus have maximum log likelihood for the saturated model: \%.2f"}\NormalTok{, ll\_sat))}
\end{Highlighting}
\end{Shaded}

\begin{verbatim}
We thus have maximum log likelihood for the saturated model: -15022.23
\end{verbatim}

\begin{Shaded}
\begin{Highlighting}[]
\FunctionTok{cat}\NormalTok{(}\FunctionTok{sprintf}\NormalTok{(}\StringTok{"And minimum log likelihood for the null model: \%.2f"}\NormalTok{, }\FunctionTok{logLik}\NormalTok{(m\_null)))}
\end{Highlighting}
\end{Shaded}

\begin{verbatim}
And minimum log likelihood for the null model: -50868.24
\end{verbatim}

The anova test will be comparing the deviance of the nested models to
evaluate whether the addition of new covariates significantly improves
the model fit.

\begin{Shaded}
\begin{Highlighting}[]
\FunctionTok{anova}\NormalTok{(m\_null, m\_ex, }\AttributeTok{test =} \StringTok{"Chisq"}\NormalTok{)}
\end{Highlighting}
\end{Shaded}

\begin{verbatim}
Analysis of Deviance Table

Model 1: NCLAIMS ~ offset(log(EXP))
Model 2: NCLAIMS ~ offset(log(EXP)) + BM_binned
  Resid. Df Resid. Dev Df Deviance  Pr(>Chi)    
1    130584      71692                          
2    130582      70460  2   1231.7 < 2.2e-16 ***
---
Signif. codes:  0 '***' 0.001 '**' 0.01 '*' 0.05 '.' 0.1 ' ' 1
\end{verbatim}

\subsection{Frequency Modelling}\label{frequency-modelling}

Based on all the notions explained in the Validation Framework: data
preparation, cross-validation, and performance metrics, we will now
create workflows for frequency modelling.

\subsubsection{GLM Models}\label{glm-models}

\begin{Shaded}
\begin{Highlighting}[]
\NormalTok{create\_wflow\_freq\_glm }\OtherTok{\textless{}{-}} \ControlFlowTok{function}\NormalTok{(data, formula, }\AttributeTok{v =} \DecValTok{5}\NormalTok{, }\AttributeTok{seed =} \DecValTok{123}\NormalTok{) \{}

    \FunctionTok{set.seed}\NormalTok{(seed)}
\NormalTok{    folds }\OtherTok{\textless{}{-}} \FunctionTok{vfold\_cv}\NormalTok{(}
\NormalTok{        data, }\AttributeTok{v =}\NormalTok{ v, }
        \AttributeTok{strata =}\NormalTok{ NCLAIMS}
\NormalTok{    )}

\NormalTok{    rec }\OtherTok{\textless{}{-}} \FunctionTok{recipe}\NormalTok{(}\FunctionTok{as.formula}\NormalTok{(formula), }\AttributeTok{data =}\NormalTok{ data)}

    \CommentTok{\# The Poisson family uses a log link by default}
\NormalTok{    mod }\OtherTok{\textless{}{-}} \FunctionTok{poisson\_reg}\NormalTok{() }\SpecialCharTok{\%\textgreater{}\%}
        \FunctionTok{set\_engine}\NormalTok{(}\StringTok{"glm"}\NormalTok{)}

\NormalTok{    wflow }\OtherTok{\textless{}{-}} \FunctionTok{workflow}\NormalTok{() }\SpecialCharTok{\%\textgreater{}\%}
        \FunctionTok{add\_recipe}\NormalTok{(rec) }\SpecialCharTok{\%\textgreater{}\%}
        \FunctionTok{add\_model}\NormalTok{(mod)}

\NormalTok{    cv\_results }\OtherTok{\textless{}{-}} \FunctionTok{fit\_resamples}\NormalTok{(}
\NormalTok{        wflow,}
        \AttributeTok{resamples =}\NormalTok{ folds,}
        \AttributeTok{control =} \FunctionTok{control\_resamples}\NormalTok{(}\AttributeTok{save\_pred =} \ConstantTok{TRUE}\NormalTok{)}
\NormalTok{    )}

    \CommentTok{\# Cross{-}validation does not fit a final model on all the data}
\NormalTok{    final\_fit }\OtherTok{\textless{}{-}} \FunctionTok{fit}\NormalTok{(wflow, }\AttributeTok{data =}\NormalTok{ data)}
\NormalTok{    glm\_fit }\OtherTok{\textless{}{-}} \FunctionTok{extract\_fit\_engine}\NormalTok{(final\_fit)}

    \CommentTok{\# Cross{-}validated deviance}
\NormalTok{    cv\_pred }\OtherTok{\textless{}{-}} \FunctionTok{collect\_predictions}\NormalTok{(cv\_results)}
\NormalTok{    cv\_deviance }\OtherTok{\textless{}{-}}\NormalTok{ cv\_pred }\SpecialCharTok{\%\textgreater{}\%}
        \FunctionTok{mutate}\NormalTok{(}\AttributeTok{deviance\_i =} \FunctionTok{poisson}\NormalTok{()}\SpecialCharTok{$}\FunctionTok{dev.resids}\NormalTok{(}\AttributeTok{y =}\NormalTok{ NCLAIMS, }\AttributeTok{mu =}\NormalTok{ .pred, }\AttributeTok{wt =} \DecValTok{1}\NormalTok{)) }\SpecialCharTok{\%\textgreater{}\%}
        \FunctionTok{summarise}\NormalTok{(}
            \AttributeTok{total\_deviance =} \FunctionTok{sum}\NormalTok{(deviance\_i),}
\NormalTok{        )}
\NormalTok{    cv\_log\_loss }\OtherTok{\textless{}{-}} \FunctionTok{poisson\_log\_loss}\NormalTok{(cv\_pred, }\AttributeTok{truth =}\NormalTok{ NCLAIMS, }\AttributeTok{estimate =}\NormalTok{ .pred)}

    \FunctionTok{list}\NormalTok{(}
        \AttributeTok{workflow =}\NormalTok{ wflow,}
        \AttributeTok{folds =}\NormalTok{ folds,}
        \AttributeTok{resamples =}\NormalTok{ cv\_results,}
        \AttributeTok{fit =}\NormalTok{ glm\_fit,}
        \AttributeTok{performance\_metrics =} \FunctionTok{list}\NormalTok{(}
            \AttributeTok{AIC =} \FunctionTok{AIC}\NormalTok{(glm\_fit),}
            \AttributeTok{BIC =} \FunctionTok{BIC}\NormalTok{(glm\_fit),}
            \AttributeTok{Deviance =} \FunctionTok{deviance}\NormalTok{(glm\_fit),}
            \AttributeTok{CV\_Deviance =}\NormalTok{ cv\_deviance}\SpecialCharTok{$}\NormalTok{total\_deviance,}
            \AttributeTok{CV\_LogLoss =}\NormalTok{ cv\_log\_loss}\SpecialCharTok{$}\NormalTok{.estimate}
\NormalTok{        )}
\NormalTok{    )}
\NormalTok{\}}
\end{Highlighting}
\end{Shaded}

We created a workflow to:

\begin{itemize}
\tightlist
\item
  Make the n folds (5 by default)
\item
  Add the new levels to a categorical variable ``Other'', particularly
  relevant for Municipalities with low frequency
\item
  Specify the Poisson regression model with a log link
\item
  Fit the model on the entire training dataset
\item
  Perform cross-validation to evaluate model performance
\item
  Store the workflow, cross-validation (Deviance and LogLoss) and final
  model fit
\end{itemize}

This workflow will be used for all models in the frequency analysis,
ensuring consistency in model fitting, cross-validation, and performance
evaluation.

\paragraph{Categorical Variables Only}\label{categorical-variables-only}

Let us start with adding the categorical variables only to our frequency
model:

\begin{itemize}
\tightlist
\item
  \texttt{COVERAGE\_BIN}
\item
  \texttt{FLEET}
\item
  \texttt{SEX}
\item
  \texttt{FUEL}
\end{itemize}

Note that \texttt{USE} was not significant in the preliminary analysis
and thus is not included in the categorical variables for the frequency
model.

\begin{Shaded}
\begin{Highlighting}[]
\NormalTok{perf\_metrics\_df }\OtherTok{\textless{}{-}} \FunctionTok{data.frame}\NormalTok{(}
    \AttributeTok{Model =} \FunctionTok{character}\NormalTok{(),}
    \AttributeTok{AIC =} \FunctionTok{numeric}\NormalTok{(),}
    \AttributeTok{BIC =} \FunctionTok{numeric}\NormalTok{(),}
    \AttributeTok{Deviance =} \FunctionTok{numeric}\NormalTok{(),}
    \AttributeTok{CV\_Deviance =} \FunctionTok{numeric}\NormalTok{(),}
    \AttributeTok{CV\_LogLoss =} \FunctionTok{numeric}\NormalTok{()}
\NormalTok{)}
\NormalTok{append\_performance\_metrics }\OtherTok{\textless{}{-}} \ControlFlowTok{function}\NormalTok{(name, formula, result, }\AttributeTok{perf\_df =}\NormalTok{ perf\_metrics\_df) \{}
\NormalTok{    perf }\OtherTok{\textless{}{-}}\NormalTok{ result}\SpecialCharTok{$}\NormalTok{performance\_metrics}
\NormalTok{    new\_row }\OtherTok{\textless{}{-}} \FunctionTok{data.frame}\NormalTok{(}
        \AttributeTok{Model =}\NormalTok{ name,}
        \AttributeTok{AIC =}\NormalTok{ perf}\SpecialCharTok{$}\NormalTok{AIC,}
        \AttributeTok{BIC =}\NormalTok{ perf}\SpecialCharTok{$}\NormalTok{BIC,}
        \AttributeTok{Deviance =}\NormalTok{ perf}\SpecialCharTok{$}\NormalTok{Deviance,}
        \AttributeTok{CV\_Deviance =}\NormalTok{ perf}\SpecialCharTok{$}\NormalTok{CV\_Deviance,}
        \AttributeTok{CV\_LogLoss =}\NormalTok{ perf}\SpecialCharTok{$}\NormalTok{CV\_LogLoss}
\NormalTok{    )}
    \FunctionTok{rbind}\NormalTok{(perf\_df, new\_row)}
\NormalTok{\}}

\NormalTok{f0 }\OtherTok{\textless{}{-}} \StringTok{"NCLAIMS \textasciitilde{} 1"}
\NormalTok{fcat }\OtherTok{\textless{}{-}} \FunctionTok{paste0}\NormalTok{(}\StringTok{"NCLAIMS \textasciitilde{} "}\NormalTok{, }\FunctionTok{paste}\NormalTok{(original\_categorical\_features, }\AttributeTok{collapse =} \StringTok{" + "}\NormalTok{))}
\FunctionTok{cat}\NormalTok{(}\FunctionTok{sprintf}\NormalTok{(}\StringTok{"\textgreater{} Categorical formula: \%s"}\NormalTok{, fcat))}
\end{Highlighting}
\end{Shaded}

\begin{verbatim}
> Categorical formula: NCLAIMS ~ COVERAGE_BIN + FUEL + FLEET + SEX
\end{verbatim}

\begin{Shaded}
\begin{Highlighting}[]
\NormalTok{r0 }\OtherTok{\textless{}{-}} \FunctionTok{create\_wflow\_freq\_glm}\NormalTok{(train\_data, f0)}
\NormalTok{perf\_metrics\_df }\OtherTok{\textless{}{-}} \FunctionTok{append\_performance\_metrics}\NormalTok{(}\AttributeTok{name =} \StringTok{"Intercept Only"}\NormalTok{, }\AttributeTok{formula =}\NormalTok{ f0, }\AttributeTok{result =}\NormalTok{ r0)}
\NormalTok{rcat }\OtherTok{\textless{}{-}} \FunctionTok{create\_wflow\_freq\_glm}\NormalTok{(train\_data, fcat)}
\NormalTok{perf\_metrics\_df }\OtherTok{\textless{}{-}} \FunctionTok{append\_performance\_metrics}\NormalTok{(}\AttributeTok{name =} \StringTok{"Categorical Only"}\NormalTok{, }\AttributeTok{formula =}\NormalTok{ fcat, }\AttributeTok{result =}\NormalTok{ rcat)}

\NormalTok{perf\_metrics\_df }\SpecialCharTok{\%\textgreater{}\%} \FunctionTok{arrange}\NormalTok{(AIC) }\SpecialCharTok{\%\textgreater{}\%}\NormalTok{ knitr}\SpecialCharTok{::}\FunctionTok{kable}\NormalTok{(}
    \AttributeTok{booktabs =} \ConstantTok{TRUE}\NormalTok{, }\AttributeTok{caption =} \StringTok{"Performance Metrics for Frequency Models After Adding Categorical Variables"}
\NormalTok{)}
\end{Highlighting}
\end{Shaded}

\begin{longtable}[]{@{}
  >{\raggedright\arraybackslash}p{(\linewidth - 10\tabcolsep) * \real{0.2537}}
  >{\raggedleft\arraybackslash}p{(\linewidth - 10\tabcolsep) * \real{0.1343}}
  >{\raggedleft\arraybackslash}p{(\linewidth - 10\tabcolsep) * \real{0.1343}}
  >{\raggedleft\arraybackslash}p{(\linewidth - 10\tabcolsep) * \real{0.1343}}
  >{\raggedleft\arraybackslash}p{(\linewidth - 10\tabcolsep) * \real{0.1791}}
  >{\raggedleft\arraybackslash}p{(\linewidth - 10\tabcolsep) * \real{0.1642}}@{}}
\caption{Performance Metrics for Frequency Models After Adding
Categorical Variables}\tabularnewline
\toprule\noalign{}
\begin{minipage}[b]{\linewidth}\raggedright
Model
\end{minipage} & \begin{minipage}[b]{\linewidth}\raggedleft
AIC
\end{minipage} & \begin{minipage}[b]{\linewidth}\raggedleft
BIC
\end{minipage} & \begin{minipage}[b]{\linewidth}\raggedleft
Deviance
\end{minipage} & \begin{minipage}[b]{\linewidth}\raggedleft
CV\_Deviance
\end{minipage} & \begin{minipage}[b]{\linewidth}\raggedleft
CV\_LogLoss
\end{minipage} \\
\midrule\noalign{}
\endfirsthead
\toprule\noalign{}
\begin{minipage}[b]{\linewidth}\raggedright
Model
\end{minipage} & \begin{minipage}[b]{\linewidth}\raggedleft
AIC
\end{minipage} & \begin{minipage}[b]{\linewidth}\raggedleft
BIC
\end{minipage} & \begin{minipage}[b]{\linewidth}\raggedleft
Deviance
\end{minipage} & \begin{minipage}[b]{\linewidth}\raggedleft
CV\_Deviance
\end{minipage} & \begin{minipage}[b]{\linewidth}\raggedleft
CV\_LogLoss
\end{minipage} \\
\midrule\noalign{}
\endhead
\bottomrule\noalign{}
\endlastfoot
Categorical Only & 101747.0 & 101795.9 & 71692.51 & 71701.73 &
0.3895784 \\
Intercept Only & 101916.8 & 101926.5 & 71870.30 & 71872.46 &
0.3902322 \\
\end{longtable}

\begin{Shaded}
\begin{Highlighting}[]
\NormalTok{test }\OtherTok{\textless{}{-}} \FunctionTok{anova}\NormalTok{(r0}\SpecialCharTok{$}\NormalTok{fit, rcat}\SpecialCharTok{$}\NormalTok{fit, }\AttributeTok{test =} \StringTok{"Chisq"}\NormalTok{)}
\FunctionTok{cat}\NormalTok{(}\FunctionTok{sprintf}\NormalTok{(}\StringTok{"Is the categorical model a significant improvement over the intercept{-}only model? \%s (p{-}value = \%.4f)}\SpecialCharTok{\textbackslash{}n}\StringTok{"}\NormalTok{, }\FunctionTok{ifelse}\NormalTok{(test}\SpecialCharTok{$}\StringTok{\textasciigrave{}}\AttributeTok{Pr(\textgreater{}Chi)}\StringTok{\textasciigrave{}}\NormalTok{[}\DecValTok{2}\NormalTok{] }\SpecialCharTok{\textless{}} \FloatTok{0.05}\NormalTok{, }\StringTok{"Yes"}\NormalTok{, }\StringTok{"No"}\NormalTok{), test}\SpecialCharTok{$}\StringTok{\textasciigrave{}}\AttributeTok{Pr(\textgreater{}Chi)}\StringTok{\textasciigrave{}}\NormalTok{[}\DecValTok{2}\NormalTok{]))}
\end{Highlighting}
\end{Shaded}

\begin{verbatim}
Is the categorical model a significant improvement over the intercept-only model? Yes (p-value = 0.0000)
\end{verbatim}

\begin{Shaded}
\begin{Highlighting}[]
\NormalTok{delta\_deviance }\OtherTok{\textless{}{-}}\NormalTok{ r0}\SpecialCharTok{$}\NormalTok{performance\_metrics}\SpecialCharTok{$}\NormalTok{CV\_Deviance }\SpecialCharTok{{-}}\NormalTok{ rcat}\SpecialCharTok{$}\NormalTok{performance\_metrics}\SpecialCharTok{$}\NormalTok{CV\_Deviance}
\NormalTok{diff\_degree\_of\_freedom }\OtherTok{\textless{}{-}}\NormalTok{ r0}\SpecialCharTok{$}\NormalTok{fit}\SpecialCharTok{$}\NormalTok{df.residual }\SpecialCharTok{{-}}\NormalTok{ rcat}\SpecialCharTok{$}\NormalTok{fit}\SpecialCharTok{$}\NormalTok{df.residual}
\NormalTok{pvalue }\OtherTok{\textless{}{-}} \FunctionTok{pchisq}\NormalTok{(delta\_deviance, diff\_degree\_of\_freedom, }\AttributeTok{lower.tail =} \ConstantTok{FALSE}\NormalTok{)}
\FunctionTok{cat}\NormalTok{(}\FunctionTok{sprintf}\NormalTok{(}\StringTok{"Anova test based on CV Deviance: delta\_deviance = \%.4f, diff\_degree\_of\_freedom = \%d, p{-}value = \%.4f}\SpecialCharTok{\textbackslash{}n}\StringTok{"}\NormalTok{, delta\_deviance, diff\_degree\_of\_freedom, pvalue))}
\end{Highlighting}
\end{Shaded}

\begin{verbatim}
Anova test based on CV Deviance: delta_deviance = 170.7362, diff_degree_of_freedom = 4, p-value = 0.0000
\end{verbatim}

\begin{Shaded}
\begin{Highlighting}[]
\NormalTok{check\_summary }\OtherTok{\textless{}{-}} \ControlFlowTok{function}\NormalTok{(result) \{}
\NormalTok{    sum\_mod }\OtherTok{\textless{}{-}} \FunctionTok{summary}\NormalTok{(result}\SpecialCharTok{$}\NormalTok{fit)}
\NormalTok{    coefficients }\OtherTok{\textless{}{-}}\NormalTok{ sum\_mod}\SpecialCharTok{$}\NormalTok{coefficients }\SpecialCharTok{\%\textgreater{}\%} \FunctionTok{as.data.frame}\NormalTok{()}
\NormalTok{    coefficients}\SpecialCharTok{$}\NormalTok{p.value }\OtherTok{\textless{}{-}}\NormalTok{ coefficients[, }\StringTok{"Pr(\textgreater{}|z|)"}\NormalTok{]}
\NormalTok{    coefficients }\OtherTok{\textless{}{-}}\NormalTok{ coefficients }\SpecialCharTok{\%\textgreater{}\%} \FunctionTok{select}\NormalTok{(Estimate, }\StringTok{\textasciigrave{}}\AttributeTok{Std. Error}\StringTok{\textasciigrave{}}\NormalTok{, p.value)}
\NormalTok{    coefficients}\SpecialCharTok{$}\NormalTok{significant }\OtherTok{\textless{}{-}} \FunctionTok{ifelse}\NormalTok{(coefficients}\SpecialCharTok{$}\NormalTok{p.value }\SpecialCharTok{\textless{}} \FloatTok{0.05}\NormalTok{, }\StringTok{"Yes"}\NormalTok{, }\StringTok{"No"}\NormalTok{)}
\NormalTok{    coefficients}\SpecialCharTok{$}\NormalTok{multiplier }\OtherTok{\textless{}{-}} \FunctionTok{exp}\NormalTok{(coefficients[, }\StringTok{"Estimate"}\NormalTok{])}
\NormalTok{    coefficients}\SpecialCharTok{$}\NormalTok{multiplier\_lb }\OtherTok{\textless{}{-}} \FunctionTok{exp}\NormalTok{(coefficients[, }\StringTok{"Estimate"}\NormalTok{] }\SpecialCharTok{{-}} \FloatTok{1.96} \SpecialCharTok{*}\NormalTok{ coefficients[, }\StringTok{"Std. Error"}\NormalTok{])}
\NormalTok{    coefficients}\SpecialCharTok{$}\NormalTok{multiplier\_ub }\OtherTok{\textless{}{-}} \FunctionTok{exp}\NormalTok{(coefficients[, }\StringTok{"Estimate"}\NormalTok{] }\SpecialCharTok{+} \FloatTok{1.96} \SpecialCharTok{*}\NormalTok{ coefficients[, }\StringTok{"Std. Error"}\NormalTok{])}
    \FunctionTok{return}\NormalTok{(coefficients)}
\NormalTok{\}}

\NormalTok{plot\_summary }\OtherTok{\textless{}{-}} \ControlFlowTok{function}\NormalTok{(result, }\AttributeTok{add\_text=}\ConstantTok{FALSE}\NormalTok{) \{}
\NormalTok{    coefficients }\OtherTok{\textless{}{-}} \FunctionTok{check\_summary}\NormalTok{(result) }
\NormalTok{    mult\_df }\OtherTok{\textless{}{-}}\NormalTok{ coefficients }\SpecialCharTok{\%\textgreater{}\%} \FunctionTok{filter}\NormalTok{(}\FunctionTok{rownames}\NormalTok{(coefficients) }\SpecialCharTok{!=} \StringTok{"(Intercept)"}\NormalTok{)}
    \FunctionTok{ggplot}\NormalTok{(mult\_df, }\FunctionTok{aes}\NormalTok{(}\AttributeTok{x =} \FunctionTok{rownames}\NormalTok{(mult\_df), }\AttributeTok{y =}\NormalTok{ multiplier, }\AttributeTok{ymin =}\NormalTok{ multiplier\_lb, }\AttributeTok{ymax =}\NormalTok{ multiplier\_ub)) }\SpecialCharTok{+}
        \FunctionTok{geom\_pointrange}\NormalTok{(}
            \AttributeTok{color =} \FunctionTok{ifelse}\NormalTok{(mult\_df}\SpecialCharTok{$}\NormalTok{significant }\SpecialCharTok{==} \StringTok{"Yes"}\NormalTok{, }\StringTok{"black"}\NormalTok{, }\StringTok{"firebrick"}\NormalTok{)}
\NormalTok{        ) }\SpecialCharTok{+}
        \FunctionTok{geom\_hline}\NormalTok{(}\AttributeTok{yintercept =} \DecValTok{1}\NormalTok{, }\AttributeTok{linetype =} \StringTok{"dashed"}\NormalTok{) }\SpecialCharTok{+}
\NormalTok{        \{}\ControlFlowTok{if}\NormalTok{(add\_text) }\FunctionTok{geom\_text}\NormalTok{(}
            \FunctionTok{aes}\NormalTok{(}\AttributeTok{label =} \FunctionTok{round}\NormalTok{(multiplier, }\DecValTok{2}\NormalTok{)), }\AttributeTok{vjust =} \DecValTok{2}\NormalTok{, }
            \AttributeTok{color =} \FunctionTok{ifelse}\NormalTok{(mult\_df}\SpecialCharTok{$}\NormalTok{multiplier }\SpecialCharTok{\textgreater{}} \DecValTok{1}\NormalTok{, }\StringTok{"firebrick"}\NormalTok{, }\StringTok{"darkgreen"}\NormalTok{)}
\NormalTok{        ) }\ControlFlowTok{else} \ConstantTok{NULL}\NormalTok{\} }\SpecialCharTok{+}
        \FunctionTok{coord\_flip}\NormalTok{() }\SpecialCharTok{+}
        \FunctionTok{labs}\NormalTok{(}\AttributeTok{x =} \StringTok{"Coefficient"}\NormalTok{, }\AttributeTok{y =} \StringTok{"Multiplier"}\NormalTok{, }\AttributeTok{title =} \FunctionTok{sprintf}\NormalTok{(}\StringTok{"Coefficient Multipliers and 95\%\% CI}\SpecialCharTok{\textbackslash{}n}\StringTok{Base Premium: \%.2f"}\NormalTok{, }\FunctionTok{exp}\NormalTok{(result}\SpecialCharTok{$}\NormalTok{fit}\SpecialCharTok{$}\NormalTok{coefficients[}\DecValTok{1}\NormalTok{]))) }\SpecialCharTok{+}
        \FunctionTok{theme\_minimal}\NormalTok{()}
\NormalTok{\}}
\NormalTok{coefs }\OtherTok{\textless{}{-}} \FunctionTok{check\_summary}\NormalTok{(rcat)}
\FunctionTok{plot\_summary}\NormalTok{(rcat, }\AttributeTok{add\_text=}\ConstantTok{TRUE}\NormalTok{)}
\end{Highlighting}
\end{Shaded}

\begin{center}\includegraphics{project_MD_files/figure-latex/summary_rcat_fit-1} \end{center}

We see that the anova test indicates that the more complex model with
categorical variables is a significant improvement over the
intercept-only model. We also observe that all the levels we kept are
significant. The multipliers are in line with our expectations based on
the data exploration:

\begin{itemize}
\tightlist
\item
  Female drivers have a higher risk compared to male drivers: the
  multiplier for the female category is of 1.0821361 (95\% CI: 1.0451905
  - 1.1203878)
\item
  Diesel cars are significantly more risky compared to gasoline cars:
  the multiplier for the diesel category is of 1.209678 (95\% CI:
  1.1707423 - 1.2499086)
\item
  Cars that are part of a fleet have a lower risk compared to non-fleet
  cars: the multiplier for the fleet category is of 0.8744607 (95\% CI:
  0.7946295 - 0.9623119)
\item
  Drivers with an omnium coverage are at a lower risk compared to those
  with TPL only: the multiplier for the omnium category is of 0.9062866
  (95\% CI: 0.8779691 - 0.9355175)
\end{itemize}

We can conclude that the categorical variables were well-chosen in the
EDA and have a significant impact in our model. The new model is better
than the null model and needs no adaptation, so we can continue with the
\texttt{rcat} model as our new baseline for further model development.

\begin{Shaded}
\begin{Highlighting}[]
\NormalTok{base\_mod }\OtherTok{\textless{}{-}}\NormalTok{ rcat}
\end{Highlighting}
\end{Shaded}

\paragraph{Binned Continuous
Variables}\label{binned-continuous-variables}

Next, we will add the binned continuous variables to our model to
capture the potential non-linear effects of continuous predictors on the
frequency of claims:

\begin{itemize}
\tightlist
\item
  BM\_binned
\item
  AGEPH\_binned\_merged
\item
  AGEC\_binned
\item
  POWER\_binned\_merged
\end{itemize}

\begin{Shaded}
\begin{Highlighting}[]
\NormalTok{fcontb }\OtherTok{\textless{}{-}} \FunctionTok{paste0}\NormalTok{(fcat, }\StringTok{" + "}\NormalTok{, }\FunctionTok{paste}\NormalTok{(binned\_continuous\_features, }\AttributeTok{collapse =} \StringTok{" + "}\NormalTok{))}
\FunctionTok{cat}\NormalTok{(}\FunctionTok{sprintf}\NormalTok{(}\StringTok{"\textgreater{} Full formula with binned continuous variables: \%s"}\NormalTok{, fcontb))}
\end{Highlighting}
\end{Shaded}

\begin{verbatim}
> Full formula with binned continuous variables: NCLAIMS ~ COVERAGE_BIN + FUEL + FLEET + SEX + BM_binned + AGEPH_binned_merged + AGEC_binned + POWER_binned_merged
\end{verbatim}

\begin{Shaded}
\begin{Highlighting}[]
\NormalTok{rcontb }\OtherTok{\textless{}{-}} \FunctionTok{create\_wflow\_freq\_glm}\NormalTok{(train\_data, fcontb)}
\NormalTok{perf\_metrics\_df }\OtherTok{\textless{}{-}} \FunctionTok{append\_performance\_metrics}\NormalTok{(}\AttributeTok{name =} \StringTok{"All Continuous Binned"}\NormalTok{, }\AttributeTok{formula =}\NormalTok{ fcontb, }\AttributeTok{result =}\NormalTok{ rcontb)}

\NormalTok{test }\OtherTok{\textless{}{-}} \FunctionTok{anova}\NormalTok{(rcat}\SpecialCharTok{$}\NormalTok{fit, rcontb}\SpecialCharTok{$}\NormalTok{fit, }\AttributeTok{test =} \StringTok{"Chisq"}\NormalTok{)}
\FunctionTok{cat}\NormalTok{(}\FunctionTok{sprintf}\NormalTok{(}\StringTok{"Is the full model with binned continuous variables a significant improvement over the categorical model? \%s (p{-}value = \%.4f)}\SpecialCharTok{\textbackslash{}n}\StringTok{"}\NormalTok{, }\FunctionTok{ifelse}\NormalTok{(test}\SpecialCharTok{$}\StringTok{\textasciigrave{}}\AttributeTok{Pr(\textgreater{}Chi)}\StringTok{\textasciigrave{}}\NormalTok{[}\DecValTok{2}\NormalTok{] }\SpecialCharTok{\textless{}} \FloatTok{0.05}\NormalTok{, }\StringTok{"Yes"}\NormalTok{, }\StringTok{"No"}\NormalTok{), test}\SpecialCharTok{$}\StringTok{\textasciigrave{}}\AttributeTok{Pr(\textgreater{}Chi)}\StringTok{\textasciigrave{}}\NormalTok{[}\DecValTok{2}\NormalTok{]))}
\end{Highlighting}
\end{Shaded}

\begin{verbatim}
Is the full model with binned continuous variables a significant improvement over the categorical model? Yes (p-value = 0.0000)
\end{verbatim}

\begin{Shaded}
\begin{Highlighting}[]
\NormalTok{coefs\_contb }\OtherTok{\textless{}{-}} \FunctionTok{check\_summary}\NormalTok{(rcontb)}
\FunctionTok{plot\_summary}\NormalTok{(rcontb)}
\end{Highlighting}
\end{Shaded}

\begin{center}\includegraphics{project_MD_files/figure-latex/binned_continuous_variables-1} \end{center}

There are two main takeaways from the analysis of this more complex
model:

\begin{enumerate}
\def\labelenumi{\arabic{enumi}.}
\tightlist
\item
  With the addition of binned continuous variables, the
  \texttt{SEXfemale} and \texttt{FLEETY} variables are not significant
  anymore. It suggests that they may have been capturing some of the
  information that is now explained by the binned continuous variables.
\item
  The binned continuous variables themselves are all significant at the
  95\% confidence level. The WoE-based binning and merging strategy we
  applied is therefore effective in capturing significant non-linear
  effects of continuous predictors on the frequency of claims.
\end{enumerate}

\begin{Shaded}
\begin{Highlighting}[]
\NormalTok{coefs\_contb }\SpecialCharTok{\%\textgreater{}\%} \FunctionTok{filter}\NormalTok{(significant }\SpecialCharTok{==} \StringTok{"No"}\NormalTok{) }\SpecialCharTok{\%\textgreater{}\%}
    \FunctionTok{arrange}\NormalTok{(}\FunctionTok{desc}\NormalTok{(p.value)) }\SpecialCharTok{\%\textgreater{}\%}
    \FunctionTok{kable}\NormalTok{(}\AttributeTok{booktabs =} \ConstantTok{TRUE}\NormalTok{, }\AttributeTok{caption =} \StringTok{"Non{-}significant Coefficients"}\NormalTok{)}
\end{Highlighting}
\end{Shaded}

\begin{longtable}[]{@{}
  >{\raggedright\arraybackslash}p{(\linewidth - 14\tabcolsep) * \real{0.1075}}
  >{\raggedleft\arraybackslash}p{(\linewidth - 14\tabcolsep) * \real{0.1183}}
  >{\raggedleft\arraybackslash}p{(\linewidth - 14\tabcolsep) * \real{0.1183}}
  >{\raggedleft\arraybackslash}p{(\linewidth - 14\tabcolsep) * \real{0.1075}}
  >{\raggedright\arraybackslash}p{(\linewidth - 14\tabcolsep) * \real{0.1290}}
  >{\raggedleft\arraybackslash}p{(\linewidth - 14\tabcolsep) * \real{0.1183}}
  >{\raggedleft\arraybackslash}p{(\linewidth - 14\tabcolsep) * \real{0.1505}}
  >{\raggedleft\arraybackslash}p{(\linewidth - 14\tabcolsep) * \real{0.1505}}@{}}
\caption{Non-significant Coefficients}\tabularnewline
\toprule\noalign{}
\begin{minipage}[b]{\linewidth}\raggedright
\end{minipage} & \begin{minipage}[b]{\linewidth}\raggedleft
Estimate
\end{minipage} & \begin{minipage}[b]{\linewidth}\raggedleft
Std. Error
\end{minipage} & \begin{minipage}[b]{\linewidth}\raggedleft
p.value
\end{minipage} & \begin{minipage}[b]{\linewidth}\raggedright
significant
\end{minipage} & \begin{minipage}[b]{\linewidth}\raggedleft
multiplier
\end{minipage} & \begin{minipage}[b]{\linewidth}\raggedleft
multiplier\_lb
\end{minipage} & \begin{minipage}[b]{\linewidth}\raggedleft
multiplier\_ub
\end{minipage} \\
\midrule\noalign{}
\endfirsthead
\toprule\noalign{}
\begin{minipage}[b]{\linewidth}\raggedright
\end{minipage} & \begin{minipage}[b]{\linewidth}\raggedleft
Estimate
\end{minipage} & \begin{minipage}[b]{\linewidth}\raggedleft
Std. Error
\end{minipage} & \begin{minipage}[b]{\linewidth}\raggedleft
p.value
\end{minipage} & \begin{minipage}[b]{\linewidth}\raggedright
significant
\end{minipage} & \begin{minipage}[b]{\linewidth}\raggedleft
multiplier
\end{minipage} & \begin{minipage}[b]{\linewidth}\raggedleft
multiplier\_lb
\end{minipage} & \begin{minipage}[b]{\linewidth}\raggedleft
multiplier\_ub
\end{minipage} \\
\midrule\noalign{}
\endhead
\bottomrule\noalign{}
\endlastfoot
FLEETY & -0.0775289 & 0.0489716 & 0.1133891 & No & 0.9254003 & 0.8407061
& 1.018627 \\
SEXfemale & 0.0306299 & 0.0182467 & 0.0932195 & No & 1.0311038 &
0.9948795 & 1.068647 \\
\end{longtable}

We will start by removing the \texttt{FLEET} {[}\^{}fleet\_remove{]}
variable from the model as it is less significant that the \texttt{SEX}
variable.

\begin{Shaded}
\begin{Highlighting}[]
\NormalTok{large\_mod }\OtherTok{\textless{}{-}}\NormalTok{ rcontb}\SpecialCharTok{$}\NormalTok{fit}
\NormalTok{fcontb\_p1 }\OtherTok{\textless{}{-}} \FunctionTok{gsub}\NormalTok{(}\StringTok{" }\SpecialCharTok{\textbackslash{}\textbackslash{}}\StringTok{+ FLEET"}\NormalTok{, }\StringTok{""}\NormalTok{, fcontb)}
\FunctionTok{cat}\NormalTok{(}\FunctionTok{sprintf}\NormalTok{(}\StringTok{"\textgreater{} Formula after removing FLEET: \%s"}\NormalTok{, fcontb\_p1))}
\end{Highlighting}
\end{Shaded}

\begin{verbatim}
> Formula after removing FLEET: NCLAIMS ~ COVERAGE_BIN + FUEL + SEX + BM_binned + AGEPH_binned_merged + AGEC_binned + POWER_binned_merged
\end{verbatim}

\begin{Shaded}
\begin{Highlighting}[]
\NormalTok{rcontb\_p1 }\OtherTok{\textless{}{-}} \FunctionTok{create\_wflow\_freq\_glm}\NormalTok{(train\_data, fcontb\_p1)}
\NormalTok{perf\_metrics\_df }\OtherTok{\textless{}{-}} \FunctionTok{append\_performance\_metrics}\NormalTok{(}\AttributeTok{name =} \StringTok{"Pruned{-}1 Continuous Binned"}\NormalTok{, }\AttributeTok{formula =}\NormalTok{ fcontb\_p1, }\AttributeTok{result =}\NormalTok{ rcontb\_p1)}
\NormalTok{test\_p1 }\OtherTok{\textless{}{-}} \FunctionTok{anova}\NormalTok{(large\_mod, rcontb\_p1}\SpecialCharTok{$}\NormalTok{fit, }\AttributeTok{test =} \StringTok{"Chisq"}\NormalTok{)}
\FunctionTok{cat}\NormalTok{(}\FunctionTok{sprintf}\NormalTok{(}\StringTok{"Are we loosing significant explanatory power by removing FLEET? \%s (p{-}value = \%.4f)}\SpecialCharTok{\textbackslash{}n}\StringTok{"}\NormalTok{, }\FunctionTok{ifelse}\NormalTok{(test\_p1}\SpecialCharTok{$}\StringTok{\textasciigrave{}}\AttributeTok{Pr(\textgreater{}Chi)}\StringTok{\textasciigrave{}}\NormalTok{[}\DecValTok{2}\NormalTok{] }\SpecialCharTok{\textless{}} \FloatTok{0.05}\NormalTok{, }\StringTok{"Yes"}\NormalTok{, }\StringTok{"No"}\NormalTok{), test\_p1}\SpecialCharTok{$}\StringTok{\textasciigrave{}}\AttributeTok{Pr(\textgreater{}Chi)}\StringTok{\textasciigrave{}}\NormalTok{[}\DecValTok{2}\NormalTok{]))}
\end{Highlighting}
\end{Shaded}

\begin{verbatim}
Are we loosing significant explanatory power by removing FLEET? No (p-value = 0.1091)
\end{verbatim}

\begin{Shaded}
\begin{Highlighting}[]
\NormalTok{coefs\_contb\_p1 }\OtherTok{\textless{}{-}} \FunctionTok{check\_summary}\NormalTok{(rcontb\_p1)}
\FunctionTok{plot\_summary}\NormalTok{(rcontb\_p1)}
\end{Highlighting}
\end{Shaded}

\begin{center}\includegraphics{project_MD_files/figure-latex/prune_cont_mod_1-1} \end{center}

Our conclusions hold and we will now proceed to remove the \texttt{SEX}
variable from the model as it is the other non-significant variable.

\begin{Shaded}
\begin{Highlighting}[]
\NormalTok{large\_mod }\OtherTok{\textless{}{-}}\NormalTok{ rcontb\_p1}\SpecialCharTok{$}\NormalTok{fit}
\NormalTok{fcontb\_p2 }\OtherTok{\textless{}{-}} \FunctionTok{gsub}\NormalTok{(}\StringTok{" }\SpecialCharTok{\textbackslash{}\textbackslash{}}\StringTok{+ SEX"}\NormalTok{, }\StringTok{""}\NormalTok{, fcontb\_p1)}
\FunctionTok{cat}\NormalTok{(}\FunctionTok{sprintf}\NormalTok{(}\StringTok{"\textgreater{} Formula after removing SEX: \%s"}\NormalTok{, fcontb\_p2))}
\end{Highlighting}
\end{Shaded}

\begin{verbatim}
> Formula after removing SEX: NCLAIMS ~ COVERAGE_BIN + FUEL + BM_binned + AGEPH_binned_merged + AGEC_binned + POWER_binned_merged
\end{verbatim}

\begin{Shaded}
\begin{Highlighting}[]
\NormalTok{rcontb\_p2 }\OtherTok{\textless{}{-}} \FunctionTok{create\_wflow\_freq\_glm}\NormalTok{(train\_data, fcontb\_p2)}
\NormalTok{perf\_metrics\_df }\OtherTok{\textless{}{-}} \FunctionTok{append\_performance\_metrics}\NormalTok{(}\AttributeTok{name =} \StringTok{"Pruned{-}2 Continuous Binned"}\NormalTok{, }\AttributeTok{formula =}\NormalTok{ fcontb\_p2, }\AttributeTok{result =}\NormalTok{ rcontb\_p2)}
\NormalTok{test\_p2 }\OtherTok{\textless{}{-}} \FunctionTok{anova}\NormalTok{(large\_mod, rcontb\_p2}\SpecialCharTok{$}\NormalTok{fit, }\AttributeTok{test =} \StringTok{"Chisq"}\NormalTok{)}
\FunctionTok{cat}\NormalTok{(}\FunctionTok{sprintf}\NormalTok{(}\StringTok{"Are we loosing significant explanatory power by removing SEX? \%s (p{-}value = \%.4f)}\SpecialCharTok{\textbackslash{}n}\StringTok{"}\NormalTok{, }\FunctionTok{ifelse}\NormalTok{(test\_p2}\SpecialCharTok{$}\StringTok{\textasciigrave{}}\AttributeTok{Pr(\textgreater{}Chi)}\StringTok{\textasciigrave{}}\NormalTok{[}\DecValTok{2}\NormalTok{] }\SpecialCharTok{\textless{}} \FloatTok{0.05}\NormalTok{, }\StringTok{"Yes"}\NormalTok{, }\StringTok{"No"}\NormalTok{), test\_p2}\SpecialCharTok{$}\StringTok{\textasciigrave{}}\AttributeTok{Pr(\textgreater{}Chi)}\StringTok{\textasciigrave{}}\NormalTok{[}\DecValTok{2}\NormalTok{]))}
\end{Highlighting}
\end{Shaded}

\begin{verbatim}
Are we loosing significant explanatory power by removing SEX? No (p-value = 0.1018)
\end{verbatim}

\begin{Shaded}
\begin{Highlighting}[]
\NormalTok{coefs\_contb\_p2 }\OtherTok{\textless{}{-}} \FunctionTok{check\_summary}\NormalTok{(rcontb\_p2)}
\FunctionTok{plot\_summary}\NormalTok{(rcontb\_p2)}
\end{Highlighting}
\end{Shaded}

\begin{center}\includegraphics{project_MD_files/figure-latex/prune_cont_mod_2-1} \end{center}

\begin{Shaded}
\begin{Highlighting}[]
\NormalTok{base\_mod }\OtherTok{\textless{}{-}}\NormalTok{ rcontb\_p2}\SpecialCharTok{$}\NormalTok{fit}
\end{Highlighting}
\end{Shaded}

After pruning the non-significant variables \texttt{FLEET} and
\texttt{SEX}, our final continuous binned model includes only the
significant predictors. This model will serve as the base for the next
enrichment steps in our modeling process.

\paragraph{Enriched Features}\label{enriched-features}

After including our continuous binned variables, we will now add the new
geospatial and enriched features:

\begin{itemize}
\tightlist
\item
  \texttt{REGION\_merged}
\item
  \texttt{MUNICIPALITY\_binned}
\item
  \texttt{POPULATION\_binned}
\end{itemize}

Because the Municipalities are so small, we will first focus on the two
other variables.

\begin{Shaded}
\begin{Highlighting}[]
\NormalTok{f\_enr2 }\OtherTok{\textless{}{-}} \FunctionTok{paste0}\NormalTok{(fcontb\_p2, }\StringTok{" + REGION\_merged + POPULATION\_binned"}\NormalTok{)}
\FunctionTok{cat}\NormalTok{(}\FunctionTok{sprintf}\NormalTok{(}\StringTok{"\textgreater{} Formula after adding two enriched features: \%s"}\NormalTok{, f\_enr2))}
\end{Highlighting}
\end{Shaded}

\begin{verbatim}
> Formula after adding two enriched features: NCLAIMS ~ COVERAGE_BIN + FUEL + BM_binned + AGEPH_binned_merged + AGEC_binned + POWER_binned_merged + REGION_merged + POPULATION_binned
\end{verbatim}

\begin{Shaded}
\begin{Highlighting}[]
\NormalTok{r\_enr2 }\OtherTok{\textless{}{-}} \FunctionTok{create\_wflow\_freq\_glm}\NormalTok{(train\_data, f\_enr2)}
\NormalTok{perf\_metrics\_df }\OtherTok{\textless{}{-}} \FunctionTok{append\_performance\_metrics}\NormalTok{(}\AttributeTok{name =} \StringTok{"Enriched Features"}\NormalTok{, }\AttributeTok{formula =}\NormalTok{ f\_enr2, }\AttributeTok{result =}\NormalTok{ r\_enr2)}
\NormalTok{test\_enr2 }\OtherTok{\textless{}{-}} \FunctionTok{anova}\NormalTok{(base\_mod, r\_enr2}\SpecialCharTok{$}\NormalTok{fit, }\AttributeTok{test =} \StringTok{"Chisq"}\NormalTok{)}
\FunctionTok{cat}\NormalTok{(}\FunctionTok{sprintf}\NormalTok{(}\StringTok{"Are we gaining significant explanatory power by adding enriched features? \%s (p{-}value = \%.4f)}\SpecialCharTok{\textbackslash{}n}\StringTok{"}\NormalTok{, }\FunctionTok{ifelse}\NormalTok{(test\_enr2}\SpecialCharTok{$}\StringTok{\textasciigrave{}}\AttributeTok{Pr(\textgreater{}Chi)}\StringTok{\textasciigrave{}}\NormalTok{[}\DecValTok{2}\NormalTok{] }\SpecialCharTok{\textless{}} \FloatTok{0.05}\NormalTok{, }\StringTok{"Yes"}\NormalTok{, }\StringTok{"No"}\NormalTok{), test\_enr2}\SpecialCharTok{$}\StringTok{\textasciigrave{}}\AttributeTok{Pr(\textgreater{}Chi)}\StringTok{\textasciigrave{}}\NormalTok{[}\DecValTok{2}\NormalTok{]))}
\end{Highlighting}
\end{Shaded}

\begin{verbatim}
Are we gaining significant explanatory power by adding enriched features? Yes (p-value = 0.0000)
\end{verbatim}

\begin{Shaded}
\begin{Highlighting}[]
\NormalTok{coefs\_enr2 }\OtherTok{\textless{}{-}} \FunctionTok{check\_summary}\NormalTok{(r\_enr2)}
\FunctionTok{plot\_summary}\NormalTok{(r\_enr2)}
\end{Highlighting}
\end{Shaded}

\begin{center}\includegraphics{project_MD_files/figure-latex/enriched_features-1} \end{center}

We notice that the \texttt{REGION\_merged} and
\texttt{POPULATION\_binned} variables contribute significantly to the
model, it was therefore worth including them as enriched features in our
final model.

\begin{itemize}
\tightlist
\item
  \texttt{REGION\_mergedBrussels} is actually the second highest
  multiplier of claim frequency in the model: 1.5031001 (95\% CI:
  1.4198746 - 1.5912039)
\item
  \texttt{POPULATION\_binned(1.34e+05,4.53e+05{]}}, so high population
  also stand out as third biggest multiplier of claim frequency in the
  model: 1.4615695 (95\% CI: 1.3620979 - 1.5683053)
\end{itemize}

Note that the largest multiplier of claim frequency is actually an
experience-related variable, it should not be available for a-priori
pricing. Therefore the enriched features we added are the most relevant
for improving the a-priori model. Finally, let's have a quick look at
the \texttt{MUNICIPALITY\_binned} variable to see if it can add anything
to the urban and population density we already considered. Note that we
won't use our workflow because the bins are too small and we will run
into issues with non-observed levels in the cross-validation folds.
We'll make a single GLM fit to test our hypothesis.

\begin{Shaded}
\begin{Highlighting}[]
\NormalTok{f\_enr3 }\OtherTok{\textless{}{-}} \FunctionTok{paste0}\NormalTok{(f\_enr2, }\StringTok{" + MUNICIPALITY\_binned"}\NormalTok{)}
\FunctionTok{cat}\NormalTok{(}\FunctionTok{sprintf}\NormalTok{(}\StringTok{"\textgreater{} Formula after adding MUNICIPALITY\_binned: \%s"}\NormalTok{, f\_enr3))}
\end{Highlighting}
\end{Shaded}

\begin{verbatim}
> Formula after adding MUNICIPALITY_binned: NCLAIMS ~ COVERAGE_BIN + FUEL + BM_binned + AGEPH_binned_merged + AGEC_binned + POWER_binned_merged + REGION_merged + POPULATION_binned + MUNICIPALITY_binned
\end{verbatim}

\begin{Shaded}
\begin{Highlighting}[]
\NormalTok{r\_enr3 }\OtherTok{\textless{}{-}} \FunctionTok{list}\NormalTok{()}
\NormalTok{r\_enr3}\SpecialCharTok{$}\NormalTok{fit }\OtherTok{\textless{}{-}} \FunctionTok{glm}\NormalTok{(}\FunctionTok{as.formula}\NormalTok{(f\_enr3), }\AttributeTok{data =}\NormalTok{ train\_data, }\AttributeTok{family =} \FunctionTok{poisson}\NormalTok{())}
\NormalTok{test\_enr3 }\OtherTok{\textless{}{-}} \FunctionTok{anova}\NormalTok{(r\_enr2}\SpecialCharTok{$}\NormalTok{fit, r\_enr3}\SpecialCharTok{$}\NormalTok{fit, }\AttributeTok{test =} \StringTok{"Chisq"}\NormalTok{)}
\FunctionTok{cat}\NormalTok{(}\FunctionTok{sprintf}\NormalTok{(}\StringTok{"Does MUNICIPALITY\_binned add significant explanatory power? \%s (p{-}value = \%.4f)}\SpecialCharTok{\textbackslash{}n}\StringTok{"}\NormalTok{, }\FunctionTok{ifelse}\NormalTok{(test\_enr3}\SpecialCharTok{$}\StringTok{\textasciigrave{}}\AttributeTok{Pr(\textgreater{}Chi)}\StringTok{\textasciigrave{}}\NormalTok{[}\DecValTok{2}\NormalTok{] }\SpecialCharTok{\textless{}} \FloatTok{0.05}\NormalTok{, }\StringTok{"Yes"}\NormalTok{, }\StringTok{"No"}\NormalTok{), test\_enr3}\SpecialCharTok{$}\StringTok{\textasciigrave{}}\AttributeTok{Pr(\textgreater{}Chi)}\StringTok{\textasciigrave{}}\NormalTok{[}\DecValTok{2}\NormalTok{]))}
\end{Highlighting}
\end{Shaded}

\begin{verbatim}
Does MUNICIPALITY_binned add significant explanatory power? Yes (p-value = 0.0000)
\end{verbatim}

\begin{Shaded}
\begin{Highlighting}[]
\NormalTok{coefs\_enr3 }\OtherTok{\textless{}{-}} \FunctionTok{check\_summary}\NormalTok{(r\_enr3)}
\FunctionTok{plot\_summary}\NormalTok{(r\_enr3)}
\end{Highlighting}
\end{Shaded}

\begin{center}\includegraphics{project_MD_files/figure-latex/municipality_binned-1} \end{center}

\begin{Shaded}
\begin{Highlighting}[]
\NormalTok{base\_mod }\OtherTok{\textless{}{-}}\NormalTok{ r\_enr2}\SpecialCharTok{$}\NormalTok{fit}
\end{Highlighting}
\end{Shaded}

Even if the more complex MUNICIPALITY\_binned variable adds some
explanatory power, we see that individual municipalities do not have a
significant impact on claim frequency (multiplier close to 1). We will
therefore not include \texttt{MUNICIPALITY\_binned} in our final model.

\paragraph{Interaction Effects}\label{interaction-effects}

To close this GLM estimation chapter for the claim frequency model, we
will include the interaction effects we identified with the decision
tree analysis.

\begin{itemize}
\tightlist
\item
  \texttt{BM\_binned} x \texttt{AGEPH\_binned\_merged}
\item
  \texttt{REGION\_merged} x \texttt{POWER\_binned\_merged}
\item
  \texttt{BM\_binned} x \texttt{REGION\_merged}
\item
  \texttt{BM\_01\_FUEL\_not\_gasoline}
\item
  \texttt{BM\_01\_FUEL\_gasoline}
\end{itemize}

The interaction based on expert knowledge will also be considered:

\begin{itemize}
\tightlist
\item
  \texttt{AGEPH\_binned\_merged} x \texttt{POWER\_binned\_merged}
\end{itemize}

Because interactions will create very small cross sections, we'll skip
cross validation and use GLM straightaway.

\begin{Shaded}
\begin{Highlighting}[]
\NormalTok{f\_sp\_int }\OtherTok{\textless{}{-}} \FunctionTok{paste0}\NormalTok{(f\_enr2, }\StringTok{" + "}\NormalTok{, }\FunctionTok{paste}\NormalTok{(interaction\_features, }\AttributeTok{collapse =} \StringTok{" + "}\NormalTok{))}
\FunctionTok{cat}\NormalTok{(}\FunctionTok{sprintf}\NormalTok{(}\StringTok{"\textgreater{} Formula after adding specific interaction: \%s"}\NormalTok{, f\_sp\_int))}
\end{Highlighting}
\end{Shaded}

\begin{verbatim}
> Formula after adding specific interaction: NCLAIMS ~ COVERAGE_BIN + FUEL + BM_binned + AGEPH_binned_merged + AGEC_binned + POWER_binned_merged + REGION_merged + POPULATION_binned + BM_01_FUEL_gasoline + BM_01_FUEL_not_gasoline + BM_binned:AGEPH_binned_merged + REGION_merged:AGEPH_binned_merged + BM_binned:REGION_merged
\end{verbatim}

\begin{Shaded}
\begin{Highlighting}[]
\NormalTok{r\_sp\_int }\OtherTok{\textless{}{-}} \FunctionTok{list}\NormalTok{()}
\NormalTok{r\_sp\_int}\SpecialCharTok{$}\NormalTok{fit }\OtherTok{\textless{}{-}} \FunctionTok{glm}\NormalTok{(}\FunctionTok{as.formula}\NormalTok{(f\_sp\_int), }\AttributeTok{data =}\NormalTok{ train\_data, }\AttributeTok{family =} \FunctionTok{poisson}\NormalTok{())}
\NormalTok{test\_sp\_int }\OtherTok{\textless{}{-}} \FunctionTok{anova}\NormalTok{(base\_mod, r\_sp\_int}\SpecialCharTok{$}\NormalTok{fit, }\AttributeTok{test =} \StringTok{"Chisq"}\NormalTok{)}
\FunctionTok{cat}\NormalTok{(}\FunctionTok{sprintf}\NormalTok{(}\StringTok{"Are we gaining significant explanatory power by adding interactions? \%s (p{-}value = \%.4f)}\SpecialCharTok{\textbackslash{}n}\StringTok{"}\NormalTok{, }\FunctionTok{ifelse}\NormalTok{(test\_sp\_int}\SpecialCharTok{$}\StringTok{\textasciigrave{}}\AttributeTok{Pr(\textgreater{}Chi)}\StringTok{\textasciigrave{}}\NormalTok{[}\DecValTok{2}\NormalTok{] }\SpecialCharTok{\textless{}} \FloatTok{0.05}\NormalTok{, }\StringTok{"Yes"}\NormalTok{, }\StringTok{"No"}\NormalTok{), test\_sp\_int}\SpecialCharTok{$}\StringTok{\textasciigrave{}}\AttributeTok{Pr(\textgreater{}Chi)}\StringTok{\textasciigrave{}}\NormalTok{[}\DecValTok{2}\NormalTok{]))}
\end{Highlighting}
\end{Shaded}

\begin{verbatim}
Are we gaining significant explanatory power by adding interactions? Yes (p-value = 0.0000)
\end{verbatim}

\begin{Shaded}
\begin{Highlighting}[]
\NormalTok{coefs\_sp\_int }\OtherTok{\textless{}{-}} \FunctionTok{check\_summary}\NormalTok{(r\_sp\_int)}
\FunctionTok{plot\_summary}\NormalTok{(r\_sp\_int)}
\end{Highlighting}
\end{Shaded}

\begin{center}\includegraphics{project_MD_files/figure-latex/interactions-1} \end{center}

Note that even if the interactions add some explanatory power, they make
the model more complex, less stable, and harder to interpret. Besides,
only the specific BM between 1 and 11 x Brussels Region has a
significant effect on claim frequency. We will therefore not include the
interaction terms in our final model.

\paragraph{Summary of Performance Metrics and Final GLM
Model}\label{summary-of-performance-metrics-and-final-glm-model}

We can close the GLM frequency model estimation here by looking at the
evolution of our performance metrics across the different modeling
stages:

\begin{Shaded}
\begin{Highlighting}[]
\NormalTok{plot\_leaderboard }\OtherTok{\textless{}{-}} \ControlFlowTok{function}\NormalTok{(data, metric, title, }\AttributeTok{digits =} \DecValTok{2}\NormalTok{) \{}

\NormalTok{    plot\_data }\OtherTok{\textless{}{-}}\NormalTok{ data }\SpecialCharTok{\%\textgreater{}\%}
        \FunctionTok{transmute}\NormalTok{(}
            \AttributeTok{Model =} \FunctionTok{as.character}\NormalTok{(Model),}
            \AttributeTok{value =}\NormalTok{ .data[[metric]]}
\NormalTok{        ) }\SpecialCharTok{\%\textgreater{}\%}
        \FunctionTok{arrange}\NormalTok{(}\FunctionTok{desc}\NormalTok{(value)) }\SpecialCharTok{\%\textgreater{}\%}
        \FunctionTok{mutate}\NormalTok{(}
            \AttributeTok{Model =} \FunctionTok{factor}\NormalTok{(Model, }\AttributeTok{levels =}\NormalTok{ Model),}
            \AttributeTok{best =}\NormalTok{ value }\SpecialCharTok{==} \FunctionTok{min}\NormalTok{(value, }\AttributeTok{na.rm =} \ConstantTok{TRUE}\NormalTok{)}
\NormalTok{        )}

    \FunctionTok{ggplot}\NormalTok{(plot\_data, }\FunctionTok{aes}\NormalTok{(}\AttributeTok{x =}\NormalTok{ value, }\AttributeTok{y =}\NormalTok{ Model)) }\SpecialCharTok{+}
        \FunctionTok{geom\_segment}\NormalTok{(}
            \FunctionTok{aes}\NormalTok{(}\AttributeTok{x =} \FunctionTok{min}\NormalTok{(value), }\AttributeTok{xend =}\NormalTok{ value, }\AttributeTok{yend =}\NormalTok{ Model),}
            \AttributeTok{colour =} \StringTok{"grey75"}
\NormalTok{        ) }\SpecialCharTok{+}
        \FunctionTok{geom\_point}\NormalTok{(}\FunctionTok{aes}\NormalTok{(}\AttributeTok{colour =}\NormalTok{ best), }\AttributeTok{size =} \DecValTok{3}\NormalTok{) }\SpecialCharTok{+}
        \FunctionTok{geom\_text}\NormalTok{(}
            \FunctionTok{aes}\NormalTok{(}\AttributeTok{label =} \FunctionTok{formatC}\NormalTok{(value, }\AttributeTok{format =} \StringTok{"f"}\NormalTok{, }\AttributeTok{digits =}\NormalTok{ digits)),}
            \AttributeTok{hjust =} \SpecialCharTok{{-}}\FloatTok{0.15}\NormalTok{,}
            \AttributeTok{size =} \FloatTok{3.3}
\NormalTok{        ) }\SpecialCharTok{+}
        \FunctionTok{scale\_colour\_manual}\NormalTok{(}
            \AttributeTok{values =} \FunctionTok{c}\NormalTok{(}\StringTok{"FALSE"} \OtherTok{=} \StringTok{"grey"}\NormalTok{, }\StringTok{"TRUE"} \OtherTok{=} \StringTok{"darkgreen"}\NormalTok{),}
            \AttributeTok{guide =} \StringTok{"none"}
\NormalTok{        ) }\SpecialCharTok{+}
        \FunctionTok{scale\_x\_continuous}\NormalTok{(}
            \AttributeTok{expand =} \FunctionTok{expansion}\NormalTok{(}\AttributeTok{mult =} \FunctionTok{c}\NormalTok{(}\FloatTok{0.08}\NormalTok{, }\FloatTok{0.35}\NormalTok{))}
\NormalTok{        ) }\SpecialCharTok{+}
        \FunctionTok{labs}\NormalTok{(}\AttributeTok{title =}\NormalTok{ title, }\AttributeTok{x =} \ConstantTok{NULL}\NormalTok{, }\AttributeTok{y =} \ConstantTok{NULL}\NormalTok{) }\SpecialCharTok{+}
        \CommentTok{\# remove x{-}axis labels}
        \FunctionTok{theme\_minimal}\NormalTok{(}\AttributeTok{base\_size =} \DecValTok{11}\NormalTok{) }\SpecialCharTok{+}
        \FunctionTok{theme}\NormalTok{(}
            \AttributeTok{panel.grid.major.y =} \FunctionTok{element\_blank}\NormalTok{(),}
            \AttributeTok{panel.grid.minor =} \FunctionTok{element\_blank}\NormalTok{(),}
            \AttributeTok{plot.title =} \FunctionTok{element\_text}\NormalTok{(}\AttributeTok{face =} \StringTok{"bold"}\NormalTok{),}
            \AttributeTok{axis.text.x =} \FunctionTok{element\_blank}\NormalTok{()}
\NormalTok{        )}
\NormalTok{\}}

\NormalTok{p1 }\OtherTok{\textless{}{-}} \FunctionTok{plot\_leaderboard}\NormalTok{(perf\_metrics\_df, }\StringTok{"AIC"}\NormalTok{, }\StringTok{"AIC"}\NormalTok{)}
\NormalTok{p2 }\OtherTok{\textless{}{-}} \FunctionTok{plot\_leaderboard}\NormalTok{(perf\_metrics\_df, }\StringTok{"BIC"}\NormalTok{, }\StringTok{"BIC"}\NormalTok{)}
\NormalTok{p3 }\OtherTok{\textless{}{-}} \FunctionTok{plot\_leaderboard}\NormalTok{(perf\_metrics\_df, }\StringTok{"CV\_Deviance"}\NormalTok{, }\StringTok{"CV deviance"}\NormalTok{)}
\NormalTok{p4 }\OtherTok{\textless{}{-}} \FunctionTok{plot\_leaderboard}\NormalTok{(}
\NormalTok{    perf\_metrics\_df, }\StringTok{"CV\_LogLoss"}\NormalTok{, }\StringTok{"CV Poisson log loss"}\NormalTok{, }\AttributeTok{digits =} \DecValTok{4}
\NormalTok{)}

\NormalTok{(p1 }\SpecialCharTok{|}\NormalTok{ p2) }\SpecialCharTok{/}\NormalTok{ (p3 }\SpecialCharTok{|}\NormalTok{ p4) }\SpecialCharTok{+}
    \FunctionTok{plot\_annotation}\NormalTok{(}
        \AttributeTok{title =} \StringTok{"Model performance leaderboard"}\NormalTok{,}
        \AttributeTok{subtitle =} \StringTok{"Best model highlighted in green"}
\NormalTok{)}
\end{Highlighting}
\end{Shaded}

\begin{center}\includegraphics{project_MD_files/figure-latex/performance_metrics-1} \end{center}

Based on this leaderboard, we can confirm that our final model with
categorical, significant binned continuous features and significant
enriched features performs the best across all evaluated metrics. The
other models have the same ranking across AIC, cross-validated deviance,
and cross-validated Poisson log loss. The only difference we noticed is
in the BIC metric that penalizes more model complexity. For the three
models with continuous features, this penalization difference is
sufficient to alter their ranking and favor the simpler models.

\subsubsection{GAM Models}\label{gam-models}

\subsubsection{Final Frequency Model}\label{final-frequency-model}

\subsection{Severity Modelling}\label{severity-modelling}

\subsubsection{GLM Models}\label{glm-models-1}

We went through a lot of details in the data preparation and model
selection process for the frequency component of our analysis. For the
severity component, we will go through a similar process in much lesser
depth to keep this report concise. We will nevertheless go through the
following steps:

\begin{itemize}
\tightlist
\item
  Filter out non-relevant observations for the severity analysis
\item
  Look at the distribution of the severity across individual features
\item
  Fit step-by-step GLM models to find our best model. We'll skip
  cross-validation for brevity.
\end{itemize}

Watch out that we are not using the canonical link function for the
Gamma distribution, which is the inverse link. This means that the sum
of predicted values might not equal the total observed cost. In a
real-world setting, we would need to account for this discrepancy when
setting premiums. The reason why we will use the log link instead of the
canonical inverse link is to ensure that our predicted values are always
positive and to simplify the interpretation of the model coefficients
(as multiplicative effects on the expected value of the response
variable).

\paragraph{Data Preparation For Severity
Analysis}\label{data-preparation-for-severity-analysis}

\begin{Shaded}
\begin{Highlighting}[]
\NormalTok{train\_severity }\OtherTok{\textless{}{-}}\NormalTok{ train\_data[train\_data}\SpecialCharTok{$}\NormalTok{NCLAIMS }\SpecialCharTok{\textgreater{}} \DecValTok{0}\NormalTok{, ]}
\NormalTok{avg\_cost\_int }\OtherTok{\textless{}{-}} \FunctionTok{glm}\NormalTok{(AVG }\SpecialCharTok{\textasciitilde{}} \DecValTok{1}\NormalTok{, }\AttributeTok{data =}\NormalTok{ train\_severity, }\AttributeTok{family =} \FunctionTok{Gamma}\NormalTok{(}\AttributeTok{link =} \StringTok{"log"}\NormalTok{))}
\NormalTok{avg\_cost\_int\_sum }\OtherTok{\textless{}{-}} \FunctionTok{summary}\NormalTok{(avg\_cost\_int)}
\NormalTok{beta }\OtherTok{\textless{}{-}}\NormalTok{ avg\_cost\_int\_sum}\SpecialCharTok{$}\NormalTok{coefficients[}\DecValTok{1}\NormalTok{, }\DecValTok{1}\NormalTok{]}
\NormalTok{beta\_sd }\OtherTok{\textless{}{-}}\NormalTok{ avg\_cost\_int\_sum}\SpecialCharTok{$}\NormalTok{coefficients[}\DecValTok{1}\NormalTok{, }\DecValTok{2}\NormalTok{]}
\FunctionTok{cat}\NormalTok{(}\FunctionTok{sprintf}\NormalTok{(}\StringTok{"Estimated mean of average cost: \%.2f (95\%\% CI: \%.2f to \%.2f)"}\NormalTok{, }\FunctionTok{exp}\NormalTok{(beta), }\FunctionTok{exp}\NormalTok{(beta }\SpecialCharTok{{-}} \FloatTok{1.96} \SpecialCharTok{*}\NormalTok{ beta\_sd), }\FunctionTok{exp}\NormalTok{(beta }\SpecialCharTok{+} \FloatTok{1.96} \SpecialCharTok{*}\NormalTok{ beta\_sd)))}
\end{Highlighting}
\end{Shaded}

\begin{verbatim}
Estimated mean of average cost: 1698.68 (95% CI: 1418.70 to 2033.92)
\end{verbatim}

\begin{Shaded}
\begin{Highlighting}[]
\NormalTok{mean\_avg }\OtherTok{\textless{}{-}} \FunctionTok{mean}\NormalTok{(train\_severity}\SpecialCharTok{$}\NormalTok{AVG)}
\NormalTok{sd\_avg }\OtherTok{\textless{}{-}} \FunctionTok{sd}\NormalTok{(train\_severity}\SpecialCharTok{$}\NormalTok{AVG)}
\FunctionTok{cat}\NormalTok{(}\FunctionTok{sprintf}\NormalTok{(}\StringTok{"Observed mean of average cost: \%.2f (sd{-}based 95\%\% CI: \%.2f to \%.2f)"}\NormalTok{, mean\_avg, mean\_avg }\SpecialCharTok{{-}} \FloatTok{1.96} \SpecialCharTok{*}\NormalTok{ sd\_avg, mean\_avg }\SpecialCharTok{+} \FloatTok{1.96} \SpecialCharTok{*}\NormalTok{ sd\_avg))}
\end{Highlighting}
\end{Shaded}

\begin{verbatim}
Observed mean of average cost: 1698.68 (sd-based 95% CI: -35232.05 to 38629.42)
\end{verbatim}

Why are t

\subsubsection{Final Severity Model}\label{final-severity-model}

\section{Global Overview}\label{global-overview}

\subsection{Importance of Features}\label{importance-of-features}

\subsection{Final Premium Structure}\label{final-premium-structure}

\section{Conclusion}\label{conclusion}

\end{document}
